% Gerado por tools/tei2latex.py a partir de entrada.xml. Regenerar, não editar.

%%>>> parte-1

%%>>> cap-3

%%>>> art-15

%%>>> t78

c) Emquanto á \emph{vantagem}, é necessario distinguir. Ou se trata da vantagem relativa á vida \emph{intellectual} do homem, ou da vantagem relativa á sua vida \emph{organica}. --- Se se trata da vantagem relativa á vida \emph{intellectual}, \emph{á vista} e o \emph{ouvido} são superiores aos outros sentidos: pois a \emph{vista} ajuda a conhecer as differenças das coisas, e o \emph{ouvido} serve para perceber o \pg{113} \pg{114}  ensino dos sabios. --- Se se trata da vantagem relativa á vida \emph{organica}, o \emph{tacto} e o \emph{gosto} são superiores aos outros sentidos \footnote{Dissemos que, se se trata da vantagem relativa á vida \emph{intellectual}, a \emph{vista} e o \emph{ouvido} são superiores aos outros sentidos. E, sempre sob este aspecto, o \emph{ouvido} é superior á \emph{vista}. Por quanto o homem é essencialmente \emph{discipulo}, o numero dos conhecimentos que recebemos dos outros é superior ao numero do conhecimentos que adquirimos pelos nossos esforços, e por isso o sentido do \emph{ouvido} é mais necessario para o desenvolvimento da intelligencia do que o sentido da \emph{vista}. (S. Thomaz, \emph{de sens}., lect. ID).\par Emquanto ao \emph{desenvolvimento}, primeiramente é o \emph{tacto}, que é o fundamento de todos os sentidos, depois o \emph{gosto}, em seguida o \emph{olfacto} e a \emph{vista}, e finalmente o \emph{ouvido}. Com relação ao \emph{tacto} diz S. Thomaz: « Cum sensus tactus sit quasi fundamentum aliorum sensuum, in organo uniuscujusque sensus non solum invenitur proprietas illius sensus, cujus est organum proprium, sed etiam proprietas tactus » (\emph{De} Ver., q. 22, a. 5).}. \textbar{}

%%>>> art-16

\section*{Artigo II.\\ Sentidos internos.}
\addcontentsline{toc}{section}{Artigo II. --- Sentidos internos.}

%%>>> t79

\tese{79}{Sentidos internos.}
\emph{Sentidos internos} são os que existem no \emph{interior} do organismo, e por cujo intermedio o homem --- \emph{directa} e \emph{immediatamente} attinge, por varios modos, as proprias \emph{sensações} ou reproduz as \emph{imagens sensiveis} dos objectos \emph{externos}, --- e \emph{indirecta} e \emph{mediatamente} percebe os mesmos objectos externos. Por isso o exercicio dos sentidos \emph{internos} suppõe o exercicio dos sentidos \emph{externos} \footnote{O homem e os animaes mais perfeitos são dotados de sentidos \emph{externos} e \emph{internos}. Se é certo que um ou outro sentido \emph{externo} é mais perfeito em alguns animaes do que em nós (assim a \emph{vista} da aguia é mais perfeita que a nossa), é tambem certo que os sentidos \emph{internos} são mais perfeitos no homem do que no animal, porque em nós estão unidos com a intelligencia.}.

%%>>> t80

\tese{80}{Numero dos sentidos internos.}
Os sentidos \emph{internos} são \emph{quatro}, a saber: o \emph{sentido commum}, a \emph{imaginação}, a \emph{estimativa}, e a \emph{memoria sensitiva} \footnote{Advirta-se a necessidade e a ordem dos \emph{quatro} sentidos \emph{internos}.\par Cada sentido \emph{externo} percebe um só objecto, que lhe é \emph{proprio}; ma não percebe o acto que elle produz, nem os objectos que são \emph{proprios} dos outros sentidos. Ora, para a perfeição da vida \emph{sensitiva}, é necessaria a existencia de uma faculdade, que perceba as sensações de todos os sentidos \emph{externos} e as differenças que existem entre ellas. Esta faculdade é o \emph{sentido commum}, que se não parcebe o proprio acto, percebe comtudo os actos de todos os sentidos \emph{externos}. --- Todavia o \emph{sentido commum}, percebendo uma sensação \emph{actual}, limita-se ao conhecimento dos objectos \emph{presentes}, mas não dos \emph{ausentes}. f, pois, necessaria uma outra faculdade, que possa conservar e reproduzir as\pgn{115}\emph{imagens} dos objectos, mesmo quando estes não estão presentes. Esta faculdade é a \emph{imaginação}. --- Mas a \emph{imaginação}, se conserva e reproduz as \emph{imagens} dos objectos percebidos, não conhece se aquelles objectos são \emph{uteis} ou nocivos para a vida. E, pois, necessaria uma faculdade que nos objectos perceba certas propriedades de \emph{vantagem} ou de \emph{damno}, que escapam ao alcance dos outros sentidos. Esta faculdade é a \emph{força estimativa}. --- A \emph{estimativa}, porém, embora perceba aquellas propriedades, não póde conservar as percepções e as coisas \emph{passadas}. E, pois, necessaria uma outra faculdade, que conserve as percepções e as coisas \emph{passadas}. Esta faculdade é a \emph{memoria sensitiva}. (Cf. \emph{Sum}. \emph{Theol}., p. I, q. 78, a. 4; QQ. DD. \emph{De anima}, a. 18).}.

%%>>> t81

\pg{115} \tese{81}{Sentido commum.}
\emph{Sentido commum} é a faculdade pela qual o homem percebe \emph{as sensações dos sentidos externos} e as \emph{suas differengas}. --- Esta faculdade, pois, é o centro que \emph{recolhe} as sensações dos sentidos \emph{externos}, \emph{comparando-as} e \emph{distinguindo-as} umas das outras.

O \emph{sentido commum} existe no homem; pois a consciencia attesta que percebemos o acto dos sentidos \emph{externos}, \emph{comparamos} as sensações externas e \emph{distinguimos} uma sensação de outra, o \emph{branco} do \emph{amargo}, o \emph{som} do \emph{sabor}, ete. \footnote{Sao duas, pois, as funcções do \emph{sentido commum}; a primeira é a de \emph{perceber} as sensações ou os actos dos singulos sentidos externos (por isso chama-se tambem \emph{consciencia sensitiva}); a segunda é a de \emph{comparar} essas sensações, o de as \emph{distinguir} umas das outras. --- Esta segunda funcção é mais importante do que a primeira. Na verdade, de que nos serviria a ceniralisac\&o de todos os dados dos sentidos externos, se os não podessemos comparar entre si e distinguir um do outro, ou se não podessemos distinguir os objectos presentes, percebidos pelos sentidos externos, das representações \emph{dmaginarias} de objectos ausentes? Ora a o sentido commum que compara e distingue as qualidades dos corpos externos, percebidas pelos diversos sentidos externos, como tambem percebe e distingue os actos dos sentidos, quer internos quer externos, attribuindo a cada sentido o acto que lhe pertence (S. Th., \emph{in} 3 \emph{de anima}, lect. 3). --- Pelo facto de perceber, comparar e distinguir as sensações dos sentidos externos, o sentido commum \emph{completa} as percepções de um sentido com as de um outro, de modo que chega a representar um objecto na sua \emph{totalidade}. E esta obra de \emph{complemento} tal-a o sentido commum, attribuindo ao mesmo objecto as diversas qualidades percebidas pelos differentes sentidos (assim a uma flor attribue-se a tal \emph{cór}, o tal perfume, eco.), ou explicando os dados incertos e obscuros de um sentido pelos dados certos e claros de um outro (assim o tacto auxilia uma vista fraca).\par Objectam: Distinguindo-se as faculdades pelos objectos \emph{formaes}, como - éque os objectos dos cinco sentidos externos, \emph{formalmente} diversos entre si, podem constituir um só objecto \emph{formal} de uma só faculdade, qual é o sen- - \emph{tido commum} ?\par Resposta. A coisa é muito simples. Varios objectos \emph{formalmente} diversos com relação a faculdades inferiores podem constituir um só objecto \emph{formal} com relação a uma faculdade superior, porque podem ser considerados por esta sob um. só aspecto, mais extenso e mais elevado. E' o que acontece no\pgn{116}nosso caso. A cór, o: som, o sabor, etc., são objectos \emph{formalmente} diversos e por isso especificam os diversos sentidos externos; mas se a cór, o som, o sabor, etc., se consideram não pelo que são em si memos, mas \emph{em quanto são percebidos}, então podem referir-se e referem-se a uma só faculdade, que é o \emph{sentido commum}. De facto, o objecto proprio e directo d'esta faculdade não é a cór, mas é a \emph{sensação} da cór; não é o sabor, mas é a \emph{sensação} do sabor, etc. (Cf. \emph{Sum}. \emph{Theol}, p.l, q. (4, a. 3 ad 4).\par Continuam: O \emph{sentido commum}, como faculdade especial, é inteiramente inutil; pois, para a simultanea percepção dos objectos e das modificações dos sentidos externos, basta a alma, que, sendo o principio de que emanam todas as faculdades, percebe todas as percepções d'ellas.\par Respondemos que a existencia do \emph{sentido commum}, como faculdade especial, é absolutamente necessaria. Na verdade, a consciencia attesta que percebemos ao mesmo tempo as diversas sensações dos sentidos \emph{externos} e os objectos que as determinam. Ora o sujeito, que percebe as sensações e os seus objectos, não é a essencia da alma, mas uma faculdade que deriva da propria essencia da alma; pois a creatura não opéra \emph{immediatamente} pela sua essencia, mas só por meio das faculdades, que derivam da essencia. Essa faculdade, que percebe as sensações e os seus objectos, não póde ser \emph{intellectual}, como provaremos d'aqui a pouco. Logo deve ser \emph{sensitiva}. Sendo \emph{sensitiva}, ou identifica-se com algum dos sentidos externos, ou é distincta de cada um d'elles. Ora não póde identificar-se com algum dos sentidos externos; pois nenhum sentido externo, por não reflectir sobre si mesmo, percebe o seu proprio acto e o acto dos outros sentidos, assim \emph{a-vista} não percebe a propria \emph{visão}, etc. Logo deve ser uma faculdade distincta de todos os sentidos externos. Essa faculdade \emph{sensitiva}, distincta de cada um dos sentidos externos, chama-se \emph{sentido commum} ou \emph{sentido interno}. Logo a existencia do \emph{sentido commum}, como faculdade especial, é absolutamente necessaria.\par Insistem: Pelo menos no homem é desnecessario o \emph{sentido commuwum} ; pois o homem póde, por meio da intelligencia de que é dotado, reflectir sobre os actos dos sentidos externos.\par Resposta. O \emph{sentido commum}, por ser uma faculdade da \emph{sensibilidade}, deve encontrar-se no homem, aliás este não seria, como effectivamente é, um animal perfeito. --- De mais, a reflexão que a \emph{ntelligencia} exerce sobre os actos dos sentidos externos é \emph{especificamente} diversa da que o \emph{sentido commum} exerce sobre os mesmos actos. Na verdade, a \emph{intelligencia} percebe os actos dos sentidos externos, comprehendendo a sua \emph{intrinseca natureza} e a \emph{relação} que têm com o principio de que emanam; --- ao passo que o \emph{seztido commum} percebe aquelles actos, considerando-os como :\emph{neros} factos, que dependem dos orgãos e são acompanhados por uma mudança .corporea; e por isso o \emph{sentido commum} é uma faculdade \emph{sensitiva} e não \emph{intellectual}.\par Pergunta-se: Existirá no homem, como tambem nos animaes mais perfeitos, uma faculdade organica, que perceba o acto do \emph{sentido commum} ?\par Respondemos que essa faculdade organica não existe, nem é necessario\pgn{117}que exista, --- Não existe; porque, devendo as modificações e os actos de uma faculdade sensitiva ser percebidos por uma faculdade sensitiva superior, é claro que os actos do \emph{sentido commwum}, que não está subordinado \& um outro \emph{sentido commum} superior, não são percebidos por outra faculdade sensitiva, alias deveriamos admittir o \emph{progresso} até ao \emph{infinito}, \emph{o} que é absurdo. --- Nem é necessaria; porque ao animal basta que tenha a consciencia das suas sensações. Do que lhe serviria ter consciencia da propria consciencia?}.

%%>>> t82

\tese{82}{Objecto do sentido commum.}
O objecto do \emph{sentido commum} é duplice: um \emph{directo} e \emph{proprio}, outro \emph{indirecto} e \emph{concomitante}. --- O \emph{directo} e \emph{proprio} é constituido pelas \emph{sensações} dos sentidos \emph{externos}. --- O \emph{indirecto} e \emph{concomitante} é formado \pg{116} \pg{117}  pelo proprio \emph{organismo} do homem e pelo \emph{objecto} das \textbf{sensações} externas; pois não póde haver sensação sem um \emph{sujeito} que sente, e sem um \emph{objecto} que é sentido \footnote{Pelo \emph{sentido commum} percebemos \emph{directa} e \emph{propriamente} as \emph{sensações} dos sentidos \emph{externos}, mas \emph{indirectamente} e \emph{por concomitancia} percebemos o nosso \emph{organismo} e o \emph{objecto} das mesmas sensações. Na verdade, a sensação é um phenomeno \emph{objectivo} e \emph{subjectivo}. --- Emquanto \emph{objectivo}, representa um \emph{objecto}: pois não ha percepção sem um \emph{objecto} percebido. --- Emquanto \emph{subjectivo}, \emph{é a} modificação do sentido, ou, melhor, é o proprio sentido modificado. Ora o sentido é o proprio orgão, é o proprio corpo animado. Logo, pelo \emph{sentido commum}, percebemos, embora \emph{indirectamente} e \emph{por concomitancia}, o nosso proprio organismo e os seus diversos estados.\par A razão, por que, pelo \emph{sentido comum}, percebemos o nosso organismo e os seus diversos estados, é a seguinte. O \emph{sentido commum} percebe as sensações dos cinco sentidos \emph{externos}, e por isso percebe as do \emph{tacto}. Ora os nervos do sentido do \emph{facto} são os mesmos que os da \emph{sensibilidade tactil}, e estes --- encontram-se, mais ou menos, \emph{em toda a superficie da pelle e da} mucosa. Logo a mesma faculdade, que percebe as sensações do \emph{facto}, percebe as da \emph{sensibilidade tactil}, e por isso o organismo e os seus diversos estados.\par D'onde se vé a inutilidade d'aquella opinião, que admitte o \emph{sentido commum} para a percepção das \emph{sensações} dos sentidos externos, e um \emph{sentido intimo}, distincto do \emph{sentido commum}, para a percepção do nosso organismo e dos seus diversos estados; pois esses dois objectos são reductiveis á mesma faculdade, e não se devem multiplicar as causas sem necessidade.\par Erra, pois, \emph{Reid}, quando diz que, pelo \emph{sentido commum}, unica e exclusivamente percebemos a sensação ou modificação do sentido, e que o sujeito sensiente só se conhece por meio de um raciocinio, baseado num instincto cego e inconsciente. Este raciocinio seria o seguinte: \emph{toda a modificação suppõe um sujeito modificado}. --- Tambem os animaes percebem o proprio organismo e os seus varios estados, e todavia não raciocinam.\par Alguns escriptores dizem que, pelo \emph{sentido commum} ou \emph{intimo}, não só percebemos as sensações, mas tambem as percepções \emph{intellectuaes}. --- Esta asserção é falsa. Uma e a mesma faculdade não póde attingir objectos \emph{especifica}\pgn{118} e \emph{formalmente} diversos, como são as \emph{sensações} e as percepções \emph{intellectuaes}. O \emph{sentido commum}, como provaremos, é faculdade \emph{sensitiva}, e por isso o seu objecto é \emph{sensivel} e não \emph{intellectual}.\par Pergunta-se: As sensações dos sentidos externos são todas e sempre percebidas pelo sentido commum, ou haverá sensações externas não percebidas por esse sentido ?\par Resposta. As opiniões dos modernos são varias e contrarias. Uns, como \emph{Cousin}, \emph{Stuart Mill}, \emph{Rabier}, e, em geral, os \emph{subjectivistas}, negam que possa haver uma sensação externa não advertida pelo. sentido commum, porque, dizem, seria essa uma sensação que nada sente. Outros, como \emph{Leibnitz}, \emph{Reid} e muitos outros com quasi todos os physiologistas, sustentam que não só pode haver, mas que effectivamente ha sensações, não percebidas pelo sentido com- "mum, isto é, das quaes não temos consciencia. --- Esta segunda opinião é á verdadeira. Com effeito, se repugna, uma percepção sem objecto, não repugna um objecto sem percepção; e por isso, se não, pode dar-se uma sensação, que nada sinta, pode dar-se uma sensação que não seja sentida. --- Além d'isso, o acto dos sentidos externos pode dar-se sem o acto do sentido commum. --- Não só não repugnam essas sensações não sentidas, mas dão-se effectivamente. \emph{A} experiencia attesta: --- \emph{a}) que, quanto mais intensa fôr a applicação dos sentidos no objecto, tanto menos viva é consciencia d'essa percepção, chegando, ás vezes, a não haver consciencia nenhuma d'essa percepção, como no \emph{extasis}; --- \emph{b}) que, na memoria sensitiva, ha muitas imagens, paixões, percepções, cuja presença não é percebida actualmente pelo sentido commum; --- €) que ha muitas percepções, que, passadas desapercebidas no principio, se tornam depois presentes á consciencia, - como ha outras, que, advertidas no principio, começam a escapar gradualmente ao alcance da consciencia, até que, em occasião opportuna, reapparecem; --- \emph{d}) que, no \emph{somnambulismo}, o homem escuta e responde, mas não tem consciencia dos seus actos, nem das suas palavras (Cf. Farges, \emph{le Cerveau}, pag. 295 e seg.)}.

%%>>> t83

\tese{83}{O sentido commum é uma faculdade sensitiva.}

a) Uma faculdade é \emph{sensitiva}, quando o seu objecto é \emph{sensivel}; pois as faculdades distinguem-se e especificam-se, em ultima analyse, pelos seus objectos. Ora o objecto do \emph{sentido commum é sensivel}; pois que é constituido pelas \emph{sensações} dos sentidos \emph{externos}, e pelas suas differencas. Logo o \emph{sentido commum} é uma faculdade \emph{sensitiva}.

b) Uma faculdade que existe nos animaes não póde deixar de ser \emph{sensitiva}, sendo os animaes destituidos de faculdades intellectuaes. Ora o \emph{sentido commum} existe nos animaes; pois elles, nas suas operações, mostram perceber, ao mesmo tempo, varias sensações, relativas a diversos sentidos externos, e preferir umas ás outras, como se as tivessem comparado entre si. \textbf{Logo} o \emph{sentido commum} é uma faculdade \emph{sensitiva} \footnote{\emph{Reid} e \emph{Mallebranche} opinam que o \emph{sentido commum} não é uma faculdade \emph{sensitiva}, mas é \emph{intellectual}.\par Dizem: Se o \emph{sentido commun} fosse uma faculdade \emph{sensitiva}, seria incapaz de reflexão e de juizo; pois estas operações são proprias da intelligenocia. Ora o \emph{sentido conmum} é dotado de reflexão e julga das diversas sensações. Logo não é uma faculdade \emph{sensitiva}.\par Resposta. O \emph{sentido commuwm} seria faculdade \emph{intellectual}, se percebesse as \emph{proprias} operações por um acto \emph{reflexo}, \emph{e} se julgasse das diversas sensações por uma comparação tambem \emph{reflexa}; porque a \emph{reflexdo} é um acto exlusivo da intelligencia. Mas não é assim. O \emph{sentido commun} só percebe as sensações dos sentidos externos por um acto \emph{directo}, --- como tambem conhece a differenca, que existe entre as sensações, não emquanto exerce uma comparação\pgn{119}\emph{reflexa}, mas de um modo \emph{directo}, isto é, emquanto recebe uma modificação diversa segundo a diversidade das sensações. Logo o \emph{sentido commwum} não é uma faculdade \emph{intellectual}, mas \emph{sensitiva}.\par Todavia concedemos que o \emph{sentido commun}, embora não exceda a esphera da sensibilidade, é menos material que os sentidos externos; porque o seu objecto \emph{directo} e \emph{immediato} é a sensação, que não é materia.}.

%%>>> t84

\pg{118} \pg{119} \tese{84}{O sentido commum é faculdade distincta dos sentidos externos.}

a) \emph{O sentido commum} percebe as sensações dos sentidos \emph{externos}. Mas nenhum sentido \emph{externo} póde perceber a propria sensação. Na verdade, para isso seria necessario que o sentido reflecetisse sobre si mesmo. Ora o sentido, por ser faculdade \emph{material,não} póde reflectir sobre si mesmo. Logo o \emph{sentido commum} é uma faculdade distincta dos sentidos \emph{externos}.

b) O \emph{sentido commum} não só percebe os \emph{actos} dos sentidos \emph{externos}, mas tambem a diversidade, que existe entre os objectos \emph{proprios} dos varios sentidos. Ora nenhum dos sentidos \emph{externos} póde perceber a diversidade que existe entre os objectos proprios dos outros sentidos, pois cada um percebe \emph{exclusivamente} o objecto que Ihe é \emph{proprio}; assim a \emph{vista} não percebe a diversidade entre-a cór e o \emph{sabor}, mas só percebe a côr, que é o seu objecto \emph{proprio}. Logo o \emph{sentido commum} é uma faculdade distinecta dos sentidos \emph{externos} \footnote{Segundo \emph{Suares}, cada um dos sentidos \emph{e«ternos} percebe, com o mesmo acto, a si mesmo e o seu objecto. --- \emph{Condillac} disse que o \emph{sentido commum} não é distincto dos sentidos externos, pois, segundo elle, todo o sentido percebe as proprias percepções, distinguindo-as das dos outros sentidos.\par A theoria do \emph{sentido commun}, que os modernos costumam chamar \emph{sentido intimo} ou \emph{consciencia sensitiva} e que existe no homem e nos animaes mais perfeitos, foi completamente desvirtuada por A. \emph{Rosmini}.\par \emph{Rosmini} admitte um \emph{sentido intimo}; mas este sentido é, para elle, a percepção constante que a alma tem do corpo, como tal, --- percepção que, ao menos \emph{por natureza}, é anterior ás sensações externas, das quaes é fundamento e sujeito. Por isso, \emph{Rosmini} chama esse sentido \emph{fundamental}, e sustenta que as sensações dos sentidos externos são apenas modificações do sentido \emph{fundamental}. E não é tudo; elle chega a collocar no \emph{sentido fundamental} a união da alma com o corpo, e a propria natureza do sujeito sensiente.\par Esta theoria de \emph{Rosmini} é falsa pelas seguintes razões:\par a) \emph{Rosmini} afirma que o \emph{sentido fundamental} é \& constante percepção do proprio corpo, como tal. --- Ora esta affirmação é falsa; porque confunde a faculdade com o acto, e a percepção é acto.\par \emph{b}) \emph{Rosmini} opina que o \emph{sentido fundamental} percebe o corpo, como tal, \emph{independentemente das sensações}. --- Esta opinião é falsa. Com effeito, quando \textbar{} aquelle sentido percebe o corpo, ou percebe-o como \emph{vivo}, ou como não-vivo. Se o percebe como \emph{não-vivo}, não percebe o nosso corpo, pois este é vivo. Se o percebe como \emph{vivo}, percebe-o emquanto este é sujeito de sensações.\par \emph{c}) \emph{Rosmini} faz consistir a propria substancia do sujeito sensiente no\pgn{120}\emph{sentido fundamental}, e por isso numa \emph{sensação}; pois, para elle, o sentido \emph{fundamental} é a percepção do nosso corpo. --- Ora isto é inadmisslvel; porque toda a sensação, embora \emph{fundamental}, longe de constituir o sujeito sensiente, presuppõe-no; pois o acto presuppõe o agente.\par \emph{d}) O \emph{sentido fundamental} de \emph{Rosmini}, embora existisse, não seria comtudo, como vimos, a base da união da alma com o corpo.\par Advertimos que o \emph{sentido commum} póde chamar-se \emph{fundamental}, não em relação aos sentidos \emph{externos}, mas em relação aos outros sentidos \emph{internos}, que são a \emph{imaginação}, a \emph{força estimativa} e a \emph{memoria sensitiva}; pois o acto d'estas faculdades suppõe o do \emph{sentido commuwm}.}.

%%>>> t85

\pg{120} \tese{85}{Orgão do sentido commum.}
O \emph{sentido commum}, sendo uma faculdade \emph{material}, ou \emph{sensitiva}, deve residir, como em \emph{sujeito}, num orgão corporeo. Esse orgão é o \emph{cerebro}; pois é do cerebro que partem, e é a elle que convergem, como a centro \emph{commum}, o8 nervos sensitivos e por isso todas as sensações dos sentidos \emph{externos} \footnote{O \emph{sentido commum} está localisado nas cellulas cinzentas dos hemispherios cerebraes; pois estas colligem as sensações espalhadas na medulla.\par Este orgão central, para o qual, como para um centro commum, convergem todas as sensações externas, para ahi encontrarem a sua perfeição e o seu complemento, pode ser e é tambem a raiz de que parte para a peripheria toda a energia sensitiva, de modo que os nervos sensitivos dos orgãos externos não são que um prolongamento do orgão cerebral do sentido commum, e o cerebro é o termo e o principio de toda a sensação. Diz S. Thomaz: « Vis sentiendi diffunditur in organa quinque sensuum \emph{ab aliqua} una \emph{radice communi}, a qua procedit vis sentiendi in \emph{omnia} organa; ad quam etiam \emph{terminantur} omnes immutationes singulorum organorum ». (\emph{In} 3 \emph{de} Anima, lect. 3) --- E não é tudo. Este orgão central do sentido commum não só é um echo consciente de todas as sensações \emph{externas}, que elle ceniralisa e compara, mas é tambem a raiz e o orgão directo de todas as sensações \emph{internas}. O mesmo S. Doutor diz: « Sensus communis est \emph{radi} phaniasiae et memoriae, quae praesupponunt aétus sensus communis » (1. \emph{c}.), Ooncluimos que o sentido commum é un \emph{centro commum}, no qual, como em seu termo, se unem as sensações externas e internas, e é o \emph{principio provimo}, do qual, como de raiz, deriva toda a sensibilidade.\par Devemos advertir que as cellulas cinzentas dos hemispherios cerebraes, que são o Orgão central do \emph{sentido commnum}, ou da \emph{consciencia sensitiva}, são insensiveis ás proprias impressões, emquanto podem ser feridas, cauterisadas, torturadas, sem que experimentem a mais leve dór. Este phenomeno, que encheu de admiração os physiologistas de todos os tempos, explica-se facil e cabalmente na nossa doutrina. Por quanto, admittindo, como devemos admittir, que nenhuma faculdade organica póde perceber o proprio acto, segue-se que o \emph{sentido commum} é destituido de sensibilidade consciente.\par Da doutrina exposta resulta que, interrompida a communicação entre os nervos sensitivos dos sentidos externos e o cerebro, ou tornando-se o cerebro incapaz para receber as impressões que lhe transmittem os sentidos externos, o homem não advertirá na sensação, nem no seu objecto. Os drs. \emph{Lessana} e \emph{Lemoigne}, no seu livro --- \emph{Physiologia dos centros nervosos}, narram\pgn{121}o facto de uma pessoa, que, tendo os hemispherios cerebraes totalmente destruidos, via, cheirava e ouvia, mas não advertia nos objectos que via, nem nos cheiros, nem nos sons. Isto confirma a possibilidade e a existencia de sensações nos sentidos externos, não percebidas pelo sentido commum.}.

%%>>> t86

\pg{121} \tese{86}{Imaginação.}
\emph{Imaginação} é a faculdade pela qual o homem conserva e reproduz as \emph{imagens sensiveis} dos entes, mesmo na \emph{ausensia} d'elles. --- A \emph{imagem} de um ente é \emph{sensivel}, quando o representa cercado de propriedades materiaes; tal é á \emph{imagem} que representa um \emph{leão}, com tal côr, \emph{figura}, ete.

A existencia d'esta faculdade no homem é innegavel. Vemos, por ex., um cavallo, e, depois de o termos perdido de vista, ainda conservamos a sua imagem \emph{sensivel}, que nol-o representa, como se estivesse presente \footnote{São duas, pois, as funcções da \emph{imaginação}. A primeira é a de conservar as imagens sensiveis das coisas, que, em seguida á percepção dos sentidos externos, são \emph{formadas} nos lóbos occipitaes do cerebro, e depois \emph{coordenadas} pelo sentido commum, e finalmente \emph{recebidas} na imaginação. Por isso esta faculdade chama-se o \emph{thesouro das especies} ou \emph{imagens sensiveis}. Diz S. Thomaz: « Ad harum formarum (scilicet sensibilium) \emph{retentionem} et conservationem ordinatur \emph{phantasia} sivo \emph{imaginatio}, quasi \emph{tiesaurus} quidam formarum persensum acceptarum » (\emph{Sum}. \emph{Theol}., p. I, q. (8, a. 4). --- A segunda funcção é a de \emph{reproduzir} essas imagens sensiveis das coisas, mesmo quando estas estão ausentes; e neste caso a imaginação diz-se \emph{reproductiva}.\par A \emph{imaginação} costuma chamar-se tambem \emph{phrantasia}, porque \emph{faz apparecer} os objectos, que estão ausentes; e as \emph{imagens sensiveis} denominam-se tambem \emph{phantasmas}, porque são representações sensiveis dos objectos materiaes ausentes, precedentemente percebidos.\par Pergunta-se: Qual é a causa, que move ao acto a \emph{imaginação} ?\par Quando o objecto é \emph{actualmente} percebido pelos sentidos externos, \& causa, que excita e move ao acto a \emph{imaginação}, é o \emph{sentido commum} ; pois é este que, recebendo as imagens sensiveis dos objectos externos, determina a \emph{imaginação} a recebel-as, conserval-as e reproduzil-as.\par Quando o objecto esta ausente, a causa, que excita e move a \emph{imaginação} a reproduzir a imagem d'aquelle objecto, é duplice: o \emph{movimento} nervoso e a \emph{associação das imagens}. --- Digamos alguma coisa ácerca d'estas causas.\par O \emph{movimento nervoso} póde excitar e mover a \emph{imaginação} a reproduzir as \emph{imagens} dos objectos percebidos. Com effeito, esta faculdade, sendo sensitiva, serve-se, nas suas operações, do cerebro como de instrumento. Por isso, assim como, se pela sensação se produzir no cerebro um movimento nervoso, a \emph{imaginação receberá} a imagem dos objectos \emph{presentes}; assim tambem, se no cerebro se reproduzir um movimento nervoso similhante ao primeiro, a \emph{imaginação reproduzirá} a mesma imagem dos objectos, ainda que \emph{ausentes}.\par A \emph{associação das imagens} tambem póde excitar e mover a \emph{imaginação} a \emph{reproduzir} as imagens dos objectos percebidos, mas ausentes. Na verdade, as imagens dos objectos \emph{sémilhantes}, ou \emph{proximos}, ou \emph{simultaneamente} percebidos\pgn{122}, estão por tal forma ligadas entre si, que são como outros tantos anneis para formar uma cadeia. D'ahi se segue que a reproduccüo de uma imagem traz comsigo a reproduccüo de muitas outras.\par A. \emph{Bain} reduz a \emph{tres} as leis, que regem os phenomenos da \emph{associação das imagens}: --- uma de \emph{contiguidade}, --- outra de \emph{semelhança}, --- e a terceira de \emph{composição}. --- Estas leis regulam tambem o exercicio da \emph{memoria}.\par L\textasciicircum{} Ler. Os actos, as emoções e as sensações produzidas simultaneamente ou num pequeno intervallo, tendem para formar um todo coherente; por forma que, se mais tarde se reproduzir um d'esses estados, os outros tendem para se reproduzirem. --- E' a lei da \emph{contiguidade}.\par IL* Ler, Os actos, as sensações, e as emoções \emph{actuaes} tendem para reproduzir as impressões \emph{passadas}, ou os estados \emph{anteriores} da consciencia, que lhe são \emph{semelhantes}. --- E' a lei da \emph{semelhança}.\par IIL* Ler. Os estados \emph{passados} da consciencia, as sensações, os actos e as emoções, reproduzem-se mais facilmente, quando se relacionam, \emph{por varios motivos}, com \& impressão do momento presente. --- E' a lei da \emph{composição}.\par A. \emph{hypothese phijsiologica}, que melhor explica o phenomeno da \emph{associação das imagens}, póde resumir-se na seguinte formula: Quando dois ou mais processos nervosos se realizaram conjunctamente (no mesmo tempo ou numa successão immediata), a causa, que reproduz \emph{directamente} um d'aquelles processos, tende para reproduzir \emph{indirectamente} os outros.}.

%%>>> t87

\tese{87}{Objecto da imaginação.}
O objecto da \emph{imaginação} é constituido pelas \emph{imagens sensiveis} dos entes, percebidos \emph{immediatamente} \pg{122}  pelos sentidos \emph{externos}, e \emph{mediatamente} pelo \emph{sentido commum}. --- A \emph{imaginação}, pois, embora distineta dos sentidos \emph{externos} e do \emph{sentido commum}, todavia suppõe o exercicio d'aquelles e d'este; e o homem, destituido de um sentido \emph{externo}, não póde reter nem reproduzir a \emph{imagem} do objecto \emph{proprio} d'esse sentido \footnote{A. \emph{imaginação} e o \emph{sentido commu}. differem no \emph{objecto}, no \emph{modo de operar}, e nos \emph{effeitos}. --- Differem no \emph{objecto}; pois que a \emph{imaginação} recebe, conserva e reproduz as \emph{imagens} dos objectos, e o \emph{sentido conmmunm} percebe as \emph{sensações} dos sentidos \emph{externos}. --- Differem no \emph{modo de operar}; porque a \emph{imaginação} produz o seu acto tambem na ausencia dos objectos, ao passo que o \emph{sentido communm} só percebe as sensações externas no tempo, em que se realizam nos sentidos. --- Differem nos \emph{effeitos}; pois a \emph{imaginação}, por não ser determinada ao acto pelos proprios objectos, está sujeita a erros, o que não se dá com o \emph{sentido commuwm}, que pelas sensações attinge os proprios objectos.\par O exercicio da \emph{imaginação} suppõe o exercicio dos sentidos \emph{externos}. Na verdade, a \emph{imaginação} é uma faculdade, pela qual nos representamos as coisas externas e materiaes, quando ellas estão ausentes. Ora só podemos representar-nos o que foi percebido por \emph{algum} sentido externo. --- A experiencia não só attesta que a nossa \emph{imaginação} não póde representar senão o que foi percebido pelos sentidos externos, mas tambem que uma pessoa, a quem falte um sentido, não tem a imagem dos objectos proprios d'esse sentido; assim um cego de nascimento não tem a imagem das \emph{córes}.\par Objectam: A's vezes nos representamos com a \emph{imaginação} muitas coisas que --- ou não existem na realidade, por ex., um \emph{cavallo com asas}, --- ou são meramente espirituaes e escapam aos sentidos, como DEkus, a \emph{alma}, etc.\par Emquanto á primeira parte da difficuldade respondemos que, se na realidade\pgn{123} não existe um \emph{cavallo com} azas, existem comtudo \emph{cavallos}, e existem \emph{azas}, --- objectos distinctos, que depois unimos em virtude da imaginação \emph{creadora}. --- Emquanto á segunda parte respondemos que, quando nos representamos com a \emph{imaginação} objectos meramente espirituaes, não os representamos sob uma forma sensivel, que lhes seja propria, pois esses objectos não têm forma sensivel; mas applicamos-lhes uma forma sensivel, que é propria dos objectos materiaes, e que foi percebida por algum sentido.\par \emph{Reid} diz que só podemos representar-nos com a \emph{imaginação} os objectos materiaes, percebidos pela \emph{vista}. --- Concedemos que a \emph{imagem} dos objectos percebidos pela \emph{vista} é mais viva que a dos objectos percebidos pelos outros sentidos, pois a \emph{vista} é o mais nobre dos sentidos externos; mas negamos que não possa formar-se a imagem de objectos percebidos pelos outros sentidos, pois a experiencia attesta que temos a imagem dos \emph{sons}, dos \emph{sabores}, dos \emph{cheiros}, etc., --- imagem que reproduzimos na vigilia e nos sonhos.\par Pergunta-se: Como é que as \emph{imagens} se distinguem das \emph{sensações} ?\par Respondemos que as \emph{imagens} dos objectos \emph{ausentes} se distinguem das \emph{sensações} dos objectos \emph{presentes}, emquanto aquellas são muito menos vivas do que estas. --- Todavia, se o exercicio dos sentidos \emph{externos} ficar suspenso, como no somno, só dominará a \emph{imaginação}, e confundir-se-ha o mundo phantastico interno com o mundo real externo; pois, nesse caso, falta a sensação, com que se possam comparar as imagens, para as distinguirmos da actual percepção do objecto. Como tambem confundir-se-ha o mundo real com o phantastico, se as 2\emph{magens} se tornarem tão vivas como as sensações actuaes, o que se verifica nos lodcos e nos que estão debaixo da impressão de um grande medo: d'ahi as illusões e as allucinações.}.

%%>>> t88

\pg{123} \tese{88}{A imaginação é uma faculdade sensitiva.}

a) Toda a faculdade, eujo objecto é um ente \emph{sensivel}, não é \emph{intellectual}, mas \emph{sensitiva}; pois que os entes \emph{sensiveis} não constituem o objecto das faculdades 2\emph{ntellectuaes}, mas das \emph{sensitivas}. Ora a \emph{imaginação} é a faculdade, cujo objecto é um ente \emph{sensivel}, isto é, cercado de condições materiaes ou propriedades singulares. Logo a \emph{imaginação} é uma faculdade \emph{sensitiva}.

b) Uma faculdade, que existe nos animaes, não póde deixar de ser \emph{sensitiva}; pois os animaes são destituidos de faculdades intellectuaes. Ora a \emph{imaginação} existe nos animaes; pois elles conservam as \emph{imagens} do \emph{dono}, da \emph{casa}, do \emph{ninho} ou \emph{covil}, ete., e em virtude d'estas imagens tendem para os objectos representados, ainda quando estes estão longe ou já não existem. Logo a \emph{imaginação} é uma faculdade \emph{sensitiva} \footnote{Acerescentamos outros argumentos em confirmação da these.\par \emph{a}) O exercicio da \emph{imaginação} é muito vivo e activo nos que sonham, nos que tém febre e nos loucos, que todavia são destituidos do exercicio da razão; pois elles, não podendo comparar a imagem com o objecto, trocam este por aquella, e tomam a \emph{sombra} pela realidade. Ora isto mostra que a \emph{imaginação} é uma faculdade \emph{sensitiva}.\par \emph{b}) As edades exercem grande influencia sobre a \emph{imaginação}. Esta faculdade\pgn{124}, que é viva e brilhante nas creanças e nos jovens, é geralmente languida nos homens de edade madura e nos velhos; ao passo que a intelligencia, que é fraca e pouco desenvolvida nas creanças e nos jovens, é prompta, vigorosa e penetrante nos homens de edade madura e nos velhos. Ora isto mostra que a \emph{imaginação} não é uma faculdade \emph{intellectual}.}.

%%>>> t89

\pg{124} \tese{89}{A imaginação no homem e no animal.}
A \emph{imaginação}, embora seja uma faculdade commum ao homem e ao animal, todavia naquelle é mais nobre e elevada do que neste. Porquanto, a \emph{imaginação} no animal é somente \emph{reproductiva}, pois que simplesmente apresenta as \emph{imagens} dos objectos externos, que foram percebidos (assim o animal reproduz a imagem de um \emph{monte} que viu); ao passo que no homem a \emph{imaginação}, por estar subordinada á direcção da intelligencia e ao imperio da vontade, não é só \emph{reproductiva}, mas é tambem \emph{creadora}, porque une ou compõe varias imagens de objectos percebidos, para formar a imagem de um objecto que nunca foi percebido e que talvez nem exista (assim, tendo as imagens do \emph{monte} e do 0\%\emph{ro}, formamos a imagem de um \emph{monte d'oiro}) \footnote{Alguns philosophos pensam que a \emph{imaginação reproductiva} e a \emph{imaginação creadora} não são duas funcções de uma e a mesma faculdade, mas que são duas faculdades distinctas.\par Esta opinião é inacceitavel. A imaginação \emph{reproductiva} e a .\emph{creadora} seriam duas faculdades distinctas, se os seus \emph{objectos} fossem \emph{especificamente} diversos; pois as faculdades distinguem-se pelos seus objectos. Ora o objecto da \emph{imaginação reproductiva} não é \emph{especificamente} diverso do objecto da \emph{imaginação creadora}, mas identifica-se com elle; pois o objecto d'uma e d' outra é constituido pelas \emph{imagens sensiveis} das coisas. --- Se a \emph{imaginação} no homem é tambem \emph{creadora}, este dote não resulta da \emph{imaginação}, como tal, mas da união d'ella com a intelligencia e com a vontade.\par A \emph{imaginação}, no homem, embora seja uma faculdade \emph{sensitiva}, esta tão intimamente ligada com a intelligencia, que uma acompanha sempre outra. --- Assim como toda a fórma sensivel da imaginação desperta na intelligencia uma idea; assim tambem não ha idea da intelligencia, que a imaginação não revista de uma fórma sensivel. Por isso, a imaginação é como um espelho em que a intelligencia se reflecte, para extrahir da imagem sensivel o conceito espiritual, e tambem para dar ao conceito espiritual uma, fórma sensivel. E nada mais natural. A alma humana, por estar unida ao corpo, não póde deixar de acceitar, nos seus actos, o concurso dos sentidos. Por isso, quando queremos formar a imagem dos seres immateriaes, como são \textsc{Deus}, os \emph{Anjos}, a nossa \emph{alma}, costumamos vestil-os de uma figura e fórma sensivel, embora elles não tenham figura nem forma sensivel.\par São varias e importantes as vantagens da \emph{imaginação}. - Eis as principaes:\par 1.2) A \emph{imaginação} subministra á intelligencia os materiaes, em que esta faculdade opera e de que abstrahe as imagens \emph{immateriaes} dos objectos, chamadas \emph{especies intelligiveis}; pois não podemos entender senão o que nos foi apresentado pela \emph{imaginação}.\par 2.º) Pela força da \emph{imaginação}, vestimos de formas sensiveis as ideas\pgn{125}abstractas e geraes, tornando-as d'este modo mais proporcionadas á indole\par n da nossa intelligencia, que apprehende com maior facilidade os conceitos universaes, apresentados sob uma imagem sensivel.\par 8.º) A \emph{imaginação} exerce um influxo poderoso e benefico nas bellas artes. Na verdade, todo o empenho das bellas artes consiste em representar o ideal sob uma fórma sensivel. Ora, se o \emph{ideal} é um producto da intelligencia, a \emph{fórma sensivel} é um producto da \emph{imaginação}.\par Todavia, se são grandes as vantagens da \emph{imaginação}, tambem são grandes os damnos, que podem resultar do seu exercicio. Na verdade, se não fôr guiada pela razão, \& \emph{imaginação} --- perturba o espirito, afastando-o das occupações serias e da inquirição da verdade, --- inclina ao amor das coisas vãs e causa fastio das coisas espirituaes, --- póde levar-nos a julgar como realidade o que não passa de uma ficção: d'ahi as allucinações.\par Um dos mais tristes phenomenos, provenientes da alteração da imaginação, é a \emph{loucura}. Como vimos e veremos ainda, o acto da intelligencia suppõe e exige, no estado actual, o acto da imaginação; por isso um desarranjo na imaginação produz necessariamente, e por reflexo, um desarranjo na intelligencia. Dá-se o desarranjo na imaginação, --- ou quando a substancia cerebral se acha affectada por inflammação, por tumores, ou pela acção do alcool, --- ou quando a circulação do sangue no cerebro soffre um augmento demasiado ou uma demasiada diminuição, --- ou quando se dá uma inflammação nas meninges (\emph{meningite}), --- ou quando se torna irregular a circulação do agente nervoso no cerebro, como nas paixões vehementes e persistentes. Para esta perturbação do agente nervoso concorre muito a demasiada excitabilidade cerebral, chamada \emph{nervosismo} ou \emph{neurastenia}, que uns herdam dos paes e outros adquirem pelo excesso de trabalho physico, intellectuale moral. --- A loucura é \emph{total} ou \emph{parcial}, conforme é \emph{total} ou \emph{parcial} o desarranjo da intelligencia. --- Na loucura \emph{parcial}, que se chama \emph{monomania}, dá-se o desarranjo só numa certa ordem de imagens, e por isso na ordem correspondente de ideas. Estas imagens e estas ideas se chamam \emph{fixas} e tomam um predominio excessivo e exclusivo, de modo que não deixam formar outras imagens nem outras ideas, senão as que com ellas tenham relação. --- São quatro as formas mais frequentes da loucura \emph{parcial}: monomania \emph{das grandezas}, monomania \emph{religiosa}, monomania de \emph{perseguição}, monomania \emph{moral} ou \emph{impulsiva}. --- O desarranjo da intelligencia, na monomania, é, ás vezes, tão parcial e tão disfarçado, que só o pode reconhecer e julgar o medico alienista e, ás vezes, só depois de uma longa e accurada analyse de todos os actos do doente, pois dá-se uma especie de loucura \emph{que raciocina}.}.

%%>>> t90

\pg{125} \tese{90}{Orgão da imaginação.}
A \emph{imaginação}, sendo uma faculdade \emph{sensitiva}, deve residir num certo e determinado orgão. Esse orgão é o \emph{cerebro}; pois está averiguado que, por uma lesão do \emph{cerebro}, o exercicio da imaginação ou se torna anormal, ou cessa inteiramente, --- e, vice-versa, por um longo e intenso exercicio da \emph{imaginação}, o \emph{cerebro} se eanca \footnote{O orgão da \emph{imaginação} é o mesmo que o do \emph{sentido commum} --- a \emph{substancia cinzenta dos hemispherios cerebraes}. Na verdade, sendo a \emph{substancia cinsenta dos hemispherios cerebraes} uma continuação das innumeraveis fibras, de que é composto cada um dos nervos sensitivos e motores, é claro que\pgn{126}toda a mudança, que se verifique naquellas fibras, encontrará um centro de repercussão nas cellulas cerebraes. D'ahi a existencia e as funcções do \emph{sentido commum}. --- Mas não basta. As cellulas cerebraes, diversamente excitadas segundo a diversidade das sensações que se realizam nos sentidos externos e se reflectem nos emispherios, apresentam as \emph{imagens} dos objectos, que foram causa d'aquella excitação. Estas imagens existem nos hemispherios cerebraes num estado latente, num estado habitual; existem, segundo \& expressio de Wundt (\emph{Principios de Psychologia physiologica}), como meras disposições cerebraes, até que ou uma nova sensação torne a excitar os pontos cerebraes que constituem alguma d'aquellas imagens, ou uma onda nervosa, proveniente de outros pontos do cerebro, faça reviver a imagem toda, em que se veja o objecto, embora ausente, das sensações passadas.\par Ninguem se admire de que o cerebro possa ser a séde de innumerayeis imagens, embora o seu volume não seja extraordinariamente grande. Por quanto, o numero das cellulas cerebraes nervosas é enorme (alguns dizem que é superior a mil milhões). - Além d'isso, uma e a mesma cellula cerebral póde entrar, como parte integrante, na formação de muitas e diversas imagens; pois, combinando-se com certas e determinadas cellulas, fórma uma imagem, e, combinando-se com outras cellulas, fórma uma outra imagem, distincta e diversa da primeira. D'este modo, poucas cellulas, diversamente combinadas, formam uma innumeravel multidão de imagens; assim como poucas lettras, diversamente combinadas, formam innumeraveis palavras.}.

%%>>> t91

\pg{126} \tese{91}{Estimativa.}
\emph{Estimativa} é a faculdade pela qual o homem percebe nos objectos materiaes certas propriedades, que escapam aos sentidos externos; taes são as propriedades de \emph{vantagem} ou de \emph{damno}, de \emph{bem} ou de \emph{mal sensivel}, ete. --- A \emph{estimativa} chama-se tambem \emph{juizo énstinctivo}.

Esta faculdade existe no homem; pois que julgamos que uma coisa \emph{particular} e \emph{concreta} nos é \emph{util} ou \emph{nociva}, e este juizo \emph{particular} e \emph{concreto}, que muitas vezes precede o juizo \emph{comparativo} e \emph{reflectivo}, deriva da \emph{estimativa} \footnote{O homem deve ser dotado de \emph{estimativa}, porque elle é tambem animal,e animal perfeito, e por isso deve ter todas as faculdades proprias dos animaes perfeitos. E que elle possua esta faculdade, ve-se pelos juizos \emph{instinctivos}, que precedem, acompanham ou seguem os juizos da razão. --- Não se pode negar que a \emph{estimativa} é menos necessaria ao homem do que aos animaes irracionaes; mas devemos tambem reconhecer que ella é muito util e necessaria na infancia, nas perturbações da razão, e em tantos movimentos que precedem a reflexão e que são regulados por esta faculdade.}.

%%>>> t92

\tese{92}{Objecto da estimativa.}
O objecto da \emph{estimativa} é constituido pelas propriedades de \emph{vantagem} ou de \emph{damno}, etc., que existem nos seres materiaes e escapam ao alcance dos outros sentidos. --- Por isso a \emph{estimativa} é uma faculdade distineta da \emph{imaginação}, e das outras faculdades sensitivas \footnote{\emph{A estimativa} é uma faculdade distincta e diversa da \emph{imaginação}. Na verdade, a \emph{imaginação} tem por objecto as \emph{imagens} das coisas percebidas por intermedio dos sentidos \emph{externos}; ao passo que a \emph{estimativa} tem por objecto\pgn{127}certas propriedades de \emph{vantagem} ou de \emph{damno}, que não foram nem podem ser percebidas por nenhum sentido externo.\par Todavia, podemos affirmar que a \emph{estimativa}, conservadas as devidas proporções, se refere á \emph{imaginação}, como a razão \emph{pratica} se refere á razão \emph{especulativa}. Na verdade, a razão \emph{especulativa} contempla o objecto, só emquanto elle é \emph{verdadeiro} e sem relação ás operações; ao passo que a razão \emph{pratica} contempla o objecto, mas em ordem ás operações. O mesmo acontece com relação á \emph{imaginação} e \& \emph{estimativa}. \emph{A imaginação} conserva as \emph{imagem} dos objectos, percebidos primeiramente pelos \emph{sentidos externos} e depois pelo \emph{sentido commum}, mas sem as relacionar com as operações; ao passo que a \emph{estimativa}, conhecendo nos objectos as propriedades de \emph{vantagem} ou de \emph{damno}, dirige os seus conhecimentos para as operações.\par A \emph{estimativa} não só é distincta da \emph{imaginação}, mas tambem de todas as outras faculdades sensitivas. Na verdade, o objecto da \emph{estimativa}, que é constituido pelas propriedades concretas de \emph{vantagem} ou de \emph{damno}, existentes nos objectos materiaes, é distincto e diverso do objecto de todas as outras faculdades; pois os \emph{sentidos externos}. percebem as qualidades \emph{sensiveis} das coisas, e o \emph{sentido communm} recolhe as \emph{sensações} dos sentidos externos. Ora as faculdades conhecem-se e distinguem-se pelos seus \emph{objectos}. Logo a \emph{estimativa} é distincta das outras faculdades.}.

%%>>> t93

\pg{127} \tese{93}{A estimativa é uma faculdade sensitiva.}

a) Toda a faculdade perceptiva, cujo objecto é \emph{material} e \emph{singular}, não é \emph{intellectual}, mas \emph{sensitiva}; pois o objecto de uma faculdade \emph{intellectual} é \emph{immaterial} e \emph{universal}. Ora o objecto da \emph{estimativa é material} e \emph{singular}; pois percebemos que uma coisa \emph{material} e \emph{singular é util} ou \emph{nociva}, por ex., \emph{esta herva é-me util}. Logo a \emph{estimativa} é uma faculdade \emph{sensitiva}.

b) Uma faculdade, que se encontra tambem no animal, não é \emph{intellectual}, mas \emph{sensitiva}; pois que o animal é destituido de faculdades intellectuaes. Ora a \emph{estimativa} encontra-se tambem no animal; pois o animal procura certas coisas, e foge de outras, não porque umas são agradaveis e outras desagradaveis, mas porque umas lhe são \emph{uteis} e outras \emph{nocivas}; assim as ovelhas fogem do lobo, não porque é deforme, mas porque é seu

4 inimigo, --- e as aves recolhem as palhas, não porque são agradaveis á vista, mas porque são \emph{uteis} para a construeção do ninho. Logo a \emph{estimativa} é uma faculdade \emph{sensitiva} \footnote{A \emph{estimativa} não é uma faculdade \emph{intellectual}, mas \emph{sensitiva}; poisque os juizos, que ella fórma, não são \emph{universaes} e \emph{abstractos}, mas \emph{singulares} e \emph{concretos}, e o \emph{singular} ou \emph{concreto} não é objecto de uma faculdade \emph{intellectual}, mas de uma faculdade \emph{sensitiva}.\par Nem se diga que tambem a intelligencia fórma alguns juízos \emph{séngulaves} e \emph{concretos}, que tratam de objectos \emph{singulares} e \emph{concretos}; pois que estes juizos, formados pela intelligencia, são \emph{especificamente} diversos dos que fórma a \emph{estimativa}. Com «feito, nos juizos \emph{singulares} e \emph{concretos} da intelligencia, esta conhece a \emph{essencia} do sujeito e da propriedade; ao passo que, nos juizos\pgn{128}\emph{singulares} e \emph{concretos} da \emph{estimativa}, esta, percebe que uma certa e determinada propriedade se encontra num sujeito, sem comtudo conhecer a \emph{essencia} da propriedade nem do sujeito.\par Dizem: O objecto não percebido pelos sentidos externos deve ser percebido pela intelligencia. Ora o objecto, que se diz proprio da \emph{estimativa}, não é percebidos pelos sentidos externos. Logo deve ser percebido pela intelligencia.\par Hespondemos que o objecto, que nao se percebe pelos sentidos externos, deve ser percebido pela intelligencia, quando esse objecto é \emph{immaterial} e \emph{universal}, mas não quando é \emph{material} e \emph{singular}. Ora o objecto proprio da \emph{estimativa é material} e \emph{singular}. Logo não é percebido pela intelligencia, mas por uma faculdade sensitiva. Esta faculdade \emph{sensitiva} não é externa. Logo é interna. --- Como se vó, os adversarios partem do falso principio de que entre os sentidos externos e a intelligencia não existem outras faculdades perceptivas, distinctas d'esta e d'aquelles.\par insistem: Os animaes formam aquelles juizos por meio de um \emph{nstincto}, e não por meio da \emph{estimativa}.\par O que entendem os adversarios por \emph{instincto}? Entendem --- uma tendencia necessaria para o objecto, sem algum previo conhecimento, --- ou uma tendencia necessaria para um objecto, determinada por uma percepção antecedente, embora não \emph{intellectual}? --- Na primeira hypothese, enunciam uma falsidade; porque o acto \emph{énstinctivo}, embora não seja livre, suppõe o exercicio da faculdade perceptiva; assim o cordeiro foge \emph{instintivamente} do lobo, \emph{mas depois de o ter percebido como seu inimigo}. --- Na segunda hypothese, reconhecem no animal a \emph{estimativa}, embora lhe dêm outro nome; pois aquella percepção, antecedente e não \emph{intellectual}, que determina a tendencia necessaria do animal para o objecto, é effeito da \emph{estimativa}.\par Se dizemos que o \emph{instincto} é \emph{cego}, esta expressão não significa que o instincto prescinde de \emph{todo} o conhecimento, mas significa que o animal não conhece o \emph{fim} das suas operações, nem a \emph{razão} dos seus movimentos, isto é, á \emph{proporção} entre os \emph{meios} que emprega e o \emph{fim} que tem em vista.}.

%%>>> t94

\pg{128} \tese{94}{A estimativa no homem e no animal.}

A \emph{estimativa}, embora seja commum ao homem e aos animaes mais perfeitos, todavia naquelle é mais nobre do que nestes. Com effeito, os juizos, que os animaes formam pela \emph{estimativa}, não resultam da comparação dos termos, mas do instincto da natureza, e por isso não se chamam \emph{comparativos}, mas \emph{tinstinctivos}; pois os juizos \emph{comparativos} são effeitos da intelligencia, e não dos sentidos. Pelo contrario, os juizos, que o homem fórma pela \emph{estimativa}, embora possam por um momento ser \emph{énstinctivos}, tornam-se depois \emph{comparativos}, e isto por causa da \emph{razão}, de que o homem é dotado e com que a \emph{estimativa} está unida. Por isso a \emph{estimativa}, como existe no homem, diz-se propriamente \emph{cogitativa}, ou \emph{razão particular} \footnote{A \emph{estimativa}, como existe no homem, chama-se razão \emph{particular}, porque compara termos \emph{particulares}, distinguindo-se assim da razão \emph{universal} ou \emph{intellectual}, que compara termos \emph{universaes}.\par A \emph{estimativa}, de que o homem é dotado, não é \emph{NT s} diversa\pgn{129}da \emph{estimativa}, que existe no animal; comtudo em nós é mais nobre do que no animal. Na verdade, o animal, guiado pela \emph{estimativa}, percebe \emph{énstinctivamente} as propriedades de \emph{vantagem} ou de \emph{damno}, de \emph{bem} ou de \emph{mal sensivel}, que existem num objecto singular e \emph{material}, sem comtudo perceber a \emph{natureza} do mesmo objecto; assim a ovelha \emph{nstinctivamente} percebe esta \emph{herva}, não emquanto é \emph{Aerva}, mas emquanto é seu \emph{alimento}, e \emph{instinctivamente} conhece este cordeiro, não emquanto é este cordeiro, mas emquanto lhe é \emph{util} para se desembaraçar do leite. Pelo contrario, \emph{o} homem, pela \emph{estimativa}, percebe que um objecto particular Ihe é \emph{util} ou \emph{nocivo}, \emph{e} tambem, pela \emph{razão}, percebe a \emph{natureza} d'aquelle objecto, por ex., da \emph{Aerva}, e, formando um juizo \emph{comparativo}, diz: \emph{esta herva é-me util}, --- D este modo os juizos da \emph{estimativa} são mais perfeitos e completos no homem do que no animal (\emph{S}. Th. \emph{De anima}, lib. II, lect. 18). --- Objectam: Se a \emph{estimativa} do homem percebe a \emph{essencia} dos objectos, não é faculdade \emph{sensitiva}, mas \emph{intellectual}; pois só uma faculdade \emph{intellectual} póde perceber a \emph{natureza} ou \emph{essencia} das coisas.\par Respondemos que a \emph{estimativa} do homem é sempre uma faculdade \emph{sensitiva}; porque, se o homem conhece a \emph{essencia} das coisas, que são \emph{uteis} ou \emph{nocivas}, conhece-a --- não emquanto é dotado de \emph{estimativa}, pois a \emph{estimativa}, \emph{por si}, não percebe \emph{essencia} alguma, --- mas emquanto a sua \emph{estimativa está unida com a razão}, isto é, emquanto é dotado tambem de \emph{razão}, que completa e aperfeiçõa o juizo da \emph{estimativa}.\par Diz S. Thomaz: «lllam eminentiam habet \emph{cogitativa}, et \emph{memorativa}, in homine, non per id quod est proprium sensitivae partis, sed per aliquam affinitatem et propinquitatem ad rationem universalem, \emph{secundum quamdam refluentiam}. (\emph{Sum}. \emph{Theol} p. L, q. 78, a. 4 ad 5). \emph{E} nada mais natural. Todas as faculdades têm a raiz na unica essencia da alma, e por isso deve haver uma \emph{refluencia}, uma redundancia das faculdades superiores, que são as intellectuaes, nas inferiores, que são as sensitivas, e estas devem estar mais ou menos subordinadas á influencia da intelligencia, conforme a sua maior ou menor independencia da materia. A intelligencia conhece as razões universaes, os principios abstractos, mas ella não poderia applicar essas razões e esses principios aos casos concretos ou ás coisas singulares, e por isso não poderia regular os nossos actos particulares, se não existisse uma faculdade capaz de perceber esses actos; como tambem esta faculdade não poderia dirigir os actos particulares, se não fosse dirigida ella mesma pelas razões universaes, conhecidas pela intelligencia. Oucamos o angelico Mestre: « A intelligencia tem ingerencia nas coisas singulares mesmo emquanto dirige as potencias inferiores por meio da \emph{razão particular}; \& qual é uma potencia individual, que se chama \emph{cogitativa} e tem o seu orgão no cerebro. Ora um principio universal, que a intelligencia tenha formado com relação a coisas praticas, não se pode applicar a um caso particular, senão por meio de uma faculdade perceptiva do singular; de modo que d'ahi resulte como um syllogismo, cuja premissa maior seja universal, \emph{a} menor seja uma proposição particular, que contenha a applicação feita pela razão \emph{particular}, e \& con-\par 9 : Vor. IL\pgn{130}clusüo seja a eleição do que se ha de praticar » (QQ. DD. \emph{De verit}., q. 10, a. 5). Assim a intelligencia conhece este principio universal: \emph{os paes merecem a nossa veneração}; mas este principio não influirá na nossa vida pratica, emquanto não reconhecermos por nossos paes estas determinadas pessoas. Se ao primeiro conhecimento, que é da intelligencia e se refere.a \emph{razões uni}- '\emph{versaes}, se unir o segundo, que é da \emph{cogitativa} e se refere a \emph{razões particulares}, é claro que ha tudo o que é necessario para a vontade proceder ao acto da eleição. E, pois, evidente a utilidade e.a necessidade da \emph{estimativa} ou \emph{cogitativa} para a nossa vida pratica. Se repararmos no que dizem os homens, veremos que elles sempre, ou quasi sempre, exprimem os pensamentos da \emph{cogitativa}, e só raras vezes os da \emph{intelligencia}. Mas não devemos julgar que a intelligencia fique em ocio. Os conceitos singulares suppõem os universaes; a cogitativa não pode trabalhar, se não trabalhar a intelligencia; é pelo trabalho combinado e harmonico das duas faculdades que se chega a um e o mesmo resultado.}.

%%>>> t95

\tese{95}{Orgão da estimativa.}
A \emph{estimativa}, sendo uma faculdade \emph{sensitiva}, reside num orgão. Este orgão é o \emph{cerebro}; pois é no cerebro que existem as imagens sensiveis dos objectos \pg{129} \pg{130}  externos, em que a \emph{estimativa} apprehende as propriedades de \emph{vantagem} ou de \emph{damno}, de \emph{bem} ou de \emph{mal sensivel}, ete. \footnote{Cf. Ol. Bernard, \emph{Lecons sur les anestesiques et Pasphyxwie}; Paris, 1875. --- Ch. \emph{Des fonctions du cerveau}.}.

%%>>> t96

\tese{96}{Memoria sensitiva.}
\emph{Memoria sensitiva} é a faculdade pela qual o homem conserva as \emph{proprias} percepções sensitivas \emph{passadas}. --- Esta faculdade existe no homem, como o prova o testemunho da consciencia \footnote{As percepções sensitivas \emph{passadas}, \emph{proprias} de um sujeito, podem considerar-se --- ou de um modo \emph{directo}, emquanto constituem um facto \emph{intimo} e \emph{concreto}, --- ou de um modo \emph{reflexo}, emquanto se consideram como \emph{proprias} e como \emph{passadas}. Ora a memoria \emph{sensitiva}, se pode recordar e recorda essas percepções, emquanto ellas constituem um facto \emph{intimo} e \emph{concreto}, não as pode recordar emquanto se consideram como \emph{proprias} e como \emph{passadas}. Por quanto, para isso, seria necessario que a faculdade \emph{reflectisse} sobre si mesma, e se percebesse como principio d'essas percepções, e possuisse a idea \emph{universal} do \emph{passado}; ora uma faculdade \emph{sensitiva} não reflecte, nem tem ideas universaes. Reconhecer as proprias percepções passadas, como \emph{proprias} e como \emph{passadas}, \& privilegio exclusivo da \emph{intelligencia}.}.

%%>>> t97

\tese{97}{Objecto da memoria sensitiva.}
O objecto da \emph{memoria sensitiva} é duplice: \emph{directo} e \emph{indirecto}. O \emph{directo} é constituido pelas percepções sensitivas \emph{passadas}; o \emph{indirecto} a formado pelos \emph{objectos} d'aquellas -percepções. --- Por isso a \emph{memoria sensitiva} é uma faculdade distineta do \emph{sentido commum}, da \emph{imaginação} e da \emph{estimativa} \footnote{O \emph{objecto} da memoria \emph{sensitiva} é constituido pelas sensações \emph{passadas} e pelos objectos \emph{passados} d'essas sensações. Aquellas são o objecto \emph{directo}; estes, o \emph{indirecto}. Diz S. Thomaz: « O \emph{passado} pode referir-se a duas coisas: ao \emph{objecto}, que é percebido, e ao «\emph{cto} da percepção, e estas duas coisas acham-se reunidas na faculdade sensitiva, que percebe uma coisa sensivel emquanto é por ella actualmente transmutada. Por isso o animal lembra-se, ao mesmo tempo, de que experimentou no passado uma sensação, e de que sentiu uma\pgn{131}coisa sensivel, que tambem passou » (\emph{Sum}. \emph{Theol}. p. I, q. 79, a. 6 ad 2) --- Quando dizemos que o objecto da memoria \emph{sensitiva} é constituido pelos \emph{objectos passados} das sensações, não entendemos que esta faculdade só se lembra dos objectos que \emph{passaram}, isto é, que já não existem na realidade; mas entendemos que se lembra dos objectos que \emph{passaram}, isto é, \emph{que foram anteriormente percebidos}, quer existam ainda na realidade, quer jà não existam (Cf. S. Thom., \emph{De mem}. \emph{ét rem}, 1. I.).\par D'ahi resulta que a memoria \emph{sensitiva} se distingue realmente dos outros sentidos \emph{internos}. Distingue-se do \emph{sentido commum} e da \emph{imaginação}; porque o \emph{sentido commum} percebe as sensações e os objectos \emph{presentes}, e a \emph{imaginação} percebe as sensações e os objectos \emph{ausentes}, ao passo que a memoria \emph{sensitiva} recorda as sensações e os objectos \emph{passados}; ora: o sensivel \emph{presente}, o \emph{ausente} e o \emph{passado} são objectos formaes distinctos, que exigem faculdades organicas distinctas. Distingue-se tambem da \emph{estimativa}; pois, emquanto a memoria \emph{sensitiva} suppõe necessariamente a \emph{experiencia}, a \emph{estimativa} é anterior a toda a \emph{experiencia}; ou, melhor, emquanto o objecto da memoria \emph{sensitiva} é percebido pelos sentidos \emph{externos}, o objecto da \emph{estimativa} escapa ao alcance d'esses sentidos. :\par \emph{Galluppi} (\emph{Les}. \emph{di Filosof}. lez. XC) confundiu a \emph{memoria sensitiva} com a \emph{imaginação}. Mas é evidente que a \emph{memoria sensitiva} e a \emph{imaginação} sao duas faculdades distinctas e diversas entre si. Na verdade:\par a) A \emph{imagem} póde considerar-se --- ou emquanto é \emph{imagem} de um objecto, --- ou emquanto é \emph{imagem} de um objecto \emph{anteriormente percebido}. --- Considerada sob o primeiro aspecto, a \emph{imagem} é objecto da \emph{imaginação}; pois, pela \emph{imaginação}, conservamos a \emph{imagem} de um objecto, \emph{sem relação} ao \emph{passado}. --- Considerada sob o segundo aspecto, a \emph{imagem} é objecto da \emph{memoria sensitiva}; pois é proprio da \emph{memoria} conservar a \emph{imagem}, emquanto representa um objecto \emph{anteriormente} conhecido. Por isso a \emph{imagem} refere-se á \emph{imaginação} e \& \emph{memoria}, mas sob diverso \emph{aspecto}. Ora o objecto, que se refere á força ou energia da alma \emph{sob diverso aspecto}, \emph{é formalmente} multiplo, e por isso exige distinctas faculdades; pois as faculdades distinguem-se pelos objectos \emph{formaes}. Logo a \emph{imaginação} e a \emph{memoria} são faculdades distinctas.\par \emph{b}) Demais, o \emph{acto} da \emph{imaginação} é independente da existencia real do objecto externo, pois formamos imagens de objectos que nunca existiram; ao passo que o acto da \emph{memoria sensitiva} é sempre dependente da existencia do objecto, pois ninguem se recorda do que nunca viu. Ora as faculdades distinguem-se pelas suas operações \emph{especificamente} diversas. Logo a \emph{imaginação} e a \emph{memoria sensitiva} são duas faculdades distinctas.\par A \emph{memoria sensitiva}, embora seja distincta e diversa do \emph{sentido commum} e da \emph{imaginação}, todavia, nas suas operações, suppõe necessariamente o exercicio d'estas duas faculdades. --- Suppõe o exercicio do \emph{sentido comam s} porque este recebo dos sentidos \emph{externos} os objectos, de que depois a \emph{memoria} se recorda. --- Suppõe o exercicio da \emph{imaginação}; pois á \emph{memoria} se recorda de um objecto \emph{ausente}, e não poderia recordar-se de um objecto \emph{ausente}, senão pela \emph{imagem} do mesmo objecto, conservada na \emph{imaginação}.}.

%%>>> t98

\pg{131} \tese{98}{A memoria sensitiva é uma faculdade material.}

a) Uma faculdade é \emph{material}, quando o seu objecto é \emph{material} ou \emph{sensivel}. Ora o objecto da \emph{memoria sensitiva é material} ou \emph{sensivel}; pois \emph{sensiveis} são as percepções \emph{passadas}, que ella \emph{conserva}, \pg{132}  e \emph{sensiveis} são os \emph{objectos} d'aquellas percepções. Logo a memoria \emph{sensitiva} é uma faculdade \emph{material}.

b) Uma faculdade, que se encontra tambem no animal, não é \emph{intellectual}, mas \emph{material} ou \emph{sensitiva}; pois o animal não é dotado de faculdades intellectuaes. Ora a \emph{memoria sensitiva} encontra-se tambem no animal; pois este \emph{recorda} o \emph{seu} covil ou o \emph{seu} ninho, os beneficios recebidos no \emph{passado}, e, se foge do logar onde foi espancado, volta para o logar onde encontrou alimento. Logo a \emph{memoria sensitiva} é uma faculdade \emph{material} \footnote{Confirmam a these os seguintes argumentos:\par 1.9) Uma faculdade, em que influam varias circumstancias materiaes, é \emph{material}. Ora á \emph{memoria sensitiva} é uma faculdade, em que influem \emph{directamente} circumstancias materiaes; porque nas creancas é muito viva e nos velhos muito languida, e as doenças cerebraes chegam a perturbal-a e quasi a apagal-a. Logo a \emph{memoria sensitiva} é uma faculdade \emph{material}.\par 2.0) Se a \emph{memoria sensitiva} fosse uma faculdade \emph{intellectual} e nao \emph{material}, a sua evolução deveria ser proporcionada á evolução da intelligeneia. Mas não é assim; pois alguns são dotados de muito talento e de pouca memoria, e vice-versa. Logo a \emph{memoria sensitiva} é uma faculdade \emph{material}.\par Erraram, pois, \emph{Rosmini} (\emph{Psícol}., lib. IT) e \emph{Gioberti} (\emph{Errori Filos}., Lett. VI), affirmando, absolutamente e sem distincção, que a \emph{memoria}, cujo objecto é o \emph{passado}, devia classificar-se entre as faculdades \emph{intellectuaes}.\par Para evitar toda a confusão, devemos advertir que o homem, por ser um \emph{animal racional}, é dotado de duplice \emph{memoria}, --- uma \emph{sensitiva}, outra \emph{intellectual}. Estas duas memorias, embora conservem entre si uma perfeitu harmonia (pois a \emph{sensitiva} esta subordinada á \emph{intellectual}), todavia sao duas faculdades distinctas. Na verdade, duas faculdades são distinctas, quando os seus objectos são distinctos. Ora o objecto da memoria \emph{sensitiva}, que é constituido pelas nossas percepções \emph{sensitivas} passadas, é distincto do objecto da memoria \emph{intellectual}, que consiste nas nossas percepções \emph{én/ellectuces} passadas; porque a percepção \emph{sensitiva} é uma coisa \emph{concreta} e \emph{material}, ao passo que a percepção \emph{intellectual} é uma coisa \emph{immaterial} e \emph{abstracta}.\par Objectam: A percepção do \emph{passado} inclue a percepção de \emph{tempo}. Ora a percepção de \emph{tempo} é um acto da intelligencia. Logo a \emph{memoria} não é uma faculdade \emph{sensitiva}, mas é só \emph{intellectual}.\par Resposta. A percepção de \emph{tempo} póde ser \emph{espiritual} e \emph{immaterial}, ou \emph{sensitiva} e \emph{material}. E' \emph{espiritual} e \emph{immaterial}, quando conhecemos --- a \emph{essencia} do \emph{tempo}, --- as suas partes, que são o \emph{passado}, o \emph{presente} e o \emph{futuro}, --- e a \emph{ordem} por que se succedem. Esta percepção é acto da intelligencia. E' \emph{sensitiva} e \emph{material}, quando exprime uma circumstancia propria de uma coisa \emph{sensivel} ou \emph{material}, sem se perceber a essencia do \emph{passado}. Ora a percepção do \emph{passado}, que fórma o objecto da \emph{memoria sensitiva}, \emph{é sensivel} e \emph{material}; porque exprime apenas uma circumstancia propria de uma coisa \emph{sensivel} e \emph{material}, como é a sensação \emph{passada}, mas sem se advertir de que é \emph{passada}. Logo a memoria não é só \emph{intellectual}, mas ha tambem a memoria \emph{sensitiva},\par Sendo uma faculdade \emph{material}, \emph{á memoria sensitiva} tem tambem as suas \emph{doenças}. As principaes são a \emph{amnesia}, \emph{á ipermnesia} e a \emph{paramnesia}, --- A.\pgn{133}\emph{amnesia} é a falta ou perda da memoria. Esta doença --- emquanto á \emph{origem}, pode ser \emph{congenita} ou \emph{adventicia}, conforme depende de uma fraqueza \emph{natural} do cerebro, ou de uma fraqueza \emph{eventual}; --- emquanto á \emph{extensão}, \emph{é geral} ou \emph{parcial}, conforme se perdem \emph{todas} as lembranças, ou só uma parte d'ellas; --- emquanto ao \emph{desenvolvimento}, \emph{é temporanea} (cessa depois de algum tempo), ou \emph{periodica} (alterna-se com o exercicio regular da memoria), ou \emph{progressiva} (quando as lembranças se perdem gradualmente até á perda total) --- A amnesia \emph{parcial} é quadrupla, conforme a \emph{variedade} das lembranças, que vão desapparecendo; e temos --- a \emph{aphemia} ou \emph{aphasia}, quando se não recorda uma lingua, já conhecida e fallada, --- a \emph{surdez verbal}, quando se não recorda o sentido de uma palavra, que se pronuncia, --- a \emph{logoplegia}, 'quando se não recordam as palavras correspondentes aos conceitos, --- a \emph{agraphia}, quando se esquece o modo de escrever. --- \emph{A épermnesia} consiste numa extraordinaria excitação da memoria (assim um doente analphabeto repetia, no \emph{delirio}, um grande numero de versos de uma tragedia, ouvidos uma só vez) --- A \emph{paramnesia} é uma illusão da memoria, pela qual ou se trocam as lembranças, ou se julga como percebido o que nunca o foi. --- As causas d'estas doencas da memoria são as perturbações ou os desarranjos do \emph{systema nervoso} e \emph{sanguineo}, produzidos pelo estado pathologico (\emph{epilepsia}, \emph{anemia}, ecc.) ou por agentes artificiaes. Assim o bromuro de potassio provoca a \emph{amnesia}; o café, o chá, o opio, o alcool produzem a \emph{épermnesia}.}.

%%>>> t99

\pg{133} \tese{99}{A memoria sensitiva no homem e no animal.}
A \emph{memoria sensitiva}, embora seja commum ao homem e ao animal, todavia naquelle é mais elevada que neste. Com effeito, o animal, pela \emph{memoria sensitiva}, conserva as \emph{proprias} percepções sensitivas \emph{passadas}; mas com um ou outro elemento, de que se recorda, não póde reconstruir por inteiro uma percepção \emph{passada} ; pois que este trabalho requer uma reflexão da intelligencia e uma determinação da vontade, que o animal não possue. Pelo contrario, o homem, pela \emph{memoria sensitiva}, póde com alguns elementos reconstruir \emph{toda} a percepção \emph{passada}, pois a nossa \emph{memoria sensitiva} é fortalecida e guiada pela \emph{razão}, com que está unida. --- A \emph{memoria sensitiva}, emquanto auxiliada pela razão reconstrue por inteiro uma percepção sensitiva \emph{passada}, chama-se propriamente \emph{reméniscencia} \footnote{A memoria \emph{sensitiva} existe no homem e no animal. Ninguem pode duvidar d'isso. O animal conserva a lembranca dos beneficios e dos maus tratos recebidos, e, conforme a diversidade d'essas lembranças, mostra a sua amizade ou inimizade. Mas é tambem certo que á \emph{memoria}, como se encontra no animal, não excede a ordem sensitiva. Elle só conserva a lembrança de coisas materiaes; e essas mesmas lembrancas nunca apparecem unidas entre si por um nexo \emph{logico}, mas sempre e só por meio de associações \emph{empiricas}, innatas ou adquiridas, que podem contrafazer o nexo das nossas ideas e imitar o acerto dos nossos juizos, mas que não se podem de modo nenhum attribuir a uma faculdade intellectual. --- Como tambem a lembrança de um: objecto anteriormente conhecido, que vemos existir no animal, não demonstra nelle nem o sentimento da propria identidade, nem as ideas abstractas do\pgn{134}\emph{presente} e do \emph{passado}, do \emph{tempo} e do \emph{espaço}. Não é, pois, essa lembrança, nem pode ser, um acto formal de \emph{reconhecimento}, como o é o da memoria \emph{intellectual} no homem; mas é uma \emph{associação} de imagens totalmente \emph{empérica}, a qual, comtudo, produz os mesmos effeitos praticos, que produz o \emph{reconhecimento} da nossa intelligencia, e basta para as necessidades da vida de relação.\par Por isso a \emph{reminiscencia} ea \emph{memoria sensitiva} differem por duas razões: --- 1.9) A \emph{reminiscencia}, pelo facto de reconstituir com um ou outro elemento \emph{toda} a percepção passada, forma como que um raciocinio \emph{éndauctivo}, embora \emph{imperfeito}, que chega a conclusões \emph{particulares}; ao passo que a \emph{memoria sensitiva} é incapaz de todo o raciocinio, embora imperfeito. --- 2.9) A \emph{reminiscencia}, nos seus actos, importa tambem o exercicio da \emph{razão} e da \emph{vontade}; mas á \emph{memoria sensitiva} não é acompanhada ou auxiliada pelas operações das faculdades intellectuaes (Cf. S. Thom., \emph{De mem}. et \emph{remin}., lect. 5).\par \emph{Aristoteles} escreve na \emph{historia dos animaes} (\emph{lib}. I, c. 1): « Muitos animaes, diversos do homem, possuem tambem a faculdade de se recordar e de perceber, mas sô o homem tem o dom de se lembrar \emph{quando quer} ». E noutro logar (\emph{Da memoria}, c. 9) diz: « Dá-se esta differenca, porque á \emph{reméniscencia} é uma especie de \emph{raciocinio} » (\emph{Cf}. S. Thom., \emph{De mem}. \emph{et rem}., lect. 5).\par \emph{Stewart} afirma que á \emph{reminiscencia} é \& propria memoria \emph{ntellectual}, e por isso differe da memoria \emph{espontanea} ou \emph{sensitiva}.\par Esta affirmação é falsa. A \emph{reminiscencia} não differe da \emph{memoria sensitiva}; pois o objecto de uma e de outra é constituido pela percepção sensitiva \emph{passada}. Se entre a \emph{reminiscencia} e a \emph{memoria sensitiva} existem as differencas apontadas, estas são «\emph{ccidentaes}, e derivam da \emph{rasdo}, com que á \emph{memoria sensitiva} do homem está unida.\par No homem, o desenvolvimento da \emph{memoria sensitiva} --- não só depende da \emph{natureza}, --- mas tambem da \emph{arte}, chamada «\emph{nemotechnia}.\par \emph{a}) Depende da \emph{natureza}; por quanto, sendo a \emph{memoria sensitiva} uma faculdade organica, é claro que será tanto mais perfeita, quanto mais perfeito fôr o seu orgão. D'ahi a diversidade da \emph{memoria} segundo o desenvolvimento do organismo, e segundo a edade e o temperamento.\par \emph{b}) Depende da \emph{arte}; porque em nós o \emph{exercicio} da \emph{memoria} não é determinado exclusivamente, como nos animaes, pelo instincto e pelas causas naturaes. --- Os principaes meios \emph{artificiaes}, que constituem a \emph{mnemotechnia}, isto é, a arte de fortificar e cultivar a \emph{memoria}, são \emph{tres}: --- 1.9) representar-se, por meio de semelhanças, os objectos, que queremos reter; --- 2.9) classificar methodicamente os objectos, que se querem recordar; --- 3.9) reflectir profunda e constantemente sobre os mesmos objectos.}.

%%>>> t100

\pg{134} \tese{100}{Orgão da memoria sensitiva.}
O orgão da \emph{memoria sensitiva} é o \emph{cerebro}; pois que a experiencia demonstra que uma lesão do \emph{cerebro} produz um desarranjo mais ou menos notavel nesta faculdade \footnote{Segundo \emph{Luys} (\emph{Le Cerveau}, pag. 108), \emph{a} camada cortical do \emph{cerebro} é a séde principal, mas não exclusiva, da \emph{memoria}, Esta encontra-se, num estado mais ou menos rudimentar, em todas as regiões centraes do systema nervoso. Nos centros da medulla, á memoria apparece com um caracter tão 'accentuado, que elles sós podem presidir ás manifestações da vida automatica.\par Muitos physiologistas modernos localisam a memoria nos \emph{lobos frontaes} ; e alguns indicam pontos especiaes do cerebro para a memoria da \emph{escriptura},\pgn{135}ou da \emph{palavra}, ou do \emph{tempo}, ou do \emph{logar}. E provam estas affirmações com varios exemplos. Assim, o cão, ao qual se arrancaram os \emph{lóbos occipitaes}, esquece as lembranças relativas á \emph{vista}, mas não as relativas ao \emph{ouvido}. Se se paralysarem no homem os \emph{nervos opticos} até ao lóbo \emph{occipital}, dá-se a \emph{agraphia}, isto é, o homem esquece o modo de escrever; assim como, se se lesarem os \emph{nervos acusticos} até aos \emph{lóbos temporaes}, dá-se a \emph{aphemia}, isto é, o esquecimento da palavra. --- Mas estamos sempre no campo das hypotheses. A exacta localisação dos sentidos no cerebro está muito longe de ser um facto scientificamente certo.\par Apraz-nos resumir a doutrina exposta, citando o seguinte trecho da classica obra --- \emph{Luz e Calor} --- do P. \emph{Manuel Bernardes}, em que o mavioso e grave Oratoriano falla dos \emph{sentidos internos} e das suas funcções, empregando a mesma terminologia, que nós empregamos. --- Recordamos que a \emph{pAanasia} não é uma faculdade \emph{especificamente} distincta da \emph{émaginativa}, nem a \emph{reminiscencia} da \emph{memoria sensitiva}. Eis as suas palavras:\par « Nesta região pois sensitiva ha umas faculdades, cujo officio é \emph{conhecer}: e outras correspondentes a estas, cujo officio é \emph{querer}; por quanto a nossa alma de nada pega para si com uma mão, sem primeira palpar com a outra, e ver se lhe convem.\par 4 Das faculdades cognoscitivas umas são exteriores, e estas são os cinco sentidos externos --- \emph{ver}, \emph{ouvir}, \emph{cheirar}, \emph{gostar}, \emph{palpar}: outras são interiores, e d'estas não consta entre os filosofos numero certo. Uns dizem, que é uma só, chamada \emph{sentido interno} ou \emph{commum}, e que este só serve todos os offieios, ou funcções, que logo diremos. Outros dizem, que são duas, a saber, \emph{sentido communm} e \emph{fantasia}. S. Thomas, a quem seguem muitos, tem para si que são quatro; porque ás tres sobredictas (\emph{fantasia}, \emph{estimativa} e \emph{memoria material}) \&ccrescenta o \emph{sentido commum}; e por ventura se governam pelo numero dos ventriculos, ou concavidades, que ha no cerebro; os quaes são quatro, e parecem ser dados para officinas dos sentidos interiores. O B. AIberto Magno (a quem seguem outros) diz que são cinco, acerescentando além destas a \emph{émaginativa}. E não falta quem ponha tambem a sexta, que se chama \emph{reminiscencia}.\par 4 As funcções, ou officios d'estes sentidos interiores, são as seguintes.\par « O \emph{sentido comum} percebe por modo universal todos os objectos dos sentidos externos particulares, e como censor residindo na substancia do cerebro, discerne entre elles. Por este discernir alcançaram os filosofos haver esta faculdade na nossa alma: pois era operação mui necessaria á natureza humana, e não podia proceder de sentido algum particular. Por quanto para discernir entre duas coisas, é necessario conhecer ambas: e nenhum dos sentidos\pgn{136} particulares póde conhecer mais, que do seu particular objecto que Ihe ioca. E assim a vista poderá discernir entre duas córes, e entre duas luzes: o ouvido entre deus sons: o gosto entre deus sabores; porém nem a vista, nem o ouvido, nem o gosto póde discernir entre cor, som, e sabor: e isto é o que faz o sentido commum, por isso assim chamado... E o poder uma só faculdade, sendo material, tocar sensiveis tão diversos, lhe vem de ser potencia mais perfeita e universal, que une em si a virtude dos sentidos exteriores.\par u A \emph{fantasia} (ou \emph{imaginação}) serve de receber e guardar as especies, ou imagens dos objectos sensiveis exteriores, para depois obrar com ellas em ausencia dos taes objectos... --- A \emph{émaginativa} (\emph{creadora}) accrescenta sobre a \emph{fantasia}. o compor d'essas especies de diversos objectos, que nella se guardavam, varias cousas possiveis, e fingir outras impossiveis: como agora, da especie de ouro, fórma um monte de ouro: da especie do leão e da de aguia, fôrma um grypho. E por isso aos que se entretem, ou deleitão nestas composições e fingimentos, chamamos imaginativos, ou fantasticos.\par \& A \emph{estimativa} serve de apprehender os objectos externos sensiveis debaixo de alguma razão, que não é sensivel exteriormente, como é, v. g., á amizade, ou inimizade, a conveniencia, ou disconveniencia. O que se vé manifestamente até nos brutos; pois a ovelha sabe a herva, que lhe convém, ou não convém; e que o pastor é seu amigo, e o lobo seu inimigo.\par u A \emph{memoria} serve de conhecer as cousas, com a reflexão de que outra vez já foram d'ella conhecidas. --- A \emph{reminiscencia} serve de inquirir, rastrear, e revocar á memoria o que já esqueceo, pelos vestigios do que ainda lembra\par \& Todas estas faculdades, ou sentidos interiores são materiaes, e nellas se parece o homem com os brutos: não fallando na reminiscencia; porque esta, supposto que alguns dizem que é tambem sensitiva, a outros lhes parece envolver já actos de racional. E ou se distingam entre si, e da alma, ou não se distingam; ou sejam mais, ou menos no numero (que nisto ha varias opiniões), o certo é, que na alma ha todos estes ministerios, e officios, pela parte que toca a esta sua região intima, sensitiva e cognoscitiva » (\emph{Lus e Calor}, Primeira parte, Doutrina III).}.

%%>>> art-17

\section*{Artigo III.}
\addcontentsline{toc}{section}{Artigo III.}

\pg{135} \textbf{Appetite sensitivo}, \textbf{faculdade locomotiva}

e \textbf{vocal}.

%%>>> t101

\tese{101}{Appetite sensitivo.}
\emph{Appetite sensitivo} é a faculdade pela qual o homem tende para um bem, conveniente á sua natureza \emph{sensitiva} e percebido por meio dos sentidos.

Esta faculdade existe no homem. Com effeito, a tendencia de um ente está em proporção com o seu conhecimento. Ora o homem, emquanto dotado de vida \emph{sensitiva}, conhece bens sensiveis, necessarios para a sua conservação ou perfeição. Logo o homem deve ter e tem uma tendencia natural para os bens sensiveis. Esta tendencia é o \emph{appetite sensitivo} \footnote{\emph{Appetite}, em geral, é a tendencia ou inclinação de um ente para o seu fm, isto é, para uma coisa que lhe convém e que se chama \emph{bem}. Por isso é o \emph{bem} que todos os seres appetecem: \emph{Bonum}. \emph{est quod omnia appetunt}. Todos os entes creados tendem para o seu fim, determinado por Deus; e essa inclinação ou tendencia foi por Deus impressa ou gravada na propria , forma natural da cada um, e por isso identifica-se com a natureza d'elles e é commum a todos. Esta tendencia diz-se \emph{appetite natural}. --- Além d'esta, ha, em alguns entes, outras tendencias para o bem, que nascem não da forma\pgn{137}natural, mas de algumas formas adquiridas pelo conhecimento. Estas tendencias são actos de faculdades \emph{appetitivas}. (\emph{Sum Theol}., p. I, q. 80, aa. 1, 2).\par Faculdade \emph{appetitiva}, em geral, é a força pela qual o ser vivo tende para um bem conhecido, emquanto convém á sua natureza.\par O \emph{appetite} é duplice --- \emph{sensitivo} e \emph{intellectual}, porque duplice é a percepção --- \emph{sensitiva} e \emph{intellectual}, \emph{e} a cada percepção corresponde um \emph{appetite} da mesma ordem ou natureza. Appetite \emph{sensitivo} é o que tende para o bem \emph{sensivel}, percebido por meio dos sentidos. Appetite \emph{intellectual} é o que tende para o bem \emph{intellectual}, isto é, conveniente á vida \emph{intellectual} ou \emph{racional}, percebido por meio da intelligencia.\par A faculdade \emph{appetitiva}, como dissemos, corresponde á faculdade \emph{perceptiva}. Por isso, os animaes, cujas faculdades \emph{perceptivas} são exclusivamente \emph{sensitivas}, só possuem a faculdade \emph{appetitivo-sensitiva}. As substancias espirituaes, cujas faculdades são meramente \emph{espirituaes}, só possuem a faculdade \emph{appetitivo-intellectual}. O homem, que, partecipando dos animaes e dos espiritos, possue as faculdades \emph{perceptivas} da ordem \emph{sensitiva} e da ordem \emph{intellectual}, é dotado de uma e outra faculdade \emph{appetitiva}, --- \emph{sensitiva} e \emph{intellectual}.\par Ha, portanto, no homem um triplice appetite --- o \emph{natural}, o \emph{sensitivo} e o \emph{intellectual}. O appetite \emph{sensitivo} e o \emph{intellectual} não só não excluem o \emph{natural}, mas suppõem-no e exigem-no; porque tanto o \emph{sensitivo} como o \emph{intellectual,mesmo} na presença do bem conveniente, ficariam indifferentes, nem ienderiam para elle, se não tivessem essa inclinação ou \emph{appetite natural}, --- Deve notar-se que, inclinando o appetite \emph{natural} para um bem \emph{real} (que existe \emph{realmente} na coisa), e o appetite \emph{sensitivo} e o \emph{intellectual} para um bem \emph{percebido} pelas faculdades cognitivas, o bem do appetite \emph{natural} é sempre \emph{verdadeiro} bem, ao passo que o bem do appetite \emph{sensitivo} e do \emph{intellectual} pode ser um bem \emph{apparente} (\emph{Summ}. \emph{Theol}., 1-2, q. 8, a. 1).}.

%%>>> t102

\tese{102}{Objecto do appetite sensitivo.}
O objecto do \emph{appetite sensitivo} é o \emph{bem} corporeo, concreto e sensivel, isto é, relativo á conservação e perfeição do ser sensitivo; visto que só esse bem \pg{136} \pg{137}  é conveniente á natureza \emph{sensitiva} e é percebido pelos sentidos. Por isso o appetite \emph{sensitivo}, embora supponha o exercicio das faculdades \emph{perceptivo-sensitivas}, é distincto d'ellas \footnote{A faculdade \emph{appetitiva} differe das \emph{perceptivas} por tres razões: --- 1.2) As faculdades \emph{perceptivas attrahem} a si o objecto, que pela sua imagem entra nas proprias faculdades, e ahi \emph{fica}; ao passo que a faculdade \emph{appetitiva} tende para o objecto, e, attingindo-o, se lhe une. --- 2.º) As faculdades \emph{perceptivas} estão para a \emph{appetitiva}, como o motor para o movel; pois esta é determinada ao acto por aquellas. --- 8.º) Embora o objecto das faculdades \emph{perceptivas} e o da \emph{appetitiva} sejam \emph{materialmente} identicos, todavia \emph{formalmente} são diversos; pois as faculdades \emph{perceptivas} attingem um objecto emquanto é \emph{verdadeiro}, ao passo que a faculdade \emph{appetitiva} attinge-o emquanto é \emph{bom}. D'ahi a distineção entre a faculdade \emph{perceptiva} e á \emph{appetitiva}.}.

%%>>> t103

\tese{103}{Divisão do appetite sensitivo.}
O \emph{appetite sensitivo} divide-se em \emph{concupiscivel} e \emph{irascivel}.

a) \emph{Appetite concupiscivel} é a faculdade pela qual o homem --- ou tende para um bem sensivel, emquanto este é \emph{agradavel}, --- ou foge de um mal sensivel, emquanto é \emph{desagradavel} \footnote{\emph{A} tendencia para o bem sensivel e a fuga do mal sensivel são dois actos do mesmo appetite \emph{concupiscivel}; porque é a mesma força que, approximando-se de um termo, se afasta do opposto.}.

b) \emph{Appetite irascivel} é a faculdade pela qual o homem tende para um bem sensivel, emquanto este é \emph{arduo}, mas \emph{util} para remover os impedimentos, que se oppõem á sonsecução do bem sensivel ou á fuga do mal sensivel \footnote{O bem, que forma o objecto do appetite /\emph{rascivel}, chama-se \emph{arduo}, --- não só porque, para o attingir, o ser sensitivo encontra difficuldades, --- mas tambem porque tende, por si mesmo, para destruir os obstaculos que se oppõem á consecução do bem, que é o objecto do appetite \emph{concupiscivel}.\par Dizem: O objecto de todo o appetite \emph{sensitivo} é o bem \emph{conveniente} ao ser \emph{sensitivo}. Ora o bem \emph{arduo} não é conveniente ao ser sensitivo. Logo o beni \emph{arduo} não é objecto do appetite \emph{irascivel}.\par Resposta. Concedemos que o objecto de todo o \emph{appetite sensitivo} é um bem \emph{conveniente} ao ser \emph{sensitivo}. Mas o bem póde ser conveniente ao ser \emph{sensitivo}, --- ou emquanto lhe é \emph{agradavel}, --- ou emquanto lhe é \emph{util} para a sua defeza. Ora o objecto do appetite \emph{irascivel} é um bem conveniente ao ser sensitivo; pois que, embora \emph{arduo}, todavia lhe é \emph{util} para remover os obstaculos que lhe impedem a posse do bem \emph{agradavel}.}.

Assim, pelo appetite \emph{concupiscivel}, desejamos o alimento e rejeitamos o veneno; e, pelo \emph{zrascivel}, desejamos a lucta, como um bem que, embora \emph{arduo}, é \emph{util} para resistirmos aos que nos impedem a posse do alimento, ou subministram o veneno \footnote{E' innegavel a existencia d'este duplice \emph{appetite} nos seres \emph{sensitivos}. Um cão agarra o osso, que Ihe foi apresentado; e, se outro cão lhe disputar a posse d'aquelle osso, deixa de comer, e faz todo o esforco para repellir o adversario, e para conservar para si o alimento. Em ambos os casos, o cão é guiado pelo conhecimento, que tem --- do alimento, emquanto este convém á sua natureza, --- e do adversario, emquanto este é obstaculo para a consecução d'aquelle seu bem. Em ambos os casos, o appetite \emph{sensitivo} tende para o bem, com esta differenga que no primeiro caso tende para o bem, emquanto este é \emph{agradavel}, ao passo que no segundo caso tende para o bem, emquanto este, embora \emph{arduo}, \emph{é util} para remover os obstaculos que Ihe impedem a posse d'aquelle bem.\par Ora o \emph{appetite sensitivo}, emquanto tende para o bem «\emph{gradavel} aos sentidos, diz-se \emph{concupiscivel}, e, emquanto tende para o bem «\emph{arduo}, mas \emph{util} para vencer os obstaculos, que se oppõem á posse do bem, diz-se \emph{irascivel}.\par Ambos os \emph{appetites}, \emph{o concupiscivel} e o \emph{irascivel}, são necessarios ao ser \emph{sensitivo}; pois este não poderia conservar a propria existencia, nem alcançar a devida perfeição, se não tivesse o poder de procurar o bem e de fugir do mal,e o de resistir a tudo o que possa destruir a sua existencia ou impedir a perfeição conveniente é sua natureza.}.

%%>>> t104

\tese{104}{O appetite concupiscivel e o irascivel são duas faculdades distinctas.}
Duas faculdades são realmente distinctas, quando os seus objectos \emph{formaes} são realmente distinctos ; pois, como vimos, as faculdades distinguem-se, em ultima analyse, pelos seus objectos \emph{formaes}. Ora o objecto \emph{formal} do appetite \emph{concupiscivel} e o do appetite \emph{irascivel} são realmente distinctos. Na \pg{138} \pg{139}  verdade, embora ambos os appetites tendam para o bem, todavia não o attingem ambos sob o mesmo aspecto; visto que o \emph{concupiscivel} tende para o bem, emquanto é \emph{agradavel}, e o \emph{irascivel} tende para o bem, emquanto é \emph{arduo}, mas \emph{util} para afastar os obstaculos, que impedem a consecução de outro bem ou a fuga do mal. Logo o appetite \emph{concupiscivel} e o \emph{irascivel} são duas faculdades distinctas \footnote{Ao argumento, citado no texto, accrescentamos o seguinte:\par Não podem attribuir-se a uma faculdade actos, que são contrarios á sua inclinação. Ora os actos do appetite /\emph{rascivel} são contrarios á inclinação do appetite \emph{concupiscivel}. Na verdade, o \emph{concupiscivel} tende para o que nos [ agradavel, ao passo que o /\emph{rascivel} tende muitas vezes para o que nos causa tristeza e dor sensivel. D'ahi o antagonismo entre as paixões d'estas duas especies do appetite; pois, augmentando as paixões do \emph{concupiscivel}, diminuem as do \emph{irascivel}, e vice-versa. Logo o appetite \emph{concupiscivel} e o \emph{irascivel} são duas faculdades distinctas (\emph{Sum}. \emph{Th}. p. I, q. 81, a. 2).\par D'este modo o appetite \emph{sensitivo} é o \emph{genero provimo}, que se divide em duas \emph{especies} --- o appetite \emph{concupiscivel} e o \emph{irascivel},\par O appetite \emph{concupiscivel} e o \emph{irascivel}, posto que sejam duas faculdades distinctas, comtudo conservam entre si as seguintes relações:\par 1.5) O appetite \emph{irascivel} tem \emph{principio} e \emph{termo} no \emph{concupiscivel}. --- Tem \emph{principio} no \emph{concupiscivel}. Na verdade, o \emph{irascivel} não combate os obstaculos, senão emquanto estes impedem a consecução do bem ou a fuga do mal. Por isso os movimentos do \emph{irascivel} serão mais ou menos intensos, quanto mais ou menos intensos forem os movimentos do \emph{concupiscivel}. --- Tem o seu \emph{fermo} no \emph{concupiscivel}. Por quanto, o \emph{irascivel} ou remove effectivamente os obstaculos, ou não. No primeiro caso, teremos \emph{alegria}; no segundo, \emph{tristeza}. Ora a \emph{alegria} e a \emph{tristeza} são duas paixões do \emph{concupiscivel}.\par 2.º) O \emph{irascivel} aperfeiçõa e completa o \emph{concupiscivel}. Por quanto, de pouco ou nada serviria tender para o bem, se não houvesse uma faculdade para remover os obstaculos, que impedem a consecução do mesmo bem.\par 8.3) O \emph{irascivel}, emquanto ao \emph{objecto}, é mais nobre do que o \emph{concupiscivel}; pois o bem, que serve para a utilidade e defeza do ser sensitivo e que é o objecto do \emph{irascivel}, é mais nobre e menos material do que o bem, que agrada ás faculdades sensitivas e que é o objecto do \emph{concupiscivel}, --- Por isso o appetite \emph{irascivel} está para as faculdades \emph{appetitivas}, como a \emph{estimativa} está para as faculdades \emph{perceptivas}. Na verdade, assim como a \emph{estimativa}, percebendo certas propriedades, que escapam ao alcance dos outros sentidos, approxima-se mais da razão; assim tambem o \emph{irascivel}, tendo por objecto o que é \emph{util} para vencer as difficuldades, que se oppõem á consecução do bem e á fuga do mal, é mais proximo do appetite \emph{racional} ou \emph{intellectual}, que é a \emph{vontade}. --- Dissemos --- emquanto ao \emph{objecto}; pois o \emph{irascivel}, emquanto se refere \emph{ao concupiscivel}, como meio ao fim, é menos nobre do que este, visto que o meio é sempre menos nobre que o fim,}.

%%>>> t105

\tese{105}{O appetite sensitivo é uma faculdade material.}

a) Uma faculdade é \emph{material}, quando o seu objecto é \emph{material} ou \emph{sensivel}; pois as faculdades conhecem-se e distinguem-se \pg{140}  pelos seus objectos. Ora o objecto do \emph{appetite sensitivo é material}; pois é constituido pelo bem corporeo, individual, relativo á conservação e perfeição do ser sensitivo. Logo o \emph{appetite sensitivo} é uma faculdade \emph{material}.

b) Uma faculdade, que se encontra tambem no animal, é \emph{material}; pois o animal não é dotado de faculdades 2\emph{mmateriaes} ou \emph{espirituaes}. Ora o \emph{appetite sensitivo} encontra-se tambem nos animaes: pois elles --- não só tendem para os bens sensiveis, que aperfeicóam a sua natureza, e fogem do mal sensivel, --- mas tambem luetam para remover os obstaeulos, que lhes podem impedir a consecução do bem e a fuga do mal. Logo o \emph{appetite sensitivo} é uma faculdade \emph{material} \footnote{O P. Manuel Bernardes diz: « A esta virtude, ou força cognoscitiva e material, e ás noticias, que por sua via bebeo a alma, corresponde na mesma alma outra virtude, ou força tambem \emph{material}, dada para querer o bem, ou fugir do mal já conhecido. Esta chama-se \emph{appetite}, \emph{e} consta de partes, a saber, \emph{concupiscivel} e \emph{irascivel} ». (\emph{Lus e calor}, Primeira parte, Doutrina III).}.

%%>>> t106

\tese{106}{O appetite sensitivo no homem e no animal.}
O \emph{appetite sensitivo}, embora seja commum ao homem e ao animal, todavia naquelle existe de um modo differente do que existe neste. Na verdade, o \emph{appetite sensitivo} do animal é \emph{fatal}, isto é, opéra sempre e necessariamente por um modo determinado e singular, visto que determinadas e concretas são as formas, ou os modelos da sua imaginação; ao passo que o \emph{appetite sensitivo} do homem está subordinado ao imperio da razão e participa, de algum modo, da liberdade da vontade \footnote{O appetite \emph{concupiscivel} e o \emph{irascivel}, no homem, devem obedecer á razão e á vontade. Com effeito, o homem, embora composto de corpo organico e de alma racional, é comtudo uma só substancia completa, e por isso uno deve ser o seu bem, \emph{una} a sua perfeição, como \emph{uno} é o fim para o qual tende. Se nelle existem varias e diversas faculdades, cada uma das quaes tem o seu bem e o seu fim especial; estes bens, estes fins especiaes, e por isso as suas faculdades relativas devem ser reduzidas á unidade. A unidade na variedade alcança-se por meio da ordem, e a ordem exige uma subordinação. 1j necessario, pois, que entre as varias faculdades, de que o homem é dotado, exista a subordinação. Ora a subordinação das faculdades deve ser proporcionada á subordinação dos bens, para os quaes ellas tendem (pois o fim é o criterio de toda a ordem); de modo que uma faculdade será subordinada á outra, se o bem d'aquella for subordinado ao bem d'esta. Ora o commum e verdadeiro bem de todo o homem, a que estão subordinados todos os bens particulares de todas as faculdades, é aquelle, que se apprehende pela intelligencia e para o qual, como para o seu fim ultimo, tende a vontade; pois a intelligencia e a vontade são as faculdades \emph{especificas} do homem. Logo todas as faculdades do homem devem estar subordinadas á razão e á vontade. D'ahi se segue que, se alguma das faculdades inferiores, na sua tendencia para o\pgn{141}seu bem, obstasse á consecução do bem, que é conforme á razão, ella não operaria conforme á ordem, e a sua operação, longe de ser vantajosa ao homem, ser-lhe-hia nociva. Podemos, pois, concluir que o appetite \emph{concupiscivel} e o \emph{irascivel}, no homem, devem obedecer á razão e á vontade.\par Objectam: Se o appetite \emph{concupiscivel} e o \emph{irascivel} do homem estivessem subordinados á \emph{razão}, não se deria aquella lucta entre o appetite e a razão, que, todavia, se dá em cada um de nós. A\par Resposta. Ha duas especies de dominio: um \emph{despotico}, outro \emph{politico}. E \emph{despatico}, quando o inferior, não possuindo liberdade propria, não póde resistir nem tornar-se rebelde ao superior. E \emph{politico}, quando o inferior, embora tenha o dever de obedecer ao superior, todavia, por ser dotado de liberdade, póde desobedecer-lhe e tornar-se rebelde. Ora, --- se o dominio que a alma exerce sobre o corpo, é \emph{despotico}, porque o corpo não póde resistir ao imperio da alma, --- o dominio, que a razão exerce sobre o appetite \emph{sensitivo}, quer \emph{concupiscivel} quer \emph{irascivel}, \emph{é politico}, porque o \emph{appetite sensitivo}, podendo ser movido e determinado ao acto, não só pela \emph{estimativa} ou \emph{cogitativa}, mas tambem pela \emph{imaginacdo} e pelo \emph{sentido commum}, é de algum modo independente da razão e póde resistir-]he.\par Perguntam: Porque não é \emph{despotico} o dominio da \emph{razão} sobre o \emph{appetite} ?\par São varias e diversas as razões, por que o dominio, que á razão exerce sobre o appetite, não é \emph{despotico}. Daremos as principaes. --- 1.º) O appetite succede á percepção. Logo, assim como a percepção \emph{sensitiva} é anterior á \emph{intellectual}, tambem o acto do appetite, \emph{sensitivo} póde por tal forma seguir a percepção dos sentidos, que previna a advertencia e o imperio da razão. --- 2.º) Demais, o «\emph{appetite sensitivo} reside num orgão corporeo; e por isso, para o seu exercicio, é necessario que o orgão não só seja informado pela alma, mas tambem que seja convenientemente disposto. Ora esta disposição do orgão nem sempre depende da razão e da vontade; e por isso podem produzir-se no \emph{appetite sensitivo} algumas tendencias, independentemente do imperio da razão. --- 3.º) Finalmente, ha certas e determinadas causas (por ex. uma forte commoção nos centros nervosos, uma extraordinaria perturbação de humores), que podem produzir e produzem certas e determinadas tendencias no appetite \emph{sensitivo}. Ora estas causas, quer pela sua natureza, quer pela sua vehemencia, não acabam sempre e logo que a vontade mande. Continuando a existir essas causas, continuarüo tambem a existir as relativas tendencias, apezar da resistencia da vontade. (\emph{Sum}. \emph{T}7., p. L, q. 81, a. 3).\par Insistem: De que modo obedece o «\emph{ppetite sensitivo á razão} ?\par O \emph{appetite sensitivo} segue sempre o conhecimento. Nos animaes, o appetite segue a \emph{estimativa}, assim a ovelha, julgando o lobo seu inimigo, foge d'elle: mas no homem deve seguir a \emph{cogitativa}, ou \emph{razão particular}. Ora a \emph{cogitativa}, ou \emph{razão particular}, fanda-se na \emph{razão universal}, que a move e dirige. Se, pois, no homem, o appetite \emph{sezsitivo} é immediatamente movido pela \emph{cogitativa}, e esta é dirigida pela razão, é evidente que o \emph{appetite sensitivo} obedece á razão. --- Nem esta obediencia é sempre violenta. Por quanto,\pgn{142}um e o mesmo objecto póde ser agradavel ou desagradavel, segundo o diverso aspecto sob o qual se considera. Ora a razão póde considerar num objecto uma propriedade agradavel, e, apresentando-o sob este aspecto ao appetite \emph{sensitivo}, póde mover esta faculdade a tender para elle. Assim a ira, concebida contra uma pessoa, acaba, se recordarmos os beneficios, que d'ella noutro tempo recebemos.\par Insistem ainda: De que modo obedece o \emph{appetite sensitivo} \& \emph{vontade}?\par O modo, por que o \emph{appetite sensitivo} obedece á vontade, é o seguinte: Todas as nossas faculdades tem a sua raiz na unica e simplicissima essencia da nossa alma. Por isso, quanto mais \emph{intenso} fôr o acto de uma faculdade, tanto menos o será o acto das outras, chegando ás vezes a \emph{intensidade} actual de uma faculdade a impedir inteiramente o exercicio das outras. Ora a vontade, por ser livre, póde dar aos seus actos uma grande intensidade; e, por ter dominio sobre as faculdades inferiores, póde transmittir essa intensidade ao appetite \emph{sezsitivo}, de modo que este siga o impulso d'aquella.}.

%%>>> t107

\tese{107}{Orgão do appetite sensitivo.}
O \emph{appetite sensitivo}, sendo uma faculdade \emph{material}, deve residir num orgão corporeo. Este orgão é o \emph{cerebro}, que é a sede principal da vida \emph{sensitiva}; embora commummente se diga que o orgão do appetite \emph{sensitivo} \pg{141} \pg{142}  É O \emph{coração}, emquanto os movimentos do \emph{coração} são o reflexo physiologico mais immediato e sensivel dos actos e das paixões do appetite \emph{sensitivo} \footnote{Qual é o verdadeiro e principal orgão do appetite \emph{sensitivo} ?\par Durante muito tempo, o \emph{coracdo} foi considerado, pelos philosophos e pelos naturalistas, como o verdadeiro e principal orgão \emph{elcitivo} da sensibilidade \emph{affectiva}, ou do appetite \emph{sensitivo}; assim como o cerebro era (e ainda hoje o é) considerado como o verdadeiro e principal orgão da sensibilidade \emph{percepliva}. Fundava-se esta opinião no consenso do genero humano, que reconheceu sempre o coração como a séde dos affectos, e no testemunho da consciencia, que manifesta claramente a reciproca influencia das paixões no. coração, e d'este naquellas.\par Hoje, porém, a maior parte dos physiologistas sustenta --- que o orgão \emph{elicitivo} da sensibilidade \emph{a/fectiva} não é o \emph{coração}, mas sim o \emph{cerebro}, --- e que o \emph{coração} é apenas o orgão \emph{manifestativo} d'essa, sensibilidade. --- \emph{a}) O orgão \emph{elicitivo} não é o \emph{coração}. Por quanto, o coração é um museulo, e os museculos não servem para a sensibilidade, mas só para o movimento. Além disso, a maravillosa estructura do coração indica claramente que, se elle é o orgão principal da vida \emph{vegetativa}, não o é do appetite \emph{sensitivo}. --- \emph{b}) O orgão \emph{elicitivo} dos affectos é o \emph{cerebro}. Por quanto, o cerebro é a sede da percepção \emph{sensitiva}, e por isso deve ser tambem a sido do appetite \emph{sensitivo}; porque, se o não fosse, o appetite, separado do centro perceptivo, não poderia tender para o bem, pois não se tende para o desconhecido. Além d'isso: o appetite, sendo o principio do movimento local espontaneo, deve residir naquella parte do corpo, em que se realiza a percepção sensitiva e da qual partem os nervos d' esse movimento; ora essa parte é o cerebro. Finalmente, é de todos sabido que as lesões do cerebro, como tambem os desarranjos nas suas funcções, alteram immediatamente a sensibilidade affectiva, e, vice-versa, as grandes afflicções perturbam a imaginação, até á alienação mental; ora isto prova que o orgão dos affectos é o cerebro. --- \emph{cj} O coração é o orgão \emph{manifestativo} dos affectos e das paixões, Esta funcção exerce-a o coração, não emquanto é um musculo, mas emquanto está unido ao cerebro por ganglios e fibras nervosas. Sendo tão estreita a communicação do coração com o gerebro, é natural que um movimento produzido no cerebro por um acto\pgn{143}ou paixão do appetite sensitivo deva ter e tenha effectivamente um reflexo physiologico, uma especie de echo no coração, --- reflexo e echo proporcionado á intensidade do acto do proprio appetite. E' só neste sentido que o coração pode chamar-se a sede dos affectos e, em geral, das paixões.\par Todavia, embora se admitta que o orgão principal da sensibilidade \emph{affectiva} é o \emph{cerebro}, não se pode dizer que o genero humano errasse e erre, quando diz que esse orgão é o \emph{coração}. As funcções cerebraes, que acompanham os actos da sensibilidade «/\emph{fectiva}, escapam ao alcance da consciencia \emph{sensitiva}, e por isso não attribuimos esses actos ao cerebro; porque, para attribuirmos um phenomeno a uma causa, devemos ter consciencia d'esse phenomeno e da sua dependencia. Mas se á consciencia escapam as funcções do cerebro, não lhe escapam as do coração; e, como a consciencia nos attesta que os movimentos do coração acompanham os actos da sensibilidade affectiva, e são o reflexo mais immediato e sensivel dos affectos e das paixões, os mesmos movimentos do coração serviram para medir a intensidade dos 'affectos e das paixões e o coração foi tido como a séde, o orgão do amor, sem comtudo se decidir se o coração era a sede \emph{directa} ou \emph{indirecta} dos affectos, isto é, se era o orgão \emph{elicitivo}, ou apenas \emph{manifestativo}, da sensibilidade \emph{affectiva}.\par Orgão \emph{elecitivo} ou \emph{manifestativo} do amor \emph{sensivel}, o coração foi e é tido como o \emph{symbolo} do amor \emph{racional}. E nada mais natural. No homem ha um nexo intimo entre o appetite \emph{sensitivo} e o \emph{racional}. Quando o appetite \emph{racional}, que é a \emph{vontade}, ama e, levado pelo amor, produz ou impera outros actos, move-se tambem o appetite \emph{sensitivo}, e este movimento encontra um reflexo, um echo no coração, que, por isso, se torna uma especie de thermometro, capaz de medir a intensidade do amor \emph{sensivel} e \emph{racional}.\par Pergunta-se: De que deriva o reflexo das paixões no coração?\par A esta pergunta dão os modernos physiologistas a seguinte resposta. Embora as contracções musculares do coração não dependam inteiramente do centro cerebro-espinhal (pois o coração, separado do organismo, continua a palpitar por algum tempo), todavia a intensidade, a celeridade e a regularidade d'essas contraeções podem depender e dependem da influencia do centro cerebro-espinhal. Por quanto, o coração está preso ao grande centro nervoso por dois feixes de fibras nervosas, que pertencem respectivamente ao nervo \emph{paewmo-gastrico} e ao nervo \emph{grande-sympathico}. A excitação do nervo \emph{pnewmo-gastrico} retarda os movimentos do coração; a do \emph{grandesympathico} accelera-os. As percepções, pois, e as imagens cerebraes, que excitam o \emph{pneumo-gastrico}, retardam a circulação do sangue, diminuem a nutrição dos tecidos e das cellulas, e por conseguinte enfraquecem e abatem o organismo. Pelo contrario, os phenomenos physicos e as excitações cerebraes, que influem no nervo \emph{grande-sympathico}, tornam mais activos os movimentos do coração, fazem que o sangue afflua mais abundantemente aos tecidos, e que a nossa vitalidade seja mais energica, os desejos mais ardentes, as emoções mais vivas. Logo os centros nervosos superiores exer- “cem uma influencia real nos movimentos do coração, e motivam a variedade dos movimentos, que estão ligados ás nossas emoções. Ora, sendo os\pgn{144}centros nervosos superiores o verdadeiro orgão de toda a sensibilidade, tanto cognitiva como afectiva, é evidente que as nossas percepções e emoções sensiveis (e, por meio d'estas, as operações da intelligencia e da vontade) hão de influir sobre as funcções do coração.\par \emph{Cl}. \emph{Bernard} resume as ultimas descobertas da physiologia ácerca d'este assumpto nas seguintes palavras:\par « Os sentimentos, que experimentamos, são sempre acompanhados por acções reflexas do coração; é o coração que subministra as condições para a manifestação dos sentimentos, embora estes só residam no cerebro.\par « Nos organismos mais elevados, a vida é uma continua correspondencia entre o systema sanguineo.6 o nervoso; os nossos sentimentos exprimem-se por uma correspondencia entre o coração " o cerebro, --- as duas rodas mais perfeitas da machina viva. Esta correspondencia realiza-se por meio de relações anatomicas muito conhecidas, isto é, pelos nervos \emph{pnewmogastricos}, que trazem as influencias nervosas ao oração, e pelas arterias \emph{carotidas} e \emph{vertebraes}, que levam o sangue ao cerebro.\par \& A sciencia physiologica ensina que, por um lado, o coração recebe realmente a impressão de todos os nossos sentimentos, e que, por outro lado, o coração reage para enviar ao cerebro as condições necessarias para a manifestação d'esses sentimentos. Quando dizemos que \emph{o coração está despedacado pela dór}, affirma-se um phenomeno verdadeiro e real; pois que, se a impressão dolorosa foi repentina, o coração p'ra, e, deixando de enviar sangue ao cerebro, provoca necessariamente a syncope e crises nervosas.\par \& As impressões agradaveis tambem influer' no coração. Quando uma pessoa experimenta uma doce emoção, as palavris, que a suscitaram, atravessam o seu espirito, rapidas como relampago, para cahirem sobre o coração, prevenindo qualquer raciocinio ou reflexão. O sentimento começa a manifestar-se depois de uma breve pausa cardiaca..., e o coração, excitado pelos nervos, reage, dando fortes palpitações, e lançando sangue em abundancia para a cabeça; d'haio rubor do rosto, e certos signaes na physionomia, que traduzem o prazer da alma. Dizer, pois, que o \emph{amor faz palpitar o coração}, não é só uma phrase poetica, mas é tambem uma verdadeira realidade physiologica....\par « Em summa, no homem, o coração é o mais sensivel entre todos os orgãos da vida vegetativa, e é o primeiro a receber todas as influencias da circulação sanguinea. D'ahi se segue que ambos os orgãos principaes da vida (o \emph{cerebro} e o \emph{coração}) devem estar sempre e indissoluvelmente unidos pelos laços da acção e da reacção » (\emph{La science expérim}., Physiologie du coeur),\par Alguns escriptores conciliam a sentença dos modernos physiologistas, que fazem do cerebro o orgão da sensibilidade, ea do genero humano, que viu sempre no coração a sede dos nobres affectos e das paixões ardentes, dizendo que no \emph{cerebro} esta a \emph{origem} da sensibilidade \emph{affectiva}, mas no coração está a \emph{sede} d'essa sensibilidade, porque é para o coração que convergem todas as fibras do systema nervoso central e é nelle que se manifestam os phenomenos da sensibilidade e da vida psychica.}.

%%>>> t108

\tese{108}{Faculdade locomotiva.}
Faculdade \emph{locomotiva} é aquella pela qual o homem se transporta de um logar para outro, com o fim de alcançar o bem e fugir do mal, ou de remover \pg{143} \pg{144}  os obstaeulos que se oppõem a essa consecução OU a essa fuga. --- Os orgãos d'esta faculdade são os \emph{ossos} e os \emph{musculos}. Estes dizem-se orgãos \emph{activos}; aquelles, \emph{passivos} \footnote{\emph{Fritsch} e \emph{Hitsig}, celebres physiologistas allemaes, demonstram á evidencia\pgn{145} que os \emph{centros motores}, os quaes excitam os orgãos da faculdade \emph{locomotiva} a realizar o seu acto, residem \emph{quasi exclusivamente} no \emph{lóbo parietal} e na \emph{metade posterior do lóbo frontal}, e que a sensibilidade reside na parte posterior dos hemispherios e especialmente nos lóbos occipital e temporoesphenoidal. Dissemos --- \emph{quasi exclusivamente}, porque outros \emph{centros motores}, além dos que residem no /\emph{óbo parietal} e no \emph{frontal}, concorrem para o exercicio da faculdade \emph{locomotiva}; pois as experiencias de \emph{Flourens} mostram que, extrahidos os lóbos frontaes, não cessam inteiramente os movimentos espontaneos. --- Esses outros centros motores residem na \emph{camada cortical} do cerebro; pois, como diz Wundt,se faltar \emph{toda} \& camada cortical, não poderá manifestar-se nenhum movimento espontaneo.}.

%%>>> t109

\pg{145} \tese{109}{A faculdade locomotiva é distincta do appetite sensitivo.}
\textsc{Duas} faculdades são realmente distinctas entre si, quando uma póde estar sem outra. Ora o appetite \emph{sensitivo} póde estar sem a faculdade \emph{locomotiva}, como se vê --- não só nos animaes menos perfeitos, que são dotados de \emph{appetite} e não de \emph{locomoção}, --- mas tambem nos animaes mais perfeitos, que, embora possuam a faculdade \emph{locomotiva}, comtudo não a podem exercer por algum impedimento. Logo a faculdade \emph{locomotiva} é distincta do appetite \emph{sensitivo} \footnote{O appetite \emph{sensitivo} e os orgãos dos sentidos podem estar intactos, e todavia póde ser que o animal não se possa determinar espontaneamente ao movimento. Havendo uma lesão no cerebello ou nos centros nervosos da base do cerebro, o animal, não obstante o seu desejo de caminhar, soffrerá vertigens, cahirá a cada instante ou andará sempre em volta de si mesmo.\par S. Thomaz diz: « Ainda que o sentido e o appetite sejam principios motores nos animaes perfeitos, todavia o sentido e o appetite, emquanto aos, não chegam a produzir o movimento, se não se lhes accrescentar alguma energia. Por quanto nos animaes immoveis ha sentido e appetite, mas não ha força motriz. Esta força encontra-se não só no appetite e no sentido, emquanto comandam o movimento, mas tambem nos proprias partes do corpo, para que possam obedecer ao appetite. E que isto seja assim, mostra-o a experiencia; porque quando os membros estão naturalmente indispostos, não obedecem ao imperio do appetite e não se movem » (\emph{Sum}. \emph{Th}., p. L, q. (8, a. 1 ad 4). --- A physiologia demonstra que os musculos, para se contrahirem, devem receber dos centros nervosos um impulso, ou uma corrente do agente nervoso, transmittida pelos nervos motores.}.

%%>>> t110

\tese{110}{A faculdade locomotiva no homem e no animal.}
No animal, a faculdade \emph{locomotiva} segue immediatamente o acto do appetite \emph{sensitivo}, determinado por um juizo \emph{énstinctivo} da força \emph{estimativa}; assim a ovelha, vendo o lobo vir ao seu encontro, \emph{teme}, e immediatamente deita a fugir. No homem, porém, a faculdade \emph{locomotiva} está subordinada ao acto da vontade, embora esta seja por vezes contraria ao appetite \emph{sensitivo} \footnote{Advertimos que a faculdade \emph{locomotiva}, para obedecer ao imperio da vontade, deve ter os seus orgãos bem dispostos e em estado normal;\par 10 Vor. II.\pgn{146}aliás poderia acontecer, como muitas vezes acontece, que se não realisassem os movimentos imperados pela vontade.\par Advertimos tambem que a vontade, embora possa ser o principio de certos movimentos locaes, não é, comtudo, principio \emph{immediato}; porque ella, geralmente, só determina o movimento por meio do appetite \emph{sensitivo}. E nada mais coherente. A vontade, como a intelligencia, tende naturalmente para o \emph{universal} ou \emph{immaterial}; ao passo que o movimento \emph{local} tem por termo um bem \emph{material} ou \emph{particular}, que foi directamente percebido pelos sentidos e que é objecto proprio e proporcionado do appetite \emph{sensitivo}.\par \emph{Descartes} ensinou que a faculdade \emph{locomotiva} não é uma faculdade especial e distincta das outras; porque, dizia elle, o movimento \emph{local}, no animal, deriva immediatamente do appetite \emph{sensitivo}, e, no homem, da \emph{vontade}.\par Esta opinião é falsa. A faculdade \emph{locomotiva} distingue-se tanto do appetite \emph{sensitivo}, como da \emph{vontade}. --- Distingue-se do appetite \emph{sensitivo}; porque ás vezes nos afastamos do bem sensivel, para o qual tende o mesmo appetite. --- Distingue-se da \emph{vontade}, porque muitas vezes não podemos executar um movimento, que a vontade quereria se executasse.\par Digamos alguma coisa ácerca das operações da faculdade \emph{locomotiva}, commummente chamadas \emph{movimentos}.\par \emph{Movimento} é o acto da faculdade \emph{locomotiva}, pelo qual \emph{o} homem (e, em geral, o animal) se transporta de um logar para outro, ou desloca as differentes partes do seu organismo, em relação umas ás outras.\par Os \emph{movimentos} são \emph{espontaneos}, \emph{reflexos}, e \emph{automaticos}.\par a) O movimento \emph{espontaneo} é provocado por um acto do appetite \emph{sensitivo}, com previo conhecimento das faculdades \emph{perceptivo-sensitivas}. --- Este movimento costuma chamar-se tambem \emph{voluntario}, pois que no homem á faculdade locomotiva está subordinada á vontade; mas é mais proprio chamar \emph{espontaneo} o movimento que depende \emph{immediatamente} do \emph{appetite sensitivo}, e \emph{voluntario} o que deriva \emph{immediatamente da vontade}. --- Aos movimentos \emph{espontaneos} reduzem-se os \emph{instinciivos}, como ao genero se reduz a \emph{especie}. --- O movimento \emph{espontaneo} parte da \emph{camada cortical}.\par \emph{b}) O movimento \emph{refleso} é como que uma resposta a uma excitação \emph{peripherica} ou \emph{exterior}, sem intervenção do \emph{appetite sensitivo} ou da \emph{vontade}. --- O mais simples movimento \emph{reflexo} é o \emph{medullar}, que é provocado, por uma excitação, a qual, terminando num centro sensitivo, se \emph{reflecte} sobre o respectivo centro motor. Os mais complexos \emph{reflexos} são os \emph{coordenados}, como são, por exemplo, os que servem para a deglutição. --- Os \emph{reflexos}, tanto \emph{simples} como \emph{coordenados}, podem ser \emph{conscientes} ou \emph{inconscientes}. --- Estes movimentos partem dos centros \emph{sub-corticaes}. : c) O movimento \emph{automatico} é o que tem por causa um \emph{estimulo interno}, independentemente do acto do appetite \emph{sensitivo} ou da \emph{vontade}. Assim o movimento do coração é \emph{automatico}; como tambem \emph{automatico} é o movimento rythmico da inspiração e expiração do apparelho respiratorio. --- Como se vê, o movimento \emph{automatico} e O \emph{reflexo} concordam em serem ambos independentes do acto do appetite \emph{sensitivo} e da \emph{vontade}; mas differem em\pgn{147}que a causa do movimento \emph{reflexo} é uma excitação \emph{exterior}, ao passo que a causa do movimento \emph{automatico} é um \emph{estimulo interno}.\par * Estes tres movimentos --- \emph{espont«aneo}, \emph{reflexo} e \emph{automatico} --- encontram-se no homem, e nos animaes.\par Todavia \emph{Descartes}, \emph{Mallebranche}, \emph{Huxley}, \emph{Richet}, seguidos pelos modernos \emph{muaterialistas} e \emph{positivistas}, afirmam que \emph{todos} os movimentos, que se realisam no homem, são \emph{reflexos} ou \emph{automaticos}, e que por isso não ha nenhum movimento \emph{espontaneo}.\par Esta affirmação é falsa pelas seguintes razões:\par 1.º) A consciencia attesta que em nós se realisam movimentos, que não têm a sua razão sufficiente numa simples excitação mechanica ou physicochimiea, mas que derivam de um juizo da \emph{estimativa} e de uma correspondente inclinação ou emoção do appetite \emph{sensitivo}; assim, depois de termos percebido que um objecto nos é \emph{util} ou \emph{agradavel}, sentimos uma tendencia para elle, e, em conformidade com essa tendencia, executamos os movimentos, que são necessarios para alcançar o mesmo objecto. Ora os movimentos, que derivam da \emph{percepção} e do \emph{appetite sensitivo}, chamam-se e são \emph{espontaneos}. Logo dão-se no homem movimentos \emph{espontaneos}.\par 2.2) Se todos os movimentos do homem, como tambem do animal, fossem phenomenos meramente \emph{mecAhanicos}, deveriam todos estar sujeitos ás leis da \emph{mechanica}. Ora nos seres sensitivos realisam-se movimentos que, embora dependam de \emph{algum modo} do meio exterior, todavia não estão sujeitos ás leis fataes da \emph{mechanica}; pois que o animal, nos seus movimentos, muda de velocidade e de direcção; embora não encontre obstaculo algum, o que é contrario á lei da \emph{inercia}. Logo nem todos os movimentos do homem, como tambem do animal, são \emph{mechanicos}, mas alguns são \emph{espontaneos}.\par Dissemos que os movimentos dos seres sensitivos dependem de \emph{algum modo} do meio exterior; porque este, pelo facto de \emph{influir} nas faculdades \emph{perceptivas} do ser sensitivo, é apenas a causa remota e \emph{parcial} do movimento, pois a causa \emph{provima} e \emph{determinante} é o proprio animal.\par 8. Thomaz, analysando o movimento \emph{espontaneo}, distingue tres causas: --- a causa \emph{dirigente}, que é constituida pela \emph{imaginação} e pela \emph{estimativa}, emquanto estas duas faculdades mostram ao appetite \emph{sensitivo} o que convém, ou o que não convém, --- a \emph{causa imperante}, que é o appetite \emph{sensitivo}, emquanto aectüa nos nervos motores e assim determina ao acto a faculdade locomotiva, --- a causa \emph{instrumental}, que executa o movimento e que é a energia diffusa nos musculos, nos ossos e nos nervos dos membros. (Cf. opus. \emph{De potentiis animae}, c. V ) s}.

111; \textbf{Faculdade vocal}. --- Faculdade \emph{voca}? é aquella pela qual o homem exprime e communiea os seus \emph{conhecimentos} e as suas \emph{paixões} por meio de um certo \emph{som}, que sai da boeea.

Eo som, que sai da boeca do homem, como tambem da do ani- - mal, chama-se \emph{voz}. --- O orgão principal da \emph{voz} é a \emph{larynge} \footnote{A \emph{larynge} é uma cavidade de paredes cartilago-membranosas, aberta superiormente na \emph{pAarynge} e communieando inferiormente com a \emph{trachêa}. E' formada de quatro cartilagens: a \emph{thyroidéa}, que tem a forma de escüdete, está situada na parte anterior, e apresenta uma saliencia notavel, especialmente no sexo maseulino, que se chama vulgarmente nó \emph{da garganta} ou \emph{pomo de Adão}; --- a \emph{cricoidéa}, que tem a forma de annel, e está collo-\pgn{148}\textgreater{} cada na parte inferior; --- as duas \emph{arythenoidéas}, que estão collocadas posteriormente e representam duas pequenas pyramides, cuja base se articula com a \emph{cricoidéa}. --- Da base das \emph{arythenoidéas} nascem uns ligamentos, --- dois de cada lado, --- que se inserem na \emph{thyroidea} e se chamam \emph{cordas} vocaes, sendo duas \emph{superiores} e duas \emph{inferiores}. O espaço, comprehendido entre as quatro cordas vocaes, tem o nome de \emph{glotte}. Sobre a abertura superior da glotte ha uma lingueta cartilago-fibrosa, que serve de valvula e\par que se chama \emph{epiglotte}.\par \emph{A voz} é produzida pela acção do ar sobre as cordas vocaes. O ar expirado, vindo dos pulmões, imprime ás cordas vocaes vibrações mais ou menos rapidas, e produz sons mais agudos ou mais graves, segundo as cordas vibrantes são mais ou menos curtas e estão mais ou menos tensas; pois está demonstrado na \emph{acustica} que os sons produzidos por cordas vibrantes são tanto mais agudos, quanto mais curtas e mais tensas ellas estão. D'aqui se vê que as mulheres e as creanças, cuja larynge é mais pequena e as cordas vocaes são mais curtas do que no homem, devem ter, como effectivamente têm, uma voz muito mais aguda que a Q' elle.\par Todavia as cordas vocaes podem mais ou menos encurtar-se e alongar-se, por virtude das contracções dos musculos da larynge; produzindo assim, á vontade; sons mais agudos ou mais graves.}.

%%>>> t112

\tese{112}{A faculdade vocal no homem e no animal.}
A faculdade vocal, embora seja commum ao homem e ao animal, todavia \pg{146} \pg{147} \pg{148}  naquelle é mais elevada do que neste. Por quanto, O animal, por ser irracional, não póde articular, nem formar uma \emph{palavra}, que É voz \emph{articulada}; ao passo que o homem, dotado de razão, póde articular a sua voz e formar \emph{palavras}, que expri-

mam os seus pensamentos \footnote{O homem é o unico animal que possue a faculdade de modificar os . \textbar{} diversos sons, produzidos pelas cordas vocaes, formando \emph{palavras}, para expri- 4: mir os seus pensamentos.\par A \emph{palavra} ou \emph{linguagem articulada} compõe-se de sons produzidos ---\pgn{149}uns pelas cordas vocaes --- e outros pela passagem do ar atravez dos orificios, que estão situados acima da \emph{glotte}, e cuja forma podemos modificar á vontade. --- Os sons produzidos pelas cordas vocaes constituem as \emph{vogaes}, e o3 sons produzidos pela expulsão do ar, que atravessa a \emph{glotte}, constituem as \emph{consoantes}. Estes sons, directamente associados, formam as \emph{syllabas}, de que se compõem as \emph{palavras}.}.

--- ---------\$9---2---------------

%%>>> cap-4

\chapter*{CAPITULO TERCEIRO\\[0.3em]\large FACULDADES INTELLECTUAES DO HOMEM}
\addcontentsline{toc}{chapter}{CAPITULO TERCEIRO --- FACULDADES INTELLECTUAES DO HOMEM}

\textsc{SuwMaRIO}0: --- Intelligencia, sua natureza e varios nomes. --- Vontade, sua

natureza e liberdade. --- A liberdade da vontade humana e O determi-

nismo. --- Relações entre a intelligencia e á vontade.

\textbf{Intelligencia}, \textbf{sua natureza} e \textbf{varios nomes}.

%%>>> t113

\tese{113}{Intelligencia.}
\emph{Intelligencia} é a faculdade espiritual, pela qual o homem percebe a \emph{essencia} das coisas. --- Percebe-se a \emph{essencia} de uma coisa, não só quando se indicam \emph{todos} \pg{149}  os seus elementos \emph{especificos}, mas tambem, embora menos perfeitamente, quando se aponta um ou outro dos seus elementos mais \emph{genericos} ou \emph{transcendentaes}.

A \emph{intelligencia} existe no homem; porque a consciencia nos attesta que, emquanto as faculdades sensitivas param na superficie externa dos objectos e só percebem as qualidades sensiveis, ha em nós uma energia superior que penetra até ao intimo d' esses objectos, e descobre notas e caracteres, que escapam ao alcance dos sentidos \footnote{O proprio nome de \emph{intelligencia} significa uma força que, transpondo as qualidades \emph{materiaes}, penetra até ao fundo das coisas, e \emph{intus legit}, lé no intimo d'ellas o que não pode ser percebido pelos sentidos. Diz S. Thomaz: \& Est inter alias lata differentia sensuum ac intellectuum, quod illi in cognitione \emph{ecternorwm accidentium}. \emph{sensibilium} sistunt, intellectus vero ex acoidentium cognitione ad contemplanda \emph{ea quae sub accidentibus cadunt atque latent} ingreditur. Unde intellectus vocatur quasi \emph{éntws legens} » (\emph{De anima}, lib. 4, lect. 4. --- C£. \emph{Sum}. \emph{Theol}., p. III, q 8, a. 1).\par A intelligencia não só percebe a essencia das \emph{substancias}, mas tambem a essencia das \emph{qualidades} e, em geral, dos \emph{accidentes}. Diz o mesmo S. Thomaz: 4 Intellectus intelligit absoluto \emph{cujusque rei quidditatem}, sive \emph{essentiam} per seipsam, puta \emph{quid est homo}, vel \emph{quid est album}, vel quid aliud hujusmodi » (\emph{Perihermenias}, lib. 1, lect. 8). --- Todavia quando se diz que a intelligencia percebe a \emph{essencia} das \emph{coisas}, queremos denotar --- \emph{primariamente} a essencia das \emph{substancias}, que são coisas \emph{absolutas} e se entendem por si mesmas, --- e \emph{secundariamente} a essencia dos \emph{accidentes}, que são coisas \emph{relativas} e só se entendem pela sua relação á propria substancia, e por isso são mais \emph{do ente} do que \emph{entes} (Cf. S. Thom., \emph{Sum}. \emph{Theol}, p. I, q. 90, a. 2. --- \emph{De ente et essentia}, cc. II, VIL)}.

%%>>> t114

\tese{114}{Objecto da intelligencia.}
\emph{Objecto} da intelligencia é a \emph{essencia} das coisas. --- Divide-se em \emph{adequado} e \emph{proporcionado}. O \emph{adequado} é a \emph{essencia} de todas as coisas, considerada no seu mais fundamental e commum elemento, que é o \emph{ente}. O \emph{proporcionado} é a \emph{essencia} das coisas \emph{materiaes}, \emph{abstrahida das condições materiaes}.

a) \emph{Objecto adequado da intelligencia}. --- Objecto \emph{adequado} da intelligencia é tudo o que ella pode attingir considerada \emph{como tal}, isto é, considerada emquanto á sua natureza de faculdade \emph{perceptivo-espiritual}, sem se attender á sua actual condição ou modo de existencia. Ora a nossa intelligencia, \pg{150}  considerada como faculdade \emph{perceptivo-intellectual}, pode penetrar até ao intimo de todas as coisas, que lhe forem convenientemente propostas, e perceber a \emph{essencia} pelo menos no seu mais fundamental e commum elemento ou predicado, que é o \emph{ente}; pois, possuindo a idea abstractissima de \emph{ente}, deve esta faculdade poder attingir e perceber todas as coisas, nas quaes essa idea se realiza e resplandece. Logo o objecto \emph{adequado} da intelligencia, emquanto \emph{tal}, é a \emph{essencia} de todas as coisas, considerada no seu mais fundamental e commum elemento, que é o \emph{ente}.

Como a \emph{essencia} das coisas é objecto da intelligencia emquanto e só porque é immaterial, pode tambem dizer-se que o objecto \emph{adequado} d'esta faculdade é o \emph{immaterial} \footnote{\emph{Objecto} da intelligencia é a \emph{essencia} das coisas. Diz S. Thomaz: « Objectum intelligentiae \emph{est ipsa essentia rei} » (\emph{De verit}., a. 4). A \emph{essencia} das coisas é o objecto \emph{formal} da intelligencia, porque a \emph{essencia} é a razão pela qual a intelligencia attinge as coisas; e as \emph{coisas} mesmas, que encerram o objecto \emph{formal}, isto é, a \emph{essencia}, constituem o objecto \emph{material} essa faculdade. --- O objecto \emph{formal é adequado} ou \emph{proporcionado}. \emph{O adequado} chama-se \emph{generico}, porque é attingido pela intelligencia \emph{enquanto tal}, e por isso convém a \emph{todas} as intelligencias; e o \emph{proporcionado} diz-se \emph{especifico}, porque é proprio da intelligencia \emph{humana}, e não da divina nem da angelica.\par O objecto \emph{adequado}, ou \emph{generico}, da intelligencia humana é a \emph{essencia} das coisas, considerada no seu mais commum e universal elemento, que é o \emph{ente}, --- Desenvolvamos este ponto. Todas as coisas têm uma \emph{essencia}, pela qual são o que são e pertencem a uma certa e determinada classe de seres finitos. A essencia compõe-se naturalmente de varios elementos: uns \emph{proprios} e outros communs. Pelos elementos \emph{proprios}, uma coisa distingue-se de todas as outras; pelos elementos \emph{communs}, convém com ellas. Os elementos communs são muitos e diversos: uns são \emph{especificos}, outros \emph{genericos}, outros \emph{transcendentaes}. Pelos elementos \emph{especificos}, uma coisa convém com as outras da mesma \emph{especie}; pelos \emph{genericos}, convém com as outras do mesmo \emph{genero}; pelos \emph{transcendentaes}, convém com \emph{todas} as coisas. Ora, entre os elementos \emph{transcendentaes}, \emph{o} mais \emph{intimo} e \emph{commum} é o \emph{ente}. O \emph{ente} --- é o elemento mais \emph{intimo}, porque nada mais \emph{intimo}, ou \emph{essencial}, ás coisas do que o serem \emph{alguma coisa}, e o ser alguma coisa é significado pelo \emph{ente}; --- é o elemento mais \emph{commum}, porque as coisas, embora diversas entre si pelos seus elementos \emph{especificos}, convem todas em que são \emph{entes}.\par Por isso --- 1.9) o \emph{ente} é a razão \emph{formal} e \emph{communissima}, pela qual a intelligencia percebe todas as coisas, porque todas as coisas são intelligiveis emquanto são \emph{entes}. Diz S. Thomaz: « Secundum hoc unumquodque est intelligibile, \emph{inquantum est actu}; unde \emph{ens} est \emph{proprium objectun} intellectus » (\emph{Sum}. \emph{Th}., p. I, q. 5, a. 2); « Intellectus respicit suum objectum secundum communem rationem \emph{entis} » (\emph{Sum}. \emph{Th}., p. T, q. 79, a. 7). Note-se que o \emph{ente} é objecto \emph{proprio} da intelligencia, emquanto esta se considera \emph{genericamente}, isto é, emquanto é intelligencia. --- 2.9) O \emph{ente} é o \emph{primeiro intelligivel}, --- não so porque elle é a razão \emph{formal}, pela qual todas as coisas são intelligiveis,\pgn{151}assim como a côr, que é o objecto proprio da vista, é o primeiro visivel, porque todas as coisas são visiveis emquanto são coloridas, --- mas tambem porque a nossa intelligencia passa do que é mais universal e commum para o que o é menos, e nenhum conceito é mais universal e commum do que o de ente. S. Thomaz diz: « Ens est primum intelligibile » (\emph{Sum}. \emph{T}7., p. I, q. 5, a. 2); --- « Ens est primum, quod cadit in apprehensione simpliciter » (\emph{De veritate}, a. 1) --- 8.9) O conceito de \emph{ente} é tão simples e fundamental, que se acha incluido em todos os outros conceitos, porque todos se resolvem, em ultima analyse, no conceito de \emph{ente}. Diz S. Thomaz: « Entis intellectus includitur in omnibus quaecumque quis apprehendit » (\emph{Sum}. \emph{TA}., I-II, q. 94, a. 2).\par Para a exacta comprehensao da doutrina exposta, devemos fazer as seguintes advertencias. --- \emph{a}) Quando dizemos que a.intelligencia percebe todas as coisas emquanto são \emph{entes}, não entendemos que, na percepção do objecto, só se perceba a communissima nota de \emph{ente}, como que separada ou abstrahida das outras notas especificas ou differenciaes, porque isto seria absurdo e contrario á experiencia, mas entendemos que as coisas são \emph{Intelligiveis} emquanto são \emph{entes}, \emph{e} que por isso a esphera da \emph{intelligibilidade} se extende até onde se extende a esphera de \emph{ente}. --- \emph{b}) Quando se afirma que a inlelligencia percebe um objecto não só na communissima nota de \emph{ente}, mas tambem nas suas notas \emph{especificas} ou \emph{differenciaes}, não se deve entender que essas notas \emph{especificas} se \emph{addicionam}. á nota \emph{commum} do ente, como a differenca \emph{especifica} so addiciona ao \emph{genero} (porque a differenca especifica, que se une ao \emph{genero}, não é o \emph{genero}, mas é diversa d'elle; ora nao ha coisa ou differenca que não seja \emph{ente}); --- mas devemos entender que aquelle objecto, que no principio tinha sido percebido sob o aspecto indeterminado e confuso de \emph{ente}, depois é percebido mais claramente nos seus elementos \emph{especificos} e nos seus caracteres \emph{dZfferenciaes} (Cf. S. 'Th., \emph{Sum}. \emph{Theol}., p. I, q. 85, a. 8; --- \emph{Summa totius Logicae Aristotelis}, tract. IT, c. 2). --- \emph{c}) Sustentando que o \emph{ente} é objecto \emph{adequado} da intelligencia, não queremos dizer que a intelligencia \emph{actualmente} perceba \emph{todo} o ente e pelo \emph{mesmo modo}; mas entendemos que a nossa intelligencia \emph{pode} conhecer todas as coisas, nas quaes reconhece a nota de \emph{ente}, comtanto que lhe sejam \emph{convenientemente} propostas, e que as pode conhecer por um modo \emph{directo} ou \emph{indirecto}, pela especie \emph{propria} ou \emph{analoga}, pelos elementos \emph{constitutivos} ou pelo \emph{e/feitos}, etc. --- \emph{d}) Tambem os sentidos percebem o \emph{ente}. Mas essa percepção ó differente da percepção da intelligencia. Os sentidos percebem o \emph{ente}, mas por um modo \emph{concreto}, \emph{e} não percebem a razão ou a \emph{significação} de ente, não sabem o que é ser ente, ao passo que\pgn{152}a intelligencia conhece o \emph{ente} por um modo \emph{abstracto}, \emph{e} percebe o que quer dizer a noção de \emph{ente}. Diz S. Thomaz: « Quamvis esse sit in rebus sensibilibus, tamen \emph{rationem essendi}, vel \emph{intentionem entis} sensus non apprehendit » (1. \emph{Dist}. 19, q. 5, a. 1 ad 6).\par Notamos, por ultimo, que a \emph{essencia}, sendo a razão pela qual a intelligencia attinge as coisas, é o verdadeiro objecto \emph{formal} e \emph{adequado} desta faculdade. Todavia, sendo a \emph{essencia} das coisas objecto da intelligencia emquanto é \emph{immaterial}, porque o \emph{material}, emquanto \emph{tal}, não se pode unir com uma faculdade \emph{espiritual}, segue-se que o \emph{immaterial} se pode chamar tambem o objecto \emph{formal} e \emph{adequado} da intelligencia.}.

b) \emph{Objecto proporcionado da intelligencia}. --- Objecto \emph{proporcionado} da intelligencia é o que é connatural a esta faculdade, considerada na sua actual condição ou modo de existencia, e por isso é conhecido \emph{anteriormente} aos mais objectos e pela especie \emph{propria}, e é a razão \emph{pela qual} e \emph{segundo} a \emph{qual} todos os outros objectos são conhecidos. Ora a \emph{essencia} das coisas materiaes, \emph{abstrahida} das condições materiaes, --- a) é objecto connatural á condição actual da intelligencia humana, porque uma faculdade \emph{immaterial}, mas existente na materia, como é a nossa intelligencia no estado actual, deve ter por objecto \pg{152}  connatural uma coisa \emph{immaterial} mas existente na materia, e tal é a essencia \emph{abstracta} das coisas \emph{materiaes}, --- \emph{b}) é conhecida \emph{anteriormente} a todos os outros objectos, porque o exercicio da intelligencia segue-se immediatamente ao dos sentidos, e estes \emph{primeira} e \emph{directamente} percebem as coisas \emph{materiaes}; e é conhecida pela especie \emph{propria}, porque esta especie é produzida pela acção da mesma essencia na faculdade intellectual, e por isso não pode deixar de ser \emph{propria}; --- c) é a razão \emph{pela qual} e \emph{segundo a qual} a intelligencia percebe todas as outras coisas, porque a experiencia nos attesta que nos elevamos ao conhecimento do immaterial pelo conhecimento do material, e que não conhecemos o invisivel se não á \emph{semelhança} do visivel. Logo o objecto \emph{proporcionado} da nossa intelligencia no estado presente é a \emph{essencia} das coisas \emph{materiaes}, \emph{abstrahida} das cir-

\emph{A} cumstancias ou condições materiaes.

Como a essencia \emph{abstracta} das coisas materiaes é \emph{universal}, pode dizer-se que o objecto \emph{proporcionado} da nossa intelligencia no estado presente é o \emph{universal} \footnote{Desenvolvamos cada uma das partes da proposição \emph{menor}.\par a) A \emph{essencia} das coisas \emph{materiaes}, \&bstrahida de todas as condições \emph{materiaes}, é o objecto \emph{connatural} ou \emph{proporcionado} da intelligencia humana no estado presente. Oucamos o ensino do Angelico Doutor. Diz elle assim: « O conhecimento faz-se emquanto o objecto conhecido se une com a faculdade cognitiva. Unindo-se com a faculdade, o objecto aecomoda-se á natureza e ao modo de ser d'ella. Por isso todo o conhecimento é proporcionado á natureza e ao modo de ser da faculdade cognitiva. Se, portanto, o modo de ser de um objecto cognoscivel excede o modo de ser da faculdade cognitiva, é claro que o conhecimento d'aquelle objecto excede o alcance da faculdade. \emph{Cognitio contingit secundum quod cognitum est in cognoscente}. \emph{Cognitwm}. \emph{autem est in cognoscente secundum modum cognoscentis}. \emph{Unde cujuslibet cognoscentis cognitio est secundum modum suae naturae}. \emph{Si igitur modus essendi alicujus rei cognitae excedit modum naturae cognoscentis}, \emph{oportet quod cognitio illius vei sit supra naturam illius cognoscentis} » (\emph{Sum}. \emph{Theol}, p. I, q. 12, a. 4). Logo na opinião de S. Thomaz, a natureza do objecto e a natureza da faculdade devem convir no mesmo modo de ser; --- a immaterialidade do\pgn{153}objecto deve por tal forma corresponder ao genero de immaterialidade da faculdade, que, se a immaterialidade da faculdade é \emph{pura}, isto é, não precisa de imagens sensiveis pars entender, tambem a immaterialidade do objecto deve ser \emph{pura}, isto é, tambem o objecto deve ser \emph{puramente} immaterial, e, se a immaterialidade da faculdade é \emph{mista}, isto é, se esta unida com as imagens sensiveis, tambem \emph{mista} deve ser a immaterialidade do objecto, emquanto deve ser uma coisa \emph{immaterial} unida com algum elemento \emph{material}. --- Postos estes principios, é claro que, sendo a intelligencia humana uma faculdade \emph{immaterial}, mas existente na \emph{materia} (emquanto é uma faculdade \emph{espiritual} da alma, que é forma substancial do \emph{corpo}), o seu objecto \emph{proprio} ou \emph{proporcionado} deve ser o \emph{immaterial} existente na \emph{materia}, e tal é a \emph{essencia abstracta} das coisas \emph{materiaes}. Diz S. Thomaz no logar apontado: « Unde per intellectum connaturale est nobis cognoscere naturas, quae quidem non habent esse nisi in materia \emph{individuali}, sed secundum quod \emph{abstrahuntur ab ea} per considerationem intellectus ». E noutro logar: « \emph{Operatio proportionatur virtuti et essentiae}: intellectivum autem hominis est in sensitivo, et ideo propria operatio ejus est intelligere \emph{éntelligibilia} in \emph{phantasmatibus} » (\emph{De Mem}. \emph{et Reminis}., 1. 4).\par b) \emph{A} essencia \emph{abstracta} das coisas \emph{materiaes} --- é conhecida \emph{anteriormente} aos outros objectos, --- e é conhecida pela sua \emph{propria} especie. Na verdade, o objecto connatural ao modo de existencia de uma faculdade é o \emph{termo} para o qual a mesma faculdade está essencialmente ordenada, e por isso é o que \emph{primeiramente}, antes de todos os outros objectos, move a faculdade e \& reduz ao acto, e portanto deve ser conhecido pela \emph{propria} especie, e não por uma especie \emph{alheia} (aliás não seria elle o \emph{primeiro} a mover a faculdade e a ser por ella percebido). É o que nos attesta a experiencia interna. --- A essencia \emph{abstracta} das coisas \emph{materiaes} é conhecida \emph{anteriormente} a todos os outros objectos. O exercicio da nossa intelligencia suppõe e exige necessariamente o exercicio dos sentidos, emquanto estes fornecem o material do pensamento; por isso a primeira idea da intelligencia corresponderá á primeira percepção dos sentidos. Ora o que os sentidos primeiramente percebem é o ente material. Logo a primeira idea da intelligencia refere-se ao ente material; ou, por outras palavras, a essencia \emph{abstracta} das coisas materiaes é o \emph{primeiro} objecto da intelligencia humana. S. Thomaz diz: « \emph{Primum}, quod intelligitur a nobis, secundum statum praesentis vitae, est \emph{quidditas}\pgn{154}\emph{ rei materialis}; quae est nostri intellectus \emph{objectum}, ut multoties dictum est » (\emph{Sum}. \emph{Theol}, p. I, q. 88, a. 8; --- Cf. p. I, q. 84, a. 7, --- q. 85, a. 1, --- q. 88, a. 2 ad 2). --- E esta essencia deve ser e é percebida pela especie \emph{propria}, e não por uma especie \emph{alheia}; porque, se fosse percebida por uma especie \emph{alheia}, nao seria o \emph{primeiro} objecto do conhecimento intellectual. Com effeito, a imagem ou especie, que nos representa a essencia das coisas materiaes, é produzida .pela mesma essencia, e é tão propria de uma essencia, que não pode representar uma outra de especie differente.\par 0) A essencia \emph{abstracta} das coisas materiaes é á razão \emph{pela qual} e \emph{segundo} a \emph{qual} conhecemos todos os mais objectos. É um corollario das proposições precedentes. O objecto que corresponde é indole e ao modo de existir de uma faculdade --- não só a o \emph{primeiro} a ser conhecido pela propria faculdade, --- mas é o \emph{unico} que actua \emph{directamente} na faculdade e a reduz ao acto. D'ahi se segue que todos os outros objectos, que a faculdade conhece, não os conhece, nem os pode conhecer senão \emph{por meio} e á \emph{semelhança} esse objecto \emph{primario}. \emph{E} o que nos attesta a consciencia com relação aos nossos conhecimentos intellectuaes. --- A intelligencia humana \emph{primaria} e \emph{directamente} percebe a essencia das coisas materiaes, e todos os outros objectos, que conhece, quer superiores quer inferiores, conhece-os unicamente \emph{em relação} a esse objecto primario e directo. É certo que a nossa intelligencia conhece as coisas puramente immateriaes ou espirituaes, mas tambem é certo que emtanto conhece essas coisas, emquanto as compara com as coisas materiaes, applicando ás immateriaes, ou negando d'ellas, os elementos e as notas, que encontrou nas .materiaes; e por isso o conhecimento das coisas puramente espirituaes ou immateriaes --- nao é \emph{proprio} mas \emph{analogo}, --- não é \emph{adequado} mas \emph{inadequado}. E, como o modo por que se chamam as coisas corresponde ao modo por que ellas se conhecem, d'ahi se segue que os nomes, que significam as coisas puramente immateriaes, são tambem tirados das'coisas \textbf{materiaes}. Diz S. Thomaz: « Potentia cognoscitiva proportionatur cognoscibili. Unde intellectus humani, qui est corpori conjunctus, \emph{proprium}. objectum est quidditas sive natura in materia corporali existens; et \emph{per luiusmodi naturas}\pgn{155}\emph{visibilium rerum} etiam in \emph{invisibilium} rerum \emph{aliqualem} cognitionem \emph{ascendit} » (\emph{Sum}. \emph{Theol.,p}. I, q. 84, a. 7). --- E não só a nossa intelligencia se eleva das coisas materiaes para as immateriaes, e representa estas por meio de especies e palavras tiradas d'aquellas; mas tambem representa as coisas immaiteriaes \& \emph{semelhança} das materiaes, isto é, envolvidas numa imagem sensivel, pelo menos na imagem da \emph{palavra}.\par Dizemos no texto que o \emph{universal} se pode chamar tambem e se chama objecto \emph{proporcionado} da intelligencia. Mas, repetimos, quando se diz que o objecto da intelligencia é o \emph{universal}, não devemos entender que esta faculdade tenha por objecto uma coisa \emph{imaginaria}, \emph{ficticia}, \& qual nenhuma realidade corresponda. No conceito \emph{universal} devemos distinguir a \emph{forma} e o \emph{conteudo}. \emph{A forma} é o caracter do \emph{universalidade}; o \emph{conteudo} é á propria \emph{essencia} das coisas, revestida de \emph{universalidade} pela abstracção das condições materiaes ou individuaes. Este conteudo é \emph{real}, \emph{realissimo}; existe verdadeiramente nas coisas, embora por um modo diverso d'aquelle por que existe na intelligencia. Aqui trata-se, do universal \emph{directo}, e não do \emph{reflexo}. (Cf. \emph{Sum}. \emph{Theol}., p. 1, q. 85, a. 2 ad 2).\par Agora devemos fazer algumas advertencias.\par 1.*) O objecto \emph{proporcionado} chama-se tambem \emph{proprio} e \emph{primario}. Expliquemos os termos. --- Objecto \emph{proprio} de uma faculdade, diz S. Thomaz, é id \emph{quod est primo} et \emph{per se cognitum} (\emph{Sum}. \emph{Theol}., p. I, q. 85, a. \emph{V}). Objecto \emph{proprio} é o que se conhece \emph{primo}, antes de todos os outros objectos, de modo que a faculdade deve attingil-o \emph{imxmediatamente}, isto é, sem o intermedio de outro objecto, --- e \emph{per se}, isto é, que se conhece pela especie \emph{propria} e não por uma especie \emph{analoga}. --- Objecto \emph{primario} de uma faculdade é o que corresponde á natureza e á condição d'essa faculdade, e é a razão pela qual são attingidos os outros objectos; distingue-se do objecto \emph{secundario}, que se não conhece directamente por si mesmo, mas por outro objecto anteriormente percebido. Por isso a essencia \emph{abstracta} das coisas materiaes forma o objecto \emph{primario} da nossa intelligencia no estado actual; emquanto que todos os mais objectos, que são conhecidos por uma especie ou imagem propria das coisas materiaes, constituem o objecto \emph{secundario}.\par 2.3 Quando sustentamos que o objecto \emph{proporcionado} da intelligencia é a \emph{essencia abstracta} das coisas materiaes, não entendemos que esta nossa faculdade perceba \emph{immediata} e \emph{perfeitamente} essa essencia; porque tal modo de perceber seria contrario ao modo de operar da intelligencia humana, a\pgn{156}qual não chega ao perfeito conhecimento de uma coisa senão por degraus, passando do mais universal para o que o é menos (\emph{Cf}. \emph{Sum}. \emph{Theol}., p. T, q. 85, a. 8. A nossa intelligencia percebe num objecto primeiramente as notas ou caracteres mais universaes, como os de \emph{ente}, de \emph{uno}, de \emph{substancia}, de \emph{qualidade}, etc., e só em seguida chega ás notas mais particulares, e ás differencas essenciaes. --- Por i8s0, os primeiros conceitos, que a intelligencia forma pela abstracção dos phantasmas, nos quaes estão contidas as essencias das coisas, são os conceitos de \emph{ente}, de \emph{uno}, de \emph{bom}, etc., que se applicam a todas as coisas, que são ou podem ser. --- Por isso o \emph{ente}, que se diz ser o objecto \emph{primario} da nossa intelligencia, não é um ente qualquer, mas é o \emph{ente} considerado nas coisas materiaes e d'ellas abstrahido por uma força activa da alma (Cf. S. Thomaz, \emph{De Mag}., a. 1).\par 3.º) A primeira funcção da nossa intelligencia é a percepção da \emph{essencia} das coisas materiaes, mas \emph{abstrahida} das condições materiaes, e por isso o seu primeiro acto é uma separação, uma \emph{analyse}. --- A essencia das coisas creadas é, por si, indifferente tanto para o estado \emph{real}, quanto para o \emph{possivel}. Desde a eternidade, estava ella na mente de Deus, e não só era universal, mas participava da propria vida de Deus, emquanto era entendida por Elle. Pela creação, passou a essencia para o estado \emph{real}, e ficou sujeita á mutabilidade, a cireumsoripcao, e a todas as. condições proprias das subsistencias creadas e materiaes. Quando, porém, é percebida pela nossa intelligencia, ella fica restituida ao seu estado \emph{ideal}, reassumindo os caracteres de \emph{universalidade} e de \emph{necessidade}.}.

%%>>> t115

\pg{151} \pg{153} \tese{115}{A intelligencia é uma faculdade espiritual.}

á) Prova-se com o seguinte argumento, deduzido da \emph{natureza} da alma. --- A alma humana, oceupando o grau supremo entre as substancias sensitivas e o infimo entre as intellectuaes, deve possuir as faculdades, que são proprias das substancias sensitivas, e as faculdades, que são proprias das substancias intellectuaes e que são a \emph{intelligencia} e a \emph{vontade}. Ora as faculdades das substancias intellectuaes são \emph{espirituaes}, isto é, exercem-se sem orgãos corporeos; pois essas substancias não têm corpo. Logo a ?\emph{ntelligencia} é uma faculdade \emph{espiritual}.

b) Prova-se com argumentos, deduzidos da \emph{qualidade} dos objectos que a intelligencia percebe, e do \emph{modo} por que opéra

\emph{a}) Uma faculdade é \emph{espiritual}, se é independente do orgão material. Ora a intelligencia é independente do orgão material. Por quanto, ella percebe objectos meramente \emph{espirituaes} ou \emph{immateriaes}, --- \textsc{Deus}, a \emph{virtude}, \emph{á justiça}, a \emph{ordem}, etc. Mas, se a intelligencia fosse uma faculdade dependente da materia, não poderia perceber semelhantes objectos; porque uma faculdade dependente da materia só póde perceber objectos materiaes, devendo haver proporção entre a faculdade cognitiva e o objecto conhecido. Logo a intelligencia é uma faculdade \emph{espiritual}.

6) A intelligencia percebe a \emph{essencia} dos objectos \emph{materiaes}, --- não singular, concreta e propria de um individuo, --- mas universal, abstracta e commum a todos os individuos da mesma especie. Ora, se ella fosse uma faculdade dependente da materia, não poderia perceber a \emph{essencia} dos objectos \emph{materiaes} \pg{155}  de um modo universal e abstracto; porque uma faculdade organica só attinge os objectos que nella fazem impressão, e que são os materiaes. Logo a intelligzencia humana é uma faculdade \emph{espiritual}.

\emph{c}) A intelligencia não só percebe os objectos \emph{immateriaes}, e os objectos \emph{materiaes} de um modo \emph{immaterial}, mas tambem reflecte sobre si mesma, de modo que ella é, ao mesmo tempo, faculdade cognitiva e objecto conhecido. Ora, se esta faculdade fosse dependente da materia, não poderia reflectir sobre si mesma; porque uma faculdade organica é extensa e tem partes, e se uma parte pode reflectir sobre outra parte, não \pg{156}  pode reflectir sobre si mesma, nem o todo sobre o todo. Logo a intelligencia é uma faculdade \emph{espiritual} \footnote{Esta these, como se vê, demonstra-se --- tanto \emph{a-priori}, isto é, com argumentos tirados da propria natureza da alma, --- como \emph{a-posteriori}, isto é, com argumentos tirados das operações da mesma intelligencia. --- Os argumentos \emph{a-posteriori} foram expostos para provar a \emph{espiritualidade} da alma.\par O que deixamos dito no texto basta para mostrar que a nossa :\emph{ntelligencia} é uma faculdade \emph{essencialmente} diversa das faculdades \emph{sensitivas}, embora a intelligencia e a faculdade perceptivo-sensitiva convenham em que ambas \emph{conhecem} e \emph{passam da potencia para o acto} (S. Thom., \emph{De Anima}, lib. III, lect. 7). --- Todavia, como \emph{Rosmini} e alguns outros escriptores, pelo facto de attribuirem a sensação exclusivamente á alma e não ao composto, estabelecem que os sentidos são faculdades inorganieas, e por isso não reconhecem differenca alguma intrinseca entre a intelligencia e os sentidos, julgamos conveniente confirmar a nossa these com os argumentos seguintes:\par 4) Duas faculdades, que sio dotadas de caracteres diversos e oppostos, differem essencialmente entre si. Ora a intelligencia e os sentidos são faculdades, dotadas de caracteres oppostos. Por quanto, --- 1.º) pelos sentidos só conhecemos o \emph{material} e o \emph{singular}, pela intelligencia conhecemos o \emph{Znmnaterial}\pgn{157} e o \emph{universal}. Diz S. Thomaz: « Sensus non est cognoscitivus nisi \emph{singularium},... et non se extendit nisi ad \emph{corporalia}; intellectus autem cognoscit \emph{wniversalia},... et \emph{incorporalia}, sicut. \emph{sapientiam}, \emph{veritatem} (\emph{C}. \emph{Gentes}, 1. II, c. 66). --- 9.9) Nenhum sentido percebe o seu conhecimento, mas a in-\par ) 1 telligeneia percebe que percebe. Diz S. Thomaz: « Nullus sensus seipsum cognoscit, ac suam operationem, visus enim non videt seipsum, nec videt se videre; intellectus autem cognoscit se intelligere » (l. c.). --- 3.º) O sentido corrompe-se pela intensidade ou excellencia do seu objecto, mas a intelligencia torna-se tanto mais forte quanto mais elevado é o seu objecto. Diz S. Thomaz: Sensitivum patitur a sensibili cum corporis immutatione; unde excellentia sensibilium corrumpit sensum. Quod in intellectu non contingit, Nam intellectus intelligens maxima intelligibilium magis potest postmodum intelligere minora.Si vero in intelligendo fatigetur corpus, hoc est \emph{per accidens}, in quantum intellectus indiget operatione virium sensitivarum, per quas ei phantasmata praeparentur » (\emph{Sum}. \emph{Theol}. p. T, q. 5, a. 8 ad 2). Logo a intelligencia differe \emph{essencialmente} dos sentidos.\par )) Uma faculdade \emph{inorganica} differe \emph{essencialmente} da faculdade \emph{organica}; porque a \emph{intrinseca} independencia ou dependencia do orgão corporeo refere-se ao \emph{proprio ser} do sujeito, e por isso modifica-o na sua \emph{essencia}. Ora, como provámos, a intelligencia é uma faculdade \emph{inorganica}, porque não precisa de orgão para entender, ao passo que o sentido a uma faculdade \emph{organica}, porque precisa do orgão para sentir. Diz S. Thomaz: « Haec est differentia, qua differt cognitio intellectiva a sensitiva, quod sentire est aliquid corporeum, non enim operatio sensus est sine organo corporali, intelligere autem non est aliquid corporeum, quia operatio intellectus non est per organum corporeum » (\emph{De anima}, lib. 3, lect. 4). Logo a intelligencia differe \emph{essencialmente} dos sentidos.\par Não expomos aqui nem refutamos as objecções, com que os materialistas combatem a \emph{espiritualidade} da inteligencia; porque essas objecções foram já expostas e refutadas nesta parte da Philosophia (Sec. I, Cap. II, Art. IV).}.

\emph{A}

%%>>> t116

\pg{154} \tese{116}{A nossa intelligencia, embora espiritual, depende das faculdades sensitivas.}

a) A nossa intelligencia, embora \emph{espiritual}, está naturalmente unida com as faculdades \emph{sensitivas}. Ora esta união exige que as faculdades sensitivas sejam de algum modo uteis á intelligencia; --- não só porque, aliás, a união da alma com \pg{157}  o corpo não teria razão de ser, --- mas tambem porque, quando dois entes estão unidos, o menos nobre deve ser util e servir ao mais nobre. Mas as faculdades sensitivas, não podem ser uteis á intelligencia emquanto á \emph{existencia}; pois a intelligencia, para a sua existencia, não precisa das faculdades sensitivas. Devem, pois, ser-lhe uteis emquanto á \emph{operação}, concorrendo, na medida e segundo a natureza da sua actividade, nas operações da intelligencia. Logo a nossa intelli- . gencia depende das faculdades \emph{sensitivas} \footnote{O modo, por que um ente opéra, corresponde sempre ao modo por que existe. Os Anjos, por serem puros espiritos, não precisam, no exercicio das suas operações intellectuaes, das faculdades sensitivas. O homem, por ser composto de espirito e de materia, precisa, nas suas operações intellectuaes, do concurso das faculdades sensitivas.\par O que deixamos dicto refere-se ao estado \emph{presente} da alma; pois, quando a alma se separar do corpo, o seu modo de entender será diverso do modo por que presentemente entende, como diverso será o seu modo de existir. --- Mais tarde explicaremos o modo, por que entende a alma, separada do corpo.}.

b) O objecto da intelligencia é a \emph{essencia} dos seres, tanto \emph{immateriaes} como \emph{materiaes}. Ora a intelligencia precisa do concurso das faculdades sensitivas na percepção tanto da \emph{essencia} dos seres \emph{materiaes}, como da dos seres meramente \emph{Zmmateriaes}. --- Precisa na percepção da \emph{essencia} dos seres \emph{materiaes}; pois esta, sendo em si singular e concreta, só póde impressionar uma faculdade material ou sensitiva. --- Precisa na percepção da \emph{essencia} dos seres \emph{immateriaes}; pois que a intelligencia não póde elevar-se ao espiritual ou immaterial, senão por meio do sensivel ou material, e o sensivel ou material só é attingido por uma faculdade sensitiva. Logo a intelligencia depende das faculdades \emph{sensitivas} \footnote{Poderia objectar-se: É verdade que a intelligeúcia, na percepção da \emph{essencia} dos seres \emph{immateriaes}, precisa do concurso das faculdades sensitivas, mas é falso que precise d'ellas na percepção da \emph{essencia} das coisas \emph{materiaes}; pois a \emph{essencia} das coisas \emph{materiaães}, por ser universal e abstracta, só póde ser percebida pela intelligencia, e não pelas faculdades sensitivas.\par A resposta está no texto. A \emph{essencia} das coisas \emph{materiaes}, considerada na propria realidade, é singular e concreta; porque está cercada de propriedades materiaes ou individuaes. Sendo singular e concreta, a essencia só póde ser percebida por uma faculdade organica ou sensitiva. --- Se a \emph{essencia} das coisas \emph{materiaes} chega a tornar-se \emph{universal} e \emph{abstracta}, este phenomeno opéra-se na alma depois da percepção dos sentidos externos, e é devido a uma \emph{força abstractiva}, de que a mesma alma é dotada.}.

c) A experiencia attesta --- que os homens, destituidos de um sentido externo, não possuem a sciencia dos objectos relativos ao mesmo sentido; --- que as operações da intelligencia estão subordinadas ás vicissitudes das faculdades sensitivas, de modo que, se se perturbar ou cessar o exercicio das faculdades sensitivas, perturba-se ao mesmo tempo ou cessa o exercicio da intelligencia; --- que, quando a intelligencia pensa, o cerebro trabalha e se cança; --- que todo o objecto, que a intelligencia entende, quer \emph{immaterial} quer \emph{material}, está sempre cercado de uma imagem sensivel, ao menos da imagem sensivel da \emph{palavra}. Ora estes factos seriam inexplicaveis, se a nossa intelligencia, embora espiritual, não dependesse das faculdades \emph{sensitivas}. Logo \footnote{Desenvolvamos os factos apontados no texto.\par 1.9) Os homens, destituidos de um certo sentido, são tambem destituidos das ideas relativas ao objecto proprio do mesmo sentido. Este facto que prova? Prova que nada chega á intelligencia, se não tenha primeiramente passado pelos sentidos, e por isso que a intelligencia depende, no seu exercicio, do exercicio das faculdades sensitivas.\par 2.9) As operações da intelligencia estão subordinadas ás vicissitudes das\pgn{159}faculdades sensitivas. Este phenomeno, que está provado até á evidencia, póde indicar duas coisas: --- ou que a intelligencia é uma faculdade sensitiva, --- ou que a intelligencia precisa, no seu exercicio, do concurso das faculdades sensitivas. Ora a intelligencia não é uma faculdade sensitiva, mas espiritual. Logo deve dizer-se que a intelligencia, no seu exercicio, precisa do concurso das faculdades sensitivas.\par 3.º) Está averiguado que o cerehro experimenta fadiga ou cançaço, quando pensamos. Esta fadiga ou \emph{cançaço} --- não pôde indicar que é o cerebro que pensa, pois repugna que o cerebro pense, --- mas indica que o exercicio da intelligencia é acompanhado do exercicio das faculdades sensitivas, que têm o seu \emph{orgão} no cerebro.\par 4.9) Todo o objecto, que a intelligencia percebe, está como que envolvido numa imagem sensivel, \emph{e} sem esta imagem sensivel nada. se percebe. Podemos considerar os objectos mais immateriaes e mais abstractos, e todavia sempre os veremos cercados de condições ou propriedades individuaes e revestidos de uma forma sensivel, ao menos da forma sensivel da \emph{palavra}. (Cf, \emph{Sum}. \emph{Theol}. p. Y, q. 84, a. 7).}.

As faculdades \emph{sensitivas}, de que a intelligencia depende, são principalmente a \emph{imaginação} e os \emph{sentidos externos}; pois que o phantasma ou imagem sensivel, que envolve o objecto da intelligencia, é formado pela \emph{imaginação}, e o exercicio d'esta faculdade suppõe o exercicio dos \emph{sentidos externos} \footnote{As faculdades \emph{sensitivas}, de que a intelligencia depende no seu \textbf{exercicio}, são principalmente a \emph{imaginação} e os \emph{sentidos externos}.\par \emph{a}) A intelligencia depende da \emph{imaginação}. Com elffeito, \emph{a} imagem sensivel, em. que resplandece o ideal que a intelligencia attinge, é um producto da \emph{imaginação}. --- De mais, a intelligencia, com as suas operações, percebe objectos que \emph{actualmente} não são percebidos pelos sentidos externos, mas que todavia estão presentes á alma. Ora objectos, que não são actualmente percebidos pelos sentidos externos, sô podem estar presentes á alma pelo facto de residirem, por meio de uma especie sensivel, na \emph{imaginação}.\par \emph{b}) A intelligeneia depende dos \emph{sentidos externos}. Por quanto, como provámos, o exercicio da \emph{imaginação} suppõe o exercicio dos sentidos \emph{externos}, e um cego de nascimento nunca poderá ter a imagem sensivel das córes.\par D'onde se vê que a intelligencia depende \emph{immediatamente} da \emph{imaginação}, e \emph{mediatamente} dos sentidos \emph{externos}.\par Cousulte-se o que dissemos ácerca d'este assumpto na refutação do \emph{materialismo} (\emph{Antrop}., Sec. I, Cap. II, Art. IV --- \emph{nota}).\par Só queremos fazer uma advertencia. Como o exercicio da intelligencia depende do exercicio das faculdades sensitivas, e sobretudo da \emph{imaginação}, \emph{ratione objecti}, isto é, emquanto os sentidos subministram o material do pensamento, é claro que deve existir uma certa correlação entre o desenvolvimento do cerebro, e o desenvolvimento da intelligencia, de modo que, embora a intelligencia seja uma faculdade inorganica e não haja entre ella e o cerebro uma proporção ou equação rigorosa, todavia pode dizer-se que a perfeição do cerebro é, \emph{geralmente}, um signal, ainda que \emph{indirecto}, e uma medida, ainda que \emph{approawimativa}, da perfeição da intelligencia. E nada mais natural. E tão intrinseca a união estabelecida por Deus entre a materia e a\pgn{160}força, entre o corpo e a alma, que o corpo pode ser considerado como o espelho e a expressão da alma. Diz o S. Doutor: « Quamvis intellectus sit absolutus a corpore quantum ad propriam operationem, quae corporali organo non expletur, tam»n coniunctionem habet ad corpus dupliciter: scilicet ex parte essentiae animae, quae corporis forma est, et ex parte inferiorum potentiarum, ex quibus intellectus recepit; \emph{ef per istum}. \emph{modum}. \emph{diversitas}: \emph{corporis in diversitatem}. \emph{intellectus redundat} » (In 2 \emph{sent}., dist. 32, q. 2, a. 8). --- Noutro logar, S. Thomaz diz que o cerebro do homem é, \emph{relativamente}, o mais desenvolvido, por isso mesmo que é o orgão das faculdades sensitivas que concorrem, a seu modo, ao acto do entendimento. Eis as suas palavras: \& Necessarium fuit quod homo inter \emph{omnia} animalia \emph{respectu sui corporis} haberet \emph{maximum} cerebrum, ut liberius in eo perficerentur operationes inferiorum virium sensitivarum, quae sunt necessariae ad intellectus operationem » (\emph{Sum}. \emph{Theol}, p. 1, q. 91, a. 8 ad 1).}.

117. A \textbf{nossa} intelligencia não depende das \textbf{faculdades sensitivas} intrinseca, mas \textbf{extrinsecamente},

a) Não depende \emph{intrinsecamente}. Uma faculdade depende \emph{intrinsecamente} de outra, --- ou quando aquella deriva d'esta, --- ou quando uma é dirigida pela outra. Ora a nossa intelligencia --- não deriva das faculdades sensitivas, pois uma faculdade \emph{immaterial}, como é a intelligencia, não póde derivar de faculdades \emph{materiaes}, como são as sensitivas, --- nem é dirigida pelas faculdades sensitivas, pois é aquella que dirige estas. Logo a nossa intelligencia não depende \emph{intrinsecamente} das faculdades \emph{sensitivas}. ;

b) Mas depende \emph{extrinsecamente}. Uma faculdade depende \emph{extrinsecamente} de outra, quando aquella não póde attingir o seu objecto, se esta lh'o não apresenta. Ora a nossa intelligencia não póde perceber o seu objecto, se as faculdades sensitivas lh'o não apresentarem; pois, como dissemos, o objecto da intelligencia não chega a esta faculdade, senão passando pelas faculdades \emph{sensitivas} e envolvido numa imagem sensivel. Logo a nossa intelligencia depende \emph{extrinsecamente} das faculdades \emph{sensitivas} \footnote{Esta conclusão deduz-se da antecedente; pois pelos argumentos, que adduzimos para provar que á intelligencia depende das faculdades sensitivas, via-se que essa dependencia se basóa no facto de a intelligencia não perceber objecto algum, que não tenha sido anteriormente percebido pelos \emph{sentidos externos}, e que não seja acompanhado por uma imagem sensivel, formada pela \emph{imaginação}. O facto, pois, de a intelligencia depender das faculdades sensitivas, não demonstra a superioridade d'estas sobre aquella, mas a superioridade da intelligencia sobre as faculdades sensitivas.\par Aclaremos a doutrina exposta com uma comparação. O inferior e o superior dependem mutuamente um de outro. Mas ha differenca entre estas dependencias. O inferior depende do superior, emquanto aquelle obedece a\pgn{161}este; ao passo que o superior depende do inferior, não emquanto aquelle. obedece a este, mas emquanto o superior precisa do servico do inferior, --- serviço a que o superior tem direito e que o inferior não póde deixar de prestar. É o que acontece no nosso caso. A intelligencia depende das faculdades sensitivas, como o superior depende do inferior; pois as faculdades sensitivas apresentam o objecto, que a intelligencia, e só ella, percebe.\par Poderia objectar-se: Se a intelligencia percebesse o objecto, que as faculdades sensitivas lhe apresentam, ella só poderia perceber as qualidades materiaes das coisas e não a essencia; pois pelos sentidos só percebemos as qualidades materiaes, e não a essencia.\par Resposta. E' verdade que pelas faculdades sensitivas só percebemos as qualidades das coisas; mas estas qualidades são percebidas pelos sentidos, não de um modo abstracto e universal, mas de um modo concreto e singular, isto é, emquanto \emph{unidas e adherentes a um certo e determinado sujeito}. Sendo assim, é claro que o objecto, que as faculdades sensitivas apresentam á intelligencia, representa não só as qualidades sensiveis de um ser, mas tambem a sua essencia, embora singular e envolvida nas mesmas qualidades sensiveis. Em seguida, a essencia é abstrahida das qualidades sensiveis, tornando-se assim universal e abstracta.}.

%%>>> t118

\pg{158} \pg{159} \pg{160} \pg{161} \tese{118}{A intelligencia humana é uma faculdade potencial.}
Uma faculdade é e diz-se \emph{potencial}, quando, embora tenha sempre a força de produzir os seus actos, todavia não os está produzindo sempre, ou não os produz todos ao mesmo tempo, mas successivamente. Ora a nossa intelligencia, posto que tenha sempre a força de entender, não está comtudo entendendo sempre actualmente, nem entende, ao mesmo tempo, tudo o que póde entender. Logo a intelligencia humana é uma faculdade \emph{potencial}. --- E' neste sentido que a nossa intelligencia foi tambem chamada \emph{possivel} \footnote{Quando, pois, os philosophos dizem que a nossa intelligencia é \emph{potencial} ou \emph{possivel}, esta \emph{potencialidade} ou \emph{possibilidade} não se refere \& \emph{existencia} da propria intelligencia, mas sim á sua \emph{operação}. Por outras palavras: A nossa intelligencia chama-se \emph{potencial} ou \emph{possivel}, --- não porque não exista realmente, --- mas porque está em \emph{potencia} com relação ás suas percepções e da potencia passa para o acto. É a doutrina dos escholasticos. Para S. Thomaz, a intelligencia \emph{possivel} não é outra coisa senão o principio intellectivo, emquanto se acha no estado de \emph{potencia} relativamente \& propria intellecção. Eis as suas palavras: « Cúm inveniamur quandoque in \emph{potentia}, necesse est ponere aliquam virtutem, per quam simus intelligentes in \emph{potentia}. Et propter ' hoe vocatur \emph{intellectus possibilis}; sicut et sensus, secundum quod est in \emph{potentia}, posset vocari sensus \emph{possibilis} » (\emph{De spirit}. \emph{creat}., a. 9). E como é a mesma faculdade que está em \emph{potencia} com relação ao seu acto e depois o produz, é claro que a intellecção, considerada \emph{formalmente}, é o acto proprio da inteligencia \emph{possivel}. Diz S. Thomaz: « Intellectus \emph{possibilis} est, quo hic homo \emph{formaliter} intelligit » (\emph{In} 3 \emph{de} Anima, lect. 7).\par 1i Vor. IL}.

%%>>> t119

\tese{119}{A intelligencia humana deve ser determinada pela especie intelligivel.}
Sendo faculdade \emph{potencial}, a intelligencia não \pg{162}  póde produzir o seu acto ou perceber um certo objecto, sem que este se lhe una e assim a determine á percepção; pois a percepção faz-se, emquanto o objecto percebido se une á faculdade perceptiva. Ora um objecto não póde unir-se á intelligencia na sua realidade. Logo deve unir-se-lhe por meio de uma imagem, que represente a sua essencia. Essa imagem deve ser \emph{espiritual} ou \emph{immaterial}, --- não só porque representa a \emph{essencia} de objectos, tanto \emph{Zmmateriaes} como \emph{materiaes}, abstrahida das condições ou propriedades individuaes, e por isso \emph{universal}, --- mas tambem porque uma imagem \emph{material} n\&o poderia unir-se á intelligencia, que é espiritual, pois a forma deve ser proporcionada ao sujeito, que á recebe. A imagem \emph{immaterial}, que representa a \emph{essencia} de objectos, tanto \emph{immateriaes} como \emph{materiaes}, abst :'ahida das condições ou propriedades concretas ou individuaes, e por isso \emph{universal}, diz-se \emph{especie intelligivel}. Logo a intelligencia humana deve ser determinada pela \emph{especie intelligivel} \footnote{Como se vê, a necessidade das imagens representativas das essencias basóa-se na propria natureza da nossa intelligencia. Por quanto esta, sendo dotada de \emph{potencia} para entender tudo, mas não entendendo tudo \emph{actualmente}, é por si indifferente para entender este ou aquelle objecto. Mas o ente, que esta em \emph{potencia}, não póde passar ao \emph{acto}, se não pela acção de outro ente que já exista em \emph{acto}. Ora o ente em \emph{acto}, que, pela sua acção, faz passar a intelligencia da \emph{potencia} ao \emph{acto} de entender, só póde ser o objecto percebido; aliás seria inexplicavel o facto de a intelligencia perceber um certo e determinado objecto, \emph{e} não um outro da mesma ou de diversa especie. Mas a acção, que o objecto exerce sobre a intelligencia, consiste em imprimir nesta faculdade uma imagem, que o represente e que, unida á mesma intelligencia, a determine ao acto. Na verdade, todo o conhecimento faz-se pela uni\&o do objecto percebido com a faculdade perceptiva; pois o conhecimento é uma acção \emph{immanente}, que tem o seu termo no proprio sujeito, visto que ella não aperfeigóa o objecto que é conhecido, mas sim o sujeito que conhece. Ora o objecto, como dizemos no texto, não póde unir-se na sua physica realidade com a faculdade; aliás identificar-se-hia com a propria faculdade, o que, além de ser impossivel, levaria ao pantheismo. Logo deve unir-se-lhe por meio de uma imagem, que o represente. Esta imagem diz-se \emph{especie}.\par A \emph{especie} póde ser \emph{sensivel} e \emph{intelligivel}. --- E \emph{sensivel}, quando representa o objecto cercado de \emph{condições materiaes}; e esta especie une-se com a faculdade \emph{sensitiva}. --- É \emph{intelligivel}, quando representa a \emph{essencia} dos objectos \emph{materiaes}, abstrahida de todas as propriedades da materia; e só esta especie póde unir-se com a intelligencia. --- A faculdade perceptiva, quer sensitiva quer intellectual, informada pela \emph{especie}, produz o seu acto, e percebe \emph{directamente} o objecto representado pela mesma especie. Por isso a \emph{especie} não é o objecto, que é percebido, mas é a forma, \emph{por meio da qual} o objecto é percebido. '\par \emph{A} especie \emph{sensivel} e a \emph{intelligivel} differem entre si. Com effeito, a especie sensivel, se é \emph{immaterial}, emquanto á sua \emph{natureza}, pois que não é composta\pgn{163}de materia, não é immaterial-emquanto ao \emph{objecto} que representa, pois esse objecto está cercado de condições materiaes. Pelo contrario, a especie \emph{intelligivel é immaterial}, não só emquanto á sua \emph{natureza}, porque não é composta de materia, mas tambem emquanto ao seu \emph{objecto}, pois este abstrahe das condições ou propriedades materiaes.\par Agora podemos entender o sentido do celebre axioma da eschola: a \emph{intelligencia}, \emph{entendendo}, \emph{torna-se todas}. \emph{as coisas}. Por quanto, a intelligencia conhece um certo objecto, emquanto nella existe a forma ou especie do mesmo objecto. Ora a forma dà o proprio \emph{ser} ao sujeito, de que á forma; assim a \emph{humanidade} dá ao sujeito, que a recebe, o \emph{ser humano}. Logo a intelligencia, quando \emph{actualmente} entende, torna-se a propria coisa, que ella entende. Mas, póde entender todas as coisas. Logo póde tornar-se todas as coisas. --- Advirta-se comtudo que a intelligencia não se torna todas as coisas por um modo \emph{real}, mas por um modo \emph{intencional}; emquanto que, pelo facto de receber a imagem dos objectos, se torna semelhante a estes, embora conserve sempre a propria natureza. Por isso os escholasticos, com aquelle axioma, não querem sustentar a doutrina que a nossa alma é composta dos elementos das coisas materiaes, que ella conhece, --- doutrina que ensinam os \emph{materialistas}; nem querem affirmar que a intelligencia se identifica \emph{realmente} com os objectos \emph{entendidos}, --- afirmação que defendem os \emph{pantheistas}; mas só querem indicar que a percepção intellectual se effectua pela união da \emph{especie intelligivel} com a intelligencia, e que, por essa união, a intelligencia não se identifica com a \emph{realidade} da coisa, mas só com a sua \emph{especie intelligivel} (Cf. \emph{C}. \emph{Gentes}, 1. II, c. 98). :\par Pela definição da faculdade \emph{passiva}, e pelo que acabamos de expor ácerca da necessidade da especie \emph{intelligivel}, deve concluir-se que a nossa intelligencia é uma faculdade \emph{passiva}. O Angelico Doutor à questão --- \emph{utrum}. \emph{éntellectus sit potentia passiva} (\emph{Sum}. \emph{Theol}., P. I, q. 79, a. 2), responde que, \emph{sicut intelligere nostrum}. \emph{est pati}, a intelligencia é faculdade \emph{passiva}. Distingue 8. Thomaz tres especies de \emph{paivao}. Dá-se a primeira, quando o sujeito se muda, perdendo alguma coisa, que convém á sua natureza ou inclinação; como quando o homem adoece. Dá-se a segunda, quando o sujeito se muda, perdendo alguma coisa, que lhe convém ou pode deixar de lhe convir; e neste sentido dizemos que é \emph{passivo} não só o homem que adoece, mas tambem o que deixa de estar doente. Dá-se a terceira, quando o sujeito se muda passando da \emph{potencia} para o \emph{acto}, e diz-se \emph{passivo}, embora não perca nenhuma propriedade, mas, antes, se torne mais perfeito. É esta terceira especie de\pgn{164}\emph{passividade} que convém á intelligencia humana, \& qual, embora seja \emph{activa} com relação ao acto de entender, é \emph{passiva} com relação ao seu objecto e á\par \emph{especie intelligivel} do objecto.\par Esta theoria serve para refutar os erros dos \emph{idealistas} e dos \emph{sensualistas}.\par \emph{a}) Os \emph{idealistas} opinam que \& intelligencia é absolutamente \emph{activa} e não soffre o influxo dos objectos externos. --- Esta opinião é falsa; porque é contraria á experiencia e nega a \emph{realidade} ou \emph{objectividade} das nossas percepções, quo consiste na conformidade d'estas com os seus objectos.\par \emph{b}) Os \emph{sensualistas} afirmam que a intelligencia é completamente \emph{passiva} nas suas percepções, limitando-se a receber a impressão do objecto sensivel, --- impressão em que ella depois reflecte com a sua attenção. --- Tambem esta affirmação é falsa; porque todo o conhecimento exige um \emph{acto} do sujeito que conhece, e por isso não póde ser e não é mera \emph{passividade}.}.

%%>>> t120

\tese{120}{A intelligencia humana, no seu exercicio, precisa de uma força abstractiva.}
Esta proposição é um corollario das proposições antecedentes. --- O objecto da nossa intelligencia, que é a \emph{essencia} das coisas, immateriaes ou materiaes, chega á nossa alma envolvido num phantasma ou imagem sensivel, e por isso cercado de propriedades individuaes; pois o phantasma contém a essencia das coisas, como ella existe na realidade ou na imaginação, isto é, cercada de propriedades materiaes. Mas o phantasma, sendo material e concreto, não póde, só por si, operar na intelligencia; pois nenhuma coisa material \pg{163} \pg{164}  póde por si fazer impressão numa faculdade immaterial. Logo é necessaria uma \emph{força} que, operando no phantasma, \emph{abstrata} a essencia das propriedades materiaes, e assim o proprio phantasma, em virtude do \emph{universal}, que ele encerra e que é \& \emph{essencia} abstrahida das propriedades materiaes, se torne capaz de operar directamente na intelligencia e produzir nella a \emph{especie intelligivel}. --- Esta \emph{forga abstractiva}, que torna o phantasma capaz de operar directamente na intelligencia e produzir nella as \emph{especies intelligiveis}, chama-se \emph{inteligencia agente} \footnote{Diremos alguma coisa ácerca da \emph{intelligencia agente}, indicando os seus varios \emph{nomes}, e tratando do seu \emph{acto}, da sua \emph{natureza}, e da sua \emph{existencia}; pois o assumpto é um dos mais importantes e mais difficeis da philosophia, o é a base de toda a \emph{ideologia}. --- S. Thomaz trata d'esta questão --- na \emph{Sum}. \emph{Theol}., p. I, q. 7, aa. 8, 4, 5; --- \emph{De anima}, a. 5; --- \emph{C}. \emph{Gentes}, l. II, cc. 76, 77, 18.\par a) \textsc{VARIOS} NOMES DA INTELLIGENCIA AGENTE. --- A \emph{intelligencia agente} tem varios nomes. --- Na verdade:\par a) Chama-se \emph{intelligencia} ou \emph{entendimento agente}. --- Diz-se \emph{intelligencia} ou \emph{entendimento}, não porque entende, mas porque torna \emph{actualmente tintelligivel} um objecto. Diz-se \emph{agente}, porque pela sua \emph{acção} faz com que o objecto, que por si era \emph{potencialmente} intelligivel, o seja \emph{actualmente}.\par \emph{b}) Ohama-se tambem \emph{força abstractiva}; porque, para a producção da \emph{especie intelligivel}, se serve do phantasma \emph{emquanto á essencia} que está nelle contida, e não emquanto ás \emph{condições materiaes}, que cercam a, mesma essencia; de modo que o phantasma só imprime na \emph{intelligencia possivel} a \emph{essencia} das coisas, e não as \emph{condições anateriaes}.\par c) Chama-se finalmente \emph{luz intellectual}; por quanto; assim como a luz \emph{material}, illuminando os objectos, os torna visiveis aos olhos corporeos; assim \& \emph{força abstractiva}, illuminando as \emph{imagens sensiveis}, torna visivel á intelligencia a \emph{essencia} das coisas, que estava cercada de propriedades materiaes, concretas e singulares.\par (p) Acto DA INTELLIGENCIA AGENTE. --- Com relação ao \emph{acto} da \emph{inteligencia agente} advertimos:\par 1.9) Ha uma duplice especie de \emph{abstracgdo}: uma abstracção de \emph{simplicidade}, e uma outra de \emph{composição} e \emph{divisão}. Faz-se a abstracção de \emph{simplici}- .\pgn{165}\emph{dade}, quando consideramos uma coisa sem outra, como quando consideramos a essencia sem os accidentes. Faz-se a abstracção de \emph{composição} ou \emph{divisão}, quando percebemos que uma coisa está unida á outra, ou separada d'ella. Ora, a abstracção, feita pela intelligencia \emph{agente}, não é a de \emph{composição} ou \emph{divisão}, mas sim a de \emph{simplicidade}. --- Esta abstracção de \emph{simplicidade} não exige o conhecimento e a comparação dos termos, como falsamente ensinou \emph{Rosmini} (\emph{N}. \emph{Saggio}, Sez. X, c. 1), mas exige apenas a existencia do \emph{phantasma}, de que a intelligencia \emph{agente} extrahe a essencia das coisas, separando-a das condições materiaes (Cf. \emph{Sum}. \emph{Theol}, p. L, q. 85, a. 1 ad 1).\par 9,9) As \emph{especies intelligiveis}, que representam a essencia dos seres, abstrahida das condições materiaes, e que têm por \emph{causa principal á intelligencia agente}, não podem ser falsas. Com effeito, póde haver falsidade na abstraccáo de \emph{composição} e de \emph{divisão}, emquanto attribuimos a um sujeito uma propriedade que lhe não pertence, ou lhe negamos uma propriedade que lhe pertence; mas não póde haver falsidade na abstracção de \emph{simplicidade}, pois que a essencia das coisas, embora esteja cercada e como que envolvida pelas condições materiaes, todavia é realmente distincta d'ellas, e nada. obsta a que se considere aquella sem estas.\par 8.9) As especies, que o phantasma, auxiliado pela \emph{intelligencia agente}, produz na \emph{intelligencia possivel}, representam as essencias dos entes, não subjectivas e imaginarias, mas \emph{objectivas} e \emph{reaes}; pois aquellas essencias são as mesmas que se encontram nos entes. Só ha esta differenca: nos proprios entes, as essencias sio singulares e concretas; ao passo que, na intelligencia, são universaes e abstractas. Assim a \emph{humanidade}, que existe na intelligencia, é a mesma essencia que existe em cada homem, embora seja diverso o seu modo de existir na intelligencia e no individuo humano.\par 4.9) A \emph{especie intelligivel} é um producto da \emph{imagem sensivel} ou \emph{phantasma} e da \emph{intelligencia agente}; e por isso participa das condições d'estes dois factores. --- Participa da \emph{imagem sensivel}; pois que a \emph{essencia}, que é representada pela \emph{especie intelligivel}, \emph{é objectiva}, como objectiva é a imagem sensivel, em que a mesma essencia está contida. --- Participa da \emph{intelligencia agente}; pois que a \emph{essencia}, abstrahida, por virtude da \emph{intelligencia agente}, das condições znateriaes, de que está cercada na realidade e na propria imagem sensivel, torna-se \emph{universal} ou 4\emph{mmaterial}, como immaterial é aquella faculdade.\par Pergunta-se: De que modo a \emph{ntelligencia agente} e o \emph{phantasma} concorrem para a producção da \emph{especie intelligivel}?\par Resposta. Na producção da \emph{especie intelligivel}, a \emph{intelligencia agente} concorre como causa \emph{principal} e o \emph{phantasma} como causa \emph{instrumental}. Diz 8. Thomaz: « Na formação das especies intelligiveis, que a intelligencia \emph{possivel} recebe dos phantasmas, estes concorrem como causa \emph{tzstrumental} e \emph{secundaria}, e a inteligencia agente como causa \emph{principal} e \emph{primaria}. É por isso que o effeito produzido na intelligencia \emph{possivel} se conforma com as condições de ambas estas causas, e não de uma só; e a intelligencia \emph{possivel} recebe as especies como \emph{actualmente intelligiveis}, pelo influxo da intelligencia \emph{agente} (que é \emph{espiritual}), e recebe-as como \emph{expressões} ou \emph{semelhanças} de objectos determinados, pelo influxo dos \emph{phantasmas} (que representam coisas sensiveis) Portanto essas especies não existem actualmente na imaginação, nem\pgn{166}na intelligencia \emph{agente}, mas só na intelligencia \emph{possivel} » (\emph{De verit}., q. X, a. 6 ad 7). Expliquemos a resposta. É innegavel que a nossa intelligencia (possivel) precisa, para entender, das especies \emph{intelligiveis}, e que estas devem ser tiradas das coisas sensiveis. As coisas sensiveis devem actuar na alma, para que esta possa recolher d'ellas as especies intelligiveis. Diz ainda o nosso Angelico: « In mente, accipiente scientiam a rebus, formae existunt per quamdam \emph{actionem rerum in animam} » (\emph{l}. c. a. 6) Ora as coisas sensiveis não podem actuar na alma pela sua realidade physica, mas só por uma imagem que as represente, por um \emph{p/iantasma}. O phantasma, portanto, deve concorrer \emph{activamente} para a producção da \emph{especie intelligivel}. Ora o phantasma, sendo \emph{muterial}, isto é, cercado de condições materiaes ou sensiveis, não pode \emph{por si}, por virtude \emph{propria}, actuar ou fazer impressão na intelligencia, que é \emph{immaterial}, porque \emph{nihil corporeum imprimere potest} (\emph{per se}) \emph{in rem incorpoream} (\emph{Sum}. \emph{Theol}., p. 1, q. 84, a. 6). É necessaria, portanto, uma força que conforte o phantasma e o torne capaz de actuar na intelligencia (possivel) e assim produzir nella a especie intelligivel. Esta força é a intelligencia \emph{agente}. Sendo assim, é claro --- que o phantasma concorre, ng produecüo da especie \emph{éstelligivel}, como causa \emph{instrumental}, porque a propriedade do \emph{instrumento}, emquanto \emph{tal}, consiste em que, conservando sempre a sua natureza, pode, pela virtude de um agente superior, ser elevado, precariamente, a produzir um effeito, que por si só nunca teria podido produzir, --- e que a intelligenecia agente concorre como causa \emph{principal}. Ei o que ensina S. Thomaz: « In nostra phantasia est phantasma seu forma repraesentans /\emph{iunc} hominem secundum quod fuit in aliquo sensu exteriori. \emph{Quae forma virtute intellectus agentis agit in intellectum possibilem}, sicut colores virtute luminis agunt in potentiam visivam, et causatur tunc in intellectu possibili quaedam forma, quae dicitur \emph{species intelligibilis} » (\emph{Sum}. \emph{Log}. \emph{Arist}, \&. 1, c. 1; --- Cf. \emph{C}. \emph{Gentes}, l. II, e. 77; \emph{Quodlib}. VIII, q. 3).\par Pergunta-se ainda: Mas como é que 2 intelligencia \emph{agente} torna o phantasma capaz de actuar ou fazer impressão na intelligencia \emph{possivel}, para produzir nella a especie \emph{intelligivel} ?\par Resposta. O phantasma representa o objecto sensivel, como está em si mesmo. Ora o objecto sensivel, como está em si mesmo, contém a \emph{essencia}, pela qual elle é o que é, mas a essencia cercada de condições materiaes ou sensiveis. Logo tambem o phantasma contem \& \emph{essencia} do objecto envolvida nas condições materiaes. Ora o phantasma, assim como está, não pode actuar na intelligencia. Todavia, se houvesse uma luz, que illuminando o phantasma, pozesse a descoberto a \emph{essencia}, considerada nos seus elementos constitutivos e despida das condições materiaes, que a cercam, é evidente que o phantasma, emquanto e só emquanto contem essa \emph{essencia}, abstrahida das condições materiaes e por isso \emph{universal}, pode actuar ou fazer impressão na intelligencia, e nella produzir a especie \emph{intelligivel}, que representa á propria \emph{essencia abstracta} ou \emph{universal}. Ora esta \emph{lus intellectual} é a intelligencia \emph{agente}: u \emph{Lux} ista, \emph{qua} mens nostra intelligit, est intellectus \emph{agens} 7 (S. Thom., \emph{De spirit}. \emph{creat}., a. 10). Devemos, pois, dizer que, --- assim como a luz sensivel, illuminando um corpo, não imprime nos nossos olhos \emph{todas} as qualidades ou propriedades d'esse corpo, mas só as que são o objecto\pgn{167}proprio dos olhos, e que são as \emph{córes}, e d'este modo torna os corpos capazes de mover e determinar a faculdade visual, produzindo nella as especies visiyeis, --- assim tambem a intelligencia \emph{agente}, illuminando os phantasmas, faz 'apparecer á intelligencia \emph{possivel} a \emph{essencia abstracta} e \emph{universal}, que é o objecto proprio da mesma intelligencia \emph{possivel}, e não as condições \emph{materias}, que são o objecto proprio das faculdades sensitivas, e d'este modo eleva o phantasma e o torna capaz de actuar na intelligencia \emph{possivel}, produzindo nella as especies \emph{intelligiveis}. --- D'onde se vé que duplice é a acção da intelligencia \emph{agente} no phantasma; \emph{illumina-o} e \emph{abstrahe} d'elle a especie \emph{intelligivel}. A illuminação consiste nisto, que a intelligencia \emph{agente}, operando no phantasma, faz com que se possa abstrahir d'elle a especie intelligivel. A. \emph{abstracção} consiste nisto, que a intelligencia \emph{agente} extrahe effectivamente do phantasma, assim illuminado, a \emph{essencia}, sem se importar com as condições materiaes, e produz com o mesmo phantasma a especie intelligivel, que se imprime na intelligencia \emph{possivel}. (\emph{Of}. \emph{Sum}. \emph{Theol}., p. I, q. 85, a. 1 ad 4). A \emph{illuminacüo} precede naturalmente a \emph{abstracção}, porque só em virtude d'essa illuminacüo é que a essencia pode ser abstrahida das condições materiaes. Todavia, estes dois actos --- \emph{illwminacdo} e \emph{abstracção} --- estão por tal forma unidos, que um chama ou suppõe o outro. Por isso S. Thomaz dá ás vezes á \emph{illuminação} o nome de \emph{abstracção}, como no trecho seguinte: « Intellectus \emph{agentis} est \emph{illuminare} non quidem alium intelligentem, sed intelligibilia in potentia, in \emph{quantwm} per \emph{abstractionem} facit ea intelligibilia in actu » (\emph{Sum}. \emph{Theol}, p. 1, q. 54, a. 4 ad 2), --- como tambem chama a intelligencia \emph{agente}, indifferentemente, \emph{lus} e \emph{força abstractiva}, para significar uma e a mesma coisa. :\par Do que temos dito vê-se o absurdo da doutrina, que pretende \emph{identificar} a \emph{lus intellectual}, pela qual se descobre a \emph{essencia} das coisas, com um pretendido \emph{ente ideal}, que seria \emph{innato} na nossa intelligencia e seria o meio para o conhecimento de todas as coisas. O angelico Doutor diz: « Ipsum \emph{lumen} intellectus nostri non se habet ad intellectum nostrum sicut \emph{quod} intelligitur, sed sieut \emph{quo} intelligitur » (\emph{Sum}. \emph{Theol}., p. I, q. 88, a.8 ad 1).\par Objectam: Se a intelligencia \emph{agente} abstrahisse do phantasma a \emph{especie intelligivel}, transformaria a natureza do proprio phantasma, e o material tornar-se-hia immaterial. Ora a natureza das coisas não se transforma.\par Concedemos que a natureza das coisas não póde estar sujeita a transformações, mas negamos que, pela acção da \emph{intelligencia agente}, se transforme a natureza do phantasma. Porquanto, a abstracção, realizada pela intelligencia \emph{agente}, não é \emph{real}, mas é \emph{intencional}; visto que a mesma intelligencia \emph{agente}, operando no phantasma, não separa \emph{realmente} da essencia as condições materiaes, mas só attinge a essencia das coisas, sem se importar com as condições da materia, de que a propria essencia se encontra cercada no phaníasma. E' falso, pois, que a intelligencia \emph{agente} transforme a natureza do phantasma. --- Por isso, quando os philosophos dizem que o phantasma, pela força da \emph{intelligencia agente}, se torna \emph{immaterial}, tal expressão não se refere á \emph{nature:a} do proprio phantasma, porque esta não se transforma, mas refere-se ao \emph{termo}, que é produzido pelo phantasma e pela \emph{intelligencia agente}, e que é a \emph{especie intelligivel}; pois a \emph{especie intelligivel}, representando\pgn{168}168 METAPHYSICA --- PARTE TERCEIRA --- ANTIIROPOLOGIA verdades não conhecidas. --- Digamos alguma coisa da sua \emph{divisão} e da sua \emph{identidade real} com a intelligencia.\par \emph{gb} a \emph{essencia} das coisas e não as condições materiaes, é \emph{immaterial} (S. 'Thom., \emph{De potentiis animae}, o. VI; \emph{Sum}. \emph{Theol}, p. I, q. 85, a. 1 ad 1b \textbar{}\par Insistem: \emph{A} abstracção é uma operação propria da \emph{inteligencia possivel},\par e não da \emph{intelligencia agente}.\par Resposta. Ha duas especies de abstracção: uma de \emph{efficiencia}, outra de \emph{simples consideração}. \emph{A} abstracção de \emph{efficiencia} é operação propria da \emph{intelligencia agente}, que, operando no phantasma, abstrahe a essencia das condições materiaes. D'este modo o phantasma, em virtude do \emph{universal}, que elle encerra e que é a \emph{essencia} das coisas, abstrahida das condições materiaes, torna-se capaz de operar directamente na intelligencia, e produzir nella a \emph{especie intelligivel}. --- A abstracção de \emph{simples consideração} é propria da \emph{intelligencia possivel}, a qual, determinada pela \emph{especie intelligivel}, considera no phantasma a \emph{essencia} das coisas, sem se importar das condições materiaes. --- Como se vê, a primeira abstracção é \emph{causa} da segunda.\par c) \textsc{NATUREZA} DA INTELLIGENCIA AGENTE, --- A \emph{intelligencia agente} é uma faculdade \emph{activa}. Na verdade, uma faculdade é \emph{activa}, quando move ou produz o proprio objecto. Ora a \emph{intelligencia agente} move ou produz o proprio objecto; porque torna universal e \emph{actualmente} intelligivel a essencia, que só o etc \emph{potencialmente}; pois que á \emph{essencia} das coisas materiaes, considerada nas proprias coisas, é singular e concreta, e, como tal, não é \emph{actualmente} intelligivel, mas só o \emph{póde} ser mediante a abstracção.\par \emph{A intelligencia agente} é, portanto, uma faculdade realmente distincta da \emph{intelligencia possivel}. Na verdade, duas faculdades são realmente distinctas, quando uma é \emph{activa} e outra é \emph{passiva}. Ora, como provámos, á \emph{intelligencia agente} é faculdade \emph{activa}, porque produz ou move o seu objecto, tornando immaterial e \emph{actualmente} intelligivel o que era materiale só \emph{potencialmente} intelligivel; ao passo que 2 intelligencia \emph{possivel} é faculdade \emph{passiva}, porque é movida e determinada ao acto pelo seu objecto, que é a \emph{especie intelligivel}, actualmente existente nella. Logo a \emph{intelligencia agente} e á \emph{intelligencia possivel} são duas faculdades distinctas. Diz S. Thomaz: « \emph{Diversificatur} potentia intellectus \emph{agentis} et intellectus \emph{possibilis}: quia respectu ejusdem. objecti aliud principium oportet esse potentiam «\emph{activam}, quae \emph{facit} objectum, esse \emph{in actu}, et aliud potentiam \emph{passivam}, quao \emph{movetur} ab objecto \emph{im actu existente}. Et sic potentia \emph{activa} comparatur ad suum objectum ut \emph{ens in actu} ad \emph{ens in potentia}; potentia autem \emph{passiva} comparatur ad suum objectum e converso, ut \emph{ens in potentia} ad \emph{ens in actu} » (\emph{Sum}. \emph{Theol}., p. I, q. 719, a. \emph{V}). --- Mas, embora o S. Doutor distinga sempre á intelligencia \emph{agente} da \emph{possivel}, attribuindo a uma um acto diverso do da outra, comtudo declara expressamente que á intellecção não é \emph{duplice}, mas é \emph{una}, e que para esta unica intellecção concorrem --- \& intelligencia \emph{agente} emquanto \emph{prepara a materia} ou o \emph{objecto}, --- e á intelligencia \emph{possivel} emquanto \emph{formalmente} percebe o objecto, preparado pela intelligencia \emph{agente}. Diz o Angelico: « Duorum intellectuum, scilicet \emph{possibilis} et \emph{agentis}, sunt \emph{duae} actiones... Nec tamen sequitur quod sit \emph{duplex} intelligere in homine; quia ad \emph{unum} intelli-\par gere oportet quod \emph{utraque} istarum actionum concurrat » (\emph{De anima}, a. 4 ad 8).\par 4) \textsc{EXISTENCIA} DA INTELLIGENCIA AGENTE. --- Esta faculdade \emph{existe}, como\par fica demonstrado no texto. --- Repetimos que o nosso conhecimento intellectual\pgn{169} é sempre precedido pelo sensitivo, porque nada pode a intelligencia conkecer sem o phantasma, e este é producto da sensação. D'onde o adagio: \emph{nihil est in intellectu quod prius non fuerit in sensibus}. Com isto não se quer dizer que os sentidos conheçam \emph{tudo} o que a intelligencia conhece, mas que a intelligencia parte sempre das coisas, que os sentidos apprehenderam, para se elevar ao conhecimento das coisas, que escapam ao alcance dos sentidos e que são apprehendidas pela applicação das imagens das coisas sensiveis e á semelhança d'estas, emquanto se apprehendem envolvidas no phantasma. Diz S. Thomaz: « Dicitur cognitio mentis a sensu originem habere, non quod omne illud quod mens cognoscit, sensus apprehendat; sed quia ex his quae sensus apprehendit, mens in aliqua ulteriora manuducitur, sicut eliam sensibilia intellecta manuducunt in intelligibilia divinorum » (\emph{De Veril}, q. X, a. 6 ad 2). --- Noutro logar, o S. Doutor prova a existencia da intelligencia \emph{agente} por estas palavras: « Se os \emph{universaes} (que são o objecto da intelligencia possivel) subsistissem em si mesmos e fóra do nosso espirito, como ensinou Platão, seriam \emph{actualmente} intelligiveis, e não haveria necessidade da intelligencia \emph{agente}. Mas, como os \emph{universaes} não subsistem senão nas coisas concretas sensiveis, segundo a doutrina \textbf{de---Axrietoteles}, e por isso não são intelligiveis \emph{em acto}, mas só \emph{em potencia}, é necessario admittir uma força, que torne intelligivel \emph{em acto} o que só o era \emph{em potencia}, abstrahindo da materia e das condições da materia as especies das coisas. Esta força chama-se \emph{intelligencia agente} » (\emph{De spirit}. \emph{creat}., a. 4; Cf. \emph{De anima}, a. 4, --- \emph{Sum}. \emph{Theol}., p. I, q. 94, a. 4; q. 19, a. 8).\par Todavia alguns escriptores modernos não a querem admittir, pois que opinam, com Platão, que o objecto da nossa intelligencia é o immaterial e o universal puro e por si mesmo isento de todas as condições da materia.\par O fundamento, em que se basea esta opinião, é falso. O objecto \emph{proporcionado} da nossa intelligencia é a \emph{essencia} das coisas materiaes. E, como a \emph{essencia} das coisas materiaes, por estar envolvida nas propriedades ou condições materiaes, não póde unir-se directamente á nossa intelligencia (a qual, sendo espiritual ou immaterial, só tem por objecto o immaterial), é necessario admittir uma \emph{força} que, operando naquella imagem, \emph{abstraia} a essencia das condições ou propriedades materiaes, de modo que esta ser torne immaterial. --- Dêem a esta \emph{força abstractiva} o nome que quizerem, se não lhes agradar o de \emph{intelligencia agente}; mas ninguem poderá negar a sua existencia, sem contradizer o testemunho da razão e da consciencia. Não é questão de nome, é questão de realidade.\par Os \emph{Neoplatonicos}, \emph{Averrües}, \emph{Cousin} ensinaram que a intelligencia \emph{agente} é uma razão impessoal, uma faculdade subsistente em si e separada da alma, --- e os \emph{ontologistas} affirmam que aquella luz intellectual é a propria luz . intelligivel infinita, que é \textsc{DEus}.\par Estas opiniões são falsas.\par a) E' falsa \& opinião dos \emph{Neoplatonicos}, de -\emph{Averrües} e de \emph{Cousin}. Na verdade, a intelligencia \emph{agente}, sendo necessaria para o conhecimento intellectual, não póde deixar de ser uma faculdade subjectiva e propria da alma; porque seria absurdo que á mais nobre natureza visivel não possuisse, como as naturezas inferiores, tudo quanto é necessario para a sua operação.\par ,\pgn{170}--- Demais, á consciencia attesta que a operação, pela qual abstrahimos dos phantasmas a \emph{especie intelligivel}, é exclusivamente propria de cada um de nós. Ora esta operação não poderia ser propria de cada um de nós, se não derivasse de um principio proprio de cada um; pois um principio extrinseco, se póde levar um agente a operar, não póde fazer com que um agente, destituido de um principio intrinseco de actividade, produza uma operação sua propria. Logo a intelligencia \emph{agente} é uma faculdade intrinseca da alma. --- 4E) por isso que os escholasticos comparam a nossa intelligencia aos olhos e certos animaes, que não só têm a força de perceber os objectos, mas são dotados tambem de luz, que torna visiveis os mesmos objectos. (Cf. S. Thom. --- \emph{De spirit}. \emph{creat}., a. 10; \emph{De anima}, a. 5; \emph{Sum}. \emph{Theol}., p. I, q. 79, a. 5).\par \emph{b}) E' tambem falsa a doutrina dos \emph{Ontologistas}. \textsc{DEus} chama-se luz \emph{intelligivel}, --- não só emquanto creou as coisas, capazes de serem entendidas, --- mas tambem emquanto nos deu a luz, pela qual entendemos as mesmas coisas. Mas d'ahi não se segue que a luz, pela qual entendemos, seja a propria luz increada, --- mas só póde concluir-se que a luz, pela qual entendemos, foi creada por \textsc{Deus} e por Elle communicada á nossa alma (C£. S. Thom., \emph{De magistro}, a. 1 ad 7).\par Objectam: Assim como não é necessario um \emph{sentido agente}, que forme as \emph{especies sensiveis}, assim tambem não é necessaria a \emph{intelligencia agente}, que forme as \emph{especies intelligiveis}.\par Resposta. As \emph{especies sensiveis} existem \emph{actualmente}, produzidas pela acção dos objectos externos nos nossos sentidos, e por isso não é necessario um \emph{sentido} «\emph{gente}. Mas não se dá o mesmo com relação ás \emph{especies intelligiveis}. Estas não existem \emph{actualmente}; porque não são produzidas pela acção directa dos objectos na nossa intelligencia. E' necessario, pois, formal-as, abstrahindo do phantasma a \emph{essencia}, \emph{e} prescindindo das condições materiaes. Ora a força, que concorre para a formação das \emph{especies intelligiveis}, é a \emph{intelligencia agente}. Logo é necessario admittir a \emph{intelligencia agente}.\par O P. \emph{Manuel Bernardes} expõe a doutrina ácerca da \emph{i/ntelligencia agente} e da \emph{intelligencia possivel} nos termos seguintes:\par u O entendimento se divide em \emph{agente} e \emph{possivel}; os quaes se distinguem entre si, ou não mais que \emph{virtualmente}, como querem graves doutores; ou \emph{realmente}, como querem outros, e o dá a entender Santo Thomaz. Como quer que seja, realmente ha na alma os officios de ambos, pelos quaes se entenderão juntamente os seus nomes.\par \& Para o que é de saber, que depois que a fantasia formou a sua especie (os filosofos chamam-lhe \emph{fantasma expresso}) ou sello expressivo de qualquer objecto, cuja representação lhe havia entrado por algum dos sentidos particulares; havia de mandar ao entendimento outra especie d'este seu conhecimento, para o informar d'aquelle tal objecto, de que ella já estava informada. Mas por quanto o tal fantasma, ou sello é corporeo, como o é tambem essa fantasia, que o figurou, não póde informar o entendimento, por ser este potencia espiritual. :\par \& Que remedio, pois, para que a noticia do objecto material, que está cà\pgn{171}\emph{fóra}, passe adiante até ao entendimento espiritual? O remedio foi, que o entendimento \emph{agente} ajuntando-se com o tal fantasma corporeo, e elevando-o á ordem superior, por virtude propria que para isso tem, produzem ambos uma especie espiritual; a quem por parte do entendimento lhe veio o ser espiritual, e por parte do fantasma corporeo lhe veio o representar ainda o objecto, se bem por outro modo...\par \& Esta especie pois produzida por concurso do entendimento \emph{agente}, e do fantasma corporeo, se chama --- \emph{especie intelligivel}; e já como espiritual que é, serve nos moldes do entendimento \emph{possivel}, \emph{e} recebendo-se nelle, o faz entender o tal objecto.\par u E d'aqui se mostra a razão por que aquelle entendimento se chama \emph{agente}, e este outro \emph{possivel}. Porque o officio d'aquelle é fazer estas especies intelligiveis; e o d'este é recebel-as para entender por ellas; e como de sua natureza está a modo de um espelho liso, e nu de formas, com uma capacidade ou possibilidade universal para todas as especies intelligiveis, que lhe vierem, por isso se chama \emph{entendimento possivel}. \emph{E} assim como o espelho é todos os objectos representativamente: assim este entendimento é todas as cousas intencionalmente: \emph{Intellectus intelligendo fit omnia}.\par u Porém advirta-se, que supposto que ambas estas virtudes se chamam \emph{entendimentos}, o \emph{entendimento possivel} é o que entende, e não o \emph{agente}. Não o \emph{agente}; porque produzir especies intelligiveis, não é entender, senão fazer as cousas intelligiveis: assim como produzir o objecto externo (v. g. a parede alva, ou \& voz sonora) especies sensiveis, que são as que vemos rebater no espelho, ou ouvimos rebater nos eccos dos montes, não é sentir, senão fazer-se sensivel aos olhos, ou aos ouvidos. O \emph{possivel} sim; porque este entendimento é o que informado com a especie intelligivel produz o acto de entender; assim como a potencia visiva informada com a especie de parede alva produz o acto de ver: e a potencia de ouvir informada com a especie de voz sonora, produz o acto de ouvir.\par v De sorte que na parte intellectiva da alma racional ha duas producções: uma da especie intelligivel, outra do acto de entender: e por conseguinte ha duas recepções d'estas duas cousas: a producção da especie toca ao \emph{entendimento agente}; a do acto de entender ao \emph{entendimento possivel}: e s este mesmo tocam ambas as recepções, assim a da especie, que veio a parar nelle, como a do acto, que nasceu nelle.\par 4 Porém ainda que um só d'estes entendimentos é o que entende, ambos elles são lumes. Porque lume espiritual é tudo o que é manifestativo das cousas espirituaes: e o entendimento agente manifesta ao entendimento possivel as especies intelligiveis, que são espirituaes: o entendimento possivel manifesta á alma a verdade das cousas, e faz juizo d'ellas; e assim o seu objecto é a verdade em toda a sua amplitude, ou esfera » (\emph{Lus e Calor}, Primeira parte, Doutrina III).\par Voltaremos a este assumpto na secção seguinte. --- Ninguem exiranhe se, na presente secção que é destinada ao tratado das \emph{faculdades} do homem, dissemos alguma coisa ácerca das suas \emph{operações}. As faculdades, como provámos, distinguem-se e conhecem-se pelas suas operações. É, pois, natural que, no tratado das faculdades, se falle das operações.}.

\emph{SNC} A ada A 2 To ---Por---

%%>>> t121

\tese{121}{Varios nomes da intelligencia.}
A intelligencia, ficando sempre uma e a' mesma faculdade, toma varios nomes segundo \pg{165} \pg{166}  as diversas funegoes que exerce, e assim chama-se \emph{razão}, \emph{consciencia}, e \emph{memoria intellectual}.

%%>>> t122

\tese{122}{Razão.}
\emph{Razão} é a propria intelligencia, que, por meio do raciocinio, de verdades já conhecidas deduz outras \pg{167} \pg{168}  NE MC Án T ------\textbf{PENTT-RD EE} ---- \textbf{cost mE} Cc uL.

8) Divisão da razão: --- A \emph{razão} divide-se em \emph{theoretica} e \emph{pratica}. --- A \emph{theoretica}, ou \emph{especulativa}, limita-se á percepção \pg{170}  da verdade. A \emph{pratica} applica á operação a verdade percebida. --- A razão \emph{theoretica} e á \emph{pratica} não são duas faculdades distinctas \pg{171} , mas duas fuzegoes da mesma faculdade; pois que o objecto da razão \emph{theoretica} é o mesmo que o da \emph{pratica}, e é \pg{172}  apenas \emph{accidental} que a razão \emph{pratica} applique á operação a verdade, que a razão \emph{theoretica} considera em si mesma \footnote{\emph{Kant} (\emph{Critica da razão pura}) ensinou erradamente que á razão \emph{theoretica} e a. \emph{pratica} são duas faculdades distinctas e diversas.\par \emph{A} razão \emph{pratica}, applicando ás operações os principios universaes, forma juizos particulares ácerca de certos e determinados actos. Para isso, emprega um syllogismo. \emph{A} premissa \emph{maior} d'este syllogismo contém um principio universal; a premissa \emph{menor} indica que um certo e determinado sujeito está contido no sujeito do principio universal; a \emph{conclusão} infere que ao sujeito particular compete a propriedade, que se attribue ao sujeito universal. Exemplo: \emph{Os filhos devem respeitar seus paes}; \emph{Pedro é filho de Paulo}; \emph{logo}: \emph{Pedro deve respeitar Paulo} (\emph{Sum}. \emph{Theol}., p. I, q.. 79, a. 11).\par Com a razão \emph{pratica} está intimamente ligada a \emph{synderesis}. --- A \emph{synderesis} é o \emph{habito natural dos principios moraes} ou \emph{axiomas praticos}. Assim como a razão \emph{theoretica} discorre ácerca de assumptos \emph{especulativos}, assim a razão \emph{pratica} discorre ácerca de assumptos \emph{praticos}. Ora o discurso, como todo e qualquer movimento, deve ter o seu começo e o seu termo num ponto immoyel; e este ponto immovel deve ser constituido pelos primeiros principios. Deve, portanto, \emph{o} homem possuir os primeiros principios tanto da ordem \emph{especulativa} (todo o effeito exige uma causa), como da \emph{pratica} (deve fazer-se o bem; deve evitar-se o mal) Ora os primeiros principios da ordem \emph{especulativa} não pertencem a uma faculdade especial, mas derivam de uma disposição natural, que a intelligencia tem para formar juizos relativos aos principios \emph{theoreticos} ou \emph{especulativos}, \emph{e} que se chama \emph{habito dos principios}. Logo tambem os primeiros principios da ordem \emph{pratica} não derivam de uma faculdade especial, mas de uma disposição ou inclinação natural que a intelligencia tem para formar juizos relativos aos principios \emph{praticos}, e que se chama \emph{synderesis} (\emph{Sum}. \emph{Theol}., p. I, q. 79, a. 12). --- O habito dos primeiros principios \emph{praticos}, como tambem o dos primeiros principios \emph{especulativos}, diz-se \emph{natural}, emquanto se adquire sem esforço e por uma simplicissima abstracção da. intelligencia agente, \emph{e} não se corrompe nem \$e destroe. (\emph{Sun}. \emph{Theol}., 1-2, q. 58, a. 1).\par \emph{A razão} costuma tambem dividir-se em \emph{superior} e \emph{inferior}. A \emph{superior} julga das coisas, segundo as verdades eternas, necessarias e immutayeis da intelligencia divina, que é o exemplar de todos os entes creados. A \emph{inferior} julga das coisas, segundo as verdades temporaneas, accidentaes e mutayeis, que se conhecem pela consideração das coisas creadas. Assim a razão \emph{superior} afirma que se deve obedecer ao rei, porque \textsc{DEUS} o \emph{quer}; e á \emph{inferior} diz que se deve obedecer ao rei, porque a \emph{felicidade da nação o exige}.\par Como se vê, a divisão da \emph{razão} em \emph{superior} e \emph{inferior} basêa-se na qualidade do seu \emph{objecto}, que é a verdade; assim como a divisio da mesma faculdade em \emph{theoretica} e \emph{pratica} se funda no \emph{fim}, para que consideramos a verdade.\par \emph{A} razão \emph{inferior} e á \emph{superior} não são duas faculdades, mas duas funcçõesda mesma faculdade. Por quanto, as verdades eternas, necessarias e immutaveis, que são o objecto da razão \emph{superior}, e as verdades temporaneas, contingentes e mutaveis, que são o objecto da razão \emph{inferior}, embora sejam \emph{materialmente} muitas e diversas, todavia constituem um só objecto \emph{formal}; pois todas são attingidas emquanto \emph{émomateriaes} ou \emph{universaes}, e a unidade ou multiplicidade das faculdades depende, como dissemos, da unidade ou multiplicidade do seu objecto \emph{formal} (\emph{Sum}. \emph{Theol}. p. l, q. 79, a. 9).\par A razão \emph{superior}, ainda que seja mais nobre e mais elevada do que a \emph{inferior} (pois aquella julga das coisas por criterios mais altos do que esta), comtudo, na sua operação, é \emph{posterior} á razão \emph{inferior}; pois não podemos elevar-nos \& consideração do eterno e do immutayvel, senão pela consideração do temporaneo e do mutavel.\par \emph{Cousin} affirma que a razão \emph{superior}, sendo o meio pelo qual a razão divina communica com a razão humana, existe fôra de nós. --- Os \emph{ontologistas} dizem que a razão \emph{superior} é a propria razão \emph{absoluta} --- a divina.\par Estas opiniões são falsas. A razão \emph{superior} e a \emph{inferior} não são duas faculdades distinctas, mas indicam apenas duas funcções distinctas, embora não diversas \emph{formalmente}, de uma e a mesma faculdade. Sendo assim, é claro que a razão \emph{superior} não está acima de nós, mas está dentro de nós, é \emph{pessoal} e \emph{propria} de cada um de nós, como \emph{pessoal} e \emph{propria} de cada um de nós é a razão \emph{inferior}. :\par Objectam: O homem percebe o immutavel. Ora o immutayel não póde ser percebido pela razão humana, porque esta é mutavel. Logo deve ser percebido por uma \emph{rasdo superior} á nossa, e por isso a \emph{rasdo superior} esta \emph{fóra} ou \emph{acima} de nos.\par Respondemos que não só \emph{não repugna}, mas até é \emph{necessario} que a nossa alma perceba o \emph{immutavel} por meio de uma faculdade \emph{mautavel}. --- \emph{Não repugna}; pois que a nossa alma, na sua percepção, não \emph{cria} as propriedades ou relações das coisas, limita-se a investigal-as e descobril-as. --- \emph{E necessario}; visto que o meio, pelo qual conhecemos o immutavel, deve ser \emph{intrinseco} á nossa alma, pois o conhecimento é uma operação \emph{immanente}, e tudo o que está em a nossa alma é mutavel.}.

\pg{169}
